% Nombres complexes :Approche géométrique — Mathématiques Tle Tech — Cours signé OnBuch+
% Programme officiel MINESEC. Compiler avec : tectonic N11-nombres-complexes-approche-geometrique.tex
\newcommand{\DOCMATIERE}{Mathématiques}
\newcommand{\DOCNIVEAU}{Tle Tech}
\newcommand{\DOCMODULE}{Mathématiques – classe de Terminale technique}
\newcommand{\DOCLECON}{N11}
\newcommand{\DOCTITRE}{Nombres complexes :Approche géométrique}
% OnBuch+ — préambule « cours signés » ENRICHI (compilé avec tectonic / XeLaTeX)
% Système visuel « Pop » d'OnBuch+ : fond crème (jamais blanc), bordures encre
% épaisses, ombres « dures » décalées, titres Archivo Black en capitales.
% Version MATHÉMATIQUES : graphes (pgfplots/tikz), boîtes pédagogiques
% (définition, propriété, méthode, attention, le savais-tu, exercice résolu,
% exercice, TP, bilan), page de garde et signature OnBuch+ ancrée en bas.
%
% Macros attendues dans main.tex AVANT \input{preamble} :
%   \DOCMATIERE  \DOCNIVEAU  \DOCMODULE  \DOCLECON (numéro)  \DOCTITRE
\documentclass[11pt]{article}

\usepackage{fontspec}
\usepackage[french]{babel}
\usepackage[a4paper,margin=2cm,top=3.4cm,bottom=2.8cm,headheight=1.4cm,footskip=1.3cm]{geometry}
\usepackage{amsmath,amssymb,amsfonts}
\usepackage{mathtools} % \xmapsto, \coloneqq... souvent utilisés par les modèles
\usepackage{cancel} % simplifications barrées
% \abs{z} (module, valeur absolue) et \norm{u} (norme), utilisés par les modèles
\providecommand{\abs}{}\let\abs\relax\DeclarePairedDelimiter{\abs}{\lvert}{\rvert}
\providecommand{\norm}{}\let\norm\relax\DeclarePairedDelimiter{\norm}{\lVert}{\rVert}
\usepackage[table]{xcolor} % [table] : fournit aussi \rowcolors (alternance de lignes)
\usepackage{pagecolor}
\usepackage{tikz}
\usetikzlibrary{arrows.meta,positioning,calc,shapes.geometric,shapes.misc,shapes.arrows,decorations.pathmorphing,decorations.pathreplacing,decorations.markings,patterns,shadows,mindmap,trees,fit,backgrounds,matrix,intersections,angles,quotes}
\usepackage{pgfplots}
\usepackage{pgf-pie} % diagrammes circulaires \pie (statistiques)
\pgfplotsset{compat=1.18}
\usepgfplotslibrary{fillbetween,groupplots} % aires entre courbes (calcul intégral)
\usepackage{enumitem}
\usepackage{fancyhdr}
% [most] charge toutes les bibliothèques tcolorbox (listings, minted, raster,
% documentation...) et gonfle inutilement la mémoire TeX sur un document long
% (~300 boîtes) : on ne charge que ce qui est réellement utilisé.
\usepackage[skins,breakable]{tcolorbox}
\usepackage{graphicx}
\usepackage{colortbl}
\usepackage{booktabs}
\usepackage{tabularx}
\usepackage{diagbox} % case d'en-tête diagonale des tableaux à double entrée
% Colonnes extensibles centrée / gauche / droite (C, L, R), souvent utilisées par les modèles
\newcolumntype{C}{>{\centering\arraybackslash}X}
\newcolumntype{L}{>{\raggedright\arraybackslash}X}
\newcolumntype{R}{>{\raggedleft\arraybackslash}X}
\usepackage{array}
\usepackage{multicol}
\usepackage{siunitx}
% parse-numbers=false : accepte aussi \SI{2,5 \times 10^{-3}}{...}
\sisetup{locale=FR,output-decimal-marker={,},group-minimum-digits=4,parse-numbers=false}
\DeclareSIUnit\ohm{\ensuremath{\symup{\Omega}}} % U+2126 absent de la police
\usepackage{qrcode}
% \degree : commande fréquemment utilisée par les modèles pour un angle ou une
% température en dehors de \SI{}{} (non fournie par les paquets chargés).
\providecommand{\degree}{^{\circ}}
% \qed (fin de démonstration), utilisé par les modèles sans amsthm
\providecommand{\qed}{\hfill$\square$}
% \barre : double barre des valeurs interdites dans un tableau de variations (tkz-tab)
\providecommand{\barre}{\ensuremath{\|}}
% environnement proof (amsthm non chargé) utilisé par les modèles
\ifdefined\proof\else\newenvironment{proof}[1][Démonstration]{\par\noindent\textit{#1.}\ }{\qed\par}\fi
\usepackage{lastpage}
\usepackage{hyperref}
\hypersetup{hidelinks}

% ============================================================
% Palette « Pop » OnBuch+ — mêmes hex que la classe P (plus_ui.dart)
% ============================================================
\definecolor{poporange}{HTML}{F59321}
\definecolor{popdark}{HTML}{E07A0C}
\definecolor{popcream}{HTML}{FFF6E8}
\definecolor{popink}{HTML}{1C1714}
\definecolor{poppurple}{HTML}{7A5AE0}
\definecolor{popgreen}{HTML}{2E9E6A}
\definecolor{poppink}{HTML}{E0457B}
\definecolor{popblue}{HTML}{2F7DE1}
\definecolor{popgold}{HTML}{C98A12}
\definecolor{popmuted}{HTML}{8A7F76}
\definecolor{popline}{HTML}{EADFD2}
\definecolor{popgray}{HTML}{8A7F76}
\colorlet{popgrey}{popgray}
% Alias tolérants (le modèle nomme parfois les couleurs différemment)
\colorlet{popwhite}{white}
\colorlet{popblack}{popink}
\colorlet{popred}{poppink}
\colorlet{popyellow}{popgold}
% Teintes claires pour fonds de tableaux / zones
\colorlet{popblueL}{popblue!10!white}
\colorlet{poporangeL}{poporange!14!white}
\colorlet{popgreenL}{popgreen!12!white}
\colorlet{poppurpleL}{poppurple!12!white}
\colorlet{poppinkL}{poppink!12!white}
\colorlet{popgoldL}{popgold!14!white}

\pagecolor{popcream}

% ============================================================
% Typographie
% ============================================================
\defaultfontfeatures{Ligatures=TeX}
\setmainfont{PlusJakartaSans-Medium.ttf}[Path=../../../fonts/,BoldFont=PlusJakartaSans-Bold.ttf]
% Maths en sans-serif (Fira Math) avec chiffres et lettres droites de Plus
% Jakarta, pour que les formules se fondent dans le texte.
\usepackage[math-style=ISO,bold-style=ISO]{unicode-math}
\setmathfont{FiraMath-Regular.otf}[Path=../../../fonts/]
\setmathfont{PlusJakartaSans-Medium.ttf}[Path=../../../fonts/,range={up/num,up/{Latin,latin},\mathup}]
\usepackage{icomma} % virgule décimale sans espace en mode math
% Caractères mathématiques Unicode absents des polices de texte (Plus Jakarta,
% Archivo Black) : rendus par leur équivalent mathématique, en texte comme en
% formule — sinon ils s'affichent en carré vide (ex. « Arithmétique dans ℤ »).
\usepackage{newunicodechar}
\newunicodechar{ℝ}{\ensuremath{\mathbb{R}}}
\newunicodechar{ℤ}{\ensuremath{\mathbb{Z}}}
\newunicodechar{ℂ}{\ensuremath{\mathbb{C}}}
\newunicodechar{ℕ}{\ensuremath{\mathbb{N}}}
\newunicodechar{ℚ}{\ensuremath{\mathbb{Q}}}
\newunicodechar{𝔻}{\ensuremath{\mathbb{D}}}
\newunicodechar{α}{\ensuremath{\alpha}}
\newunicodechar{β}{\ensuremath{\beta}}
\newunicodechar{γ}{\ensuremath{\gamma}}
\newunicodechar{δ}{\ensuremath{\delta}}
\newunicodechar{ε}{\ensuremath{\varepsilon}}
\newunicodechar{θ}{\ensuremath{\theta}}
\newunicodechar{λ}{\ensuremath{\lambda}}
\newunicodechar{μ}{\ensuremath{\mu}}
\newunicodechar{σ}{\ensuremath{\sigma}}
\newunicodechar{τ}{\ensuremath{\tau}}
\newunicodechar{φ}{\ensuremath{\varphi}}
\newunicodechar{ψ}{\ensuremath{\psi}}
\newunicodechar{ω}{\ensuremath{\omega}}
\newunicodechar{Σ}{\ensuremath{\Sigma}}
% majuscules produites par \MakeUppercase dans les titres (α-aminés -> Α)
\newunicodechar{Α}{\ensuremath{\alpha}}
\newunicodechar{Β}{\ensuremath{\beta}}
\newunicodechar{Γ}{\ensuremath{\Gamma}}
\newunicodechar{Δ}{\ensuremath{\Delta}}
\newunicodechar{Ε}{\ensuremath{\varepsilon}}
\newunicodechar{Θ}{\ensuremath{\Theta}}
\newunicodechar{Λ}{\ensuremath{\Lambda}}
\newunicodechar{Μ}{\ensuremath{\mu}}
\newunicodechar{Τ}{\ensuremath{\tau}}
\newunicodechar{Φ}{\ensuremath{\Phi}}
\newunicodechar{Ψ}{\ensuremath{\Psi}}
\newunicodechar{Ω}{\ensuremath{\Omega}}
\newunicodechar{↦}{\ensuremath{\mapsto}}
\newunicodechar{∈}{\ensuremath{\in}}
\newunicodechar{∉}{\ensuremath{\notin}}
\newunicodechar{∧}{\ensuremath{\wedge}}
\newunicodechar{⇔}{\ensuremath{\Leftrightarrow}}
\newunicodechar{⟺}{\ensuremath{\Longleftrightarrow}}
\newunicodechar{⇒}{\ensuremath{\Rightarrow}}
\newunicodechar{⟹}{\ensuremath{\Longrightarrow}}
\newunicodechar{∘}{\ensuremath{\circ}}
\newunicodechar{∪}{\ensuremath{\cup}}
\newunicodechar{∩}{\ensuremath{\cap}}
\newunicodechar{≡}{\ensuremath{\equiv}}
\newunicodechar{⊂}{\ensuremath{\subset}}
\newunicodechar{⊥}{\ensuremath{\perp}}
\newunicodechar{∃}{\ensuremath{\exists}}
\newunicodechar{∀}{\ensuremath{\forall}}
\newunicodechar{∛}{\ensuremath{\sqrt[3]{\,}}}
\newunicodechar{∜}{\ensuremath{\sqrt[4]{\,}}}
\newunicodechar{⃗}{\ensuremath{{}^{\to}}}
\newunicodechar{ⁿ}{\ifmmode^{n}\else\textsuperscript{\ensuremath{n}}\fi}
\newunicodechar{ˣ}{\ifmmode^{x}\else\textsuperscript{\ensuremath{x}}\fi}
\newunicodechar{ᵇ}{\ifmmode^{b}\else\textsuperscript{\ensuremath{b}}\fi}
\newunicodechar{ʳ}{\ifmmode^{r}\else\textsuperscript{\ensuremath{r}}\fi}
\newunicodechar{ᵘ}{\ifmmode^{u}\else\textsuperscript{\ensuremath{u}}\fi}
\newunicodechar{ʸ}{\ifmmode^{y}\else\textsuperscript{\ensuremath{y}}\fi}
\newunicodechar{ᵃ}{\ifmmode^{a}\else\textsuperscript{\ensuremath{a}}\fi}
\newunicodechar{ᵉ}{\ifmmode^{e}\else\textsuperscript{\ensuremath{e}}\fi}
\newunicodechar{ᵏ}{\ifmmode^{k}\else\textsuperscript{\ensuremath{k}}\fi}
\newunicodechar{ᵖ}{\ifmmode^{p}\else\textsuperscript{\ensuremath{p}}\fi}
\newunicodechar{ᴺ}{\ifmmode^{N}\else\textsuperscript{\ensuremath{N}}\fi}
\newunicodechar{ᶜ}{\ifmmode^{c}\else\textsuperscript{\ensuremath{c}}\fi}
\newunicodechar{ʰ}{\ifmmode^{h}\else\textsuperscript{\ensuremath{h}}\fi}
\newunicodechar{ᵠ}{\ifmmode^{q}\else\textsuperscript{\ensuremath{q}}\fi}
\newunicodechar{ₙ}{\ifmmode_{n}\else\textsubscript{\ensuremath{n}}\fi}
\newunicodechar{ₐ}{\ifmmode_{a}\else\textsubscript{\ensuremath{a}}\fi}
\newunicodechar{ᵢ}{\ifmmode_{i}\else\textsubscript{\ensuremath{i}}\fi}
\newunicodechar{ᵣ}{\ifmmode_{r}\else\textsubscript{\ensuremath{r}}\fi}
\newunicodechar{ₖ}{\ifmmode_{k}\else\textsubscript{\ensuremath{k}}\fi}
\newunicodechar{ᵦ}{\ifmmode_{\beta}\else\textsubscript{\ensuremath{\beta}}\fi}
\newunicodechar{ₓ}{\ifmmode_{x}\else\textsubscript{\ensuremath{x}}\fi}
\newfontfamily\popdisplay{ArchivoBlack-Regular.ttf}[Path=../../../fonts/]
\newfontfamily\poplogo{SpaceGrotesk-Bold.ttf}[Path=../../../fonts/]
\newfontfamily\popsemi{PlusJakartaSans-SemiBold.ttf}[Path=../../../fonts/]
\newfontfamily\popxbold{PlusJakartaSans-ExtraBold.ttf}[Path=../../../fonts/]
\newfontfamily\popxboldls{PlusJakartaSans-ExtraBold.ttf}[Path=../../../fonts/,LetterSpace=3.0]

\setlist[enumerate]{leftmargin=1.7em,itemsep=5pt,topsep=4pt}
\setlist[itemize]{leftmargin=1.7em,itemsep=4pt,topsep=4pt,label={\color{poporange}\textbullet}}
\setlength{\parindent}{0pt}
\setlength{\parskip}{5pt}
\renewcommand{\arraystretch}{1.25}

% pgfplots : style Pop par défaut
\pgfplotsset{
  every axis/.append style={axis line style={popink,line width=1pt},tick style={popink},
    grid style={popline,line width=0.5pt},label style={font=\small},tick label style={font=\footnotesize}},
  popcurve/.style={poporange,line width=1.8pt,no markers,smooth},
}
% Même style disponible dans un \draw TikZ simple (hors pgfplots)
\tikzset{popcurve/.style={poporange,line width=1.8pt}}

% ============================================================
% Pill / squiggle
% ============================================================
\newcommand{\poppill}[3]{%
  \tikz[baseline=-0.55ex]\node[draw=popink,line width=1.2pt,rounded corners=2.6pt,%
    fill=#2,inner xsep=6.5pt,inner ysep=3pt]{%
    \popxboldls\fontsize{7}{7}\selectfont\color{#3}\MakeUppercase{#1}};%
}
\newcommand{\popsquiggle}[1]{%
  \tikz[baseline=-0.4ex]{
    \draw[#1,line width=2.6pt,line cap=round]
      plot [smooth] coordinates {(0,0.09) (0.34,0.01) (0.68,0.21) (1.02,0.01) (1.36,0.10)};
  }%
}
% Mot-clé surligné (marqueur orange sous le texte)
\newcommand{\cle}[1]{{\popxbold\color{popdark}#1}}

% ============================================================
% En-tête / pied
% ============================================================
\pagestyle{fancy}
\fancyhf{}
\renewcommand{\headrulewidth}{0pt}
\renewcommand{\footrulewidth}{0pt}
\lhead{\raisebox{-0.15ex}{\poplogo\fontsize{13}{13}\selectfont\color{popink}OnBuch\color{poporange}+}}
\rhead{\poppill{\DOCMATIERE\ \textbullet\ \DOCNIVEAU\ \textbullet\ Leçon \DOCLECON}{white}{popink}}
\lfoot{\popxbold\fontsize{7.6}{7.6}\selectfont\color{popmuted}\MakeUppercase{Cours signé OnBuch+}\ \textbullet\ {\popsemi\DOCTITRE}}
\rfoot{\tikz[baseline=-0.6ex]\node[rounded corners=8.5pt,fill=popink,minimum height=17pt,minimum width=17pt,inner xsep=5pt,inner sep=0]{\popxbold\fontsize{8}{8}\selectfont\color{popcream}\thepage\,/\,\pageref*{LastPage}};}

% ============================================================
% Titres
% ============================================================
\newcommand{\chaptitre}[2]{%
  \begin{tcolorbox}[enhanced,colback=poporange,colframe=popink,boxrule=2.6pt,arc=11pt,
      drop shadow=popink,left=15pt,right=15pt,top=12pt,bottom=11pt,before skip=3pt,after skip=18pt]
    \poppill{#2}{popink}{popcream}\par\vspace{9pt}
    {\raggedright\hyphenpenalty=10000\exhyphenpenalty=10000\popdisplay\fontsize{22}{24}\selectfont\color{popcream}\MakeUppercase{#1}\par}
  \end{tcolorbox}
}
% Section numérotée : pastille orange avec le numéro + titre Archivo
\newcounter{popsec}
\newcommand{\coursec}[1]{%
  \stepcounter{popsec}\setcounter{popsub}{0}%
  \par\vspace{16pt}\noindent
  \begin{minipage}{\linewidth}
  \tikz[baseline=-0.7ex]\node[draw=popink,line width=1.6pt,fill=poporange,rounded corners=4pt,minimum size=22pt,inner sep=0]{\popdisplay\fontsize{12}{12}\selectfont\color{popcream}\thepopsec};%
  \hspace{8pt}{\popdisplay\fontsize{15}{17}\selectfont\color{popink}\MakeUppercase{#1}}\par
  \vspace{4pt}\noindent{\color{poporange}\rule{36pt}{3.6pt}}
  \end{minipage}\par\nopagebreak\vspace{8pt}
}
\newcounter{popsub}
\newcommand{\courssub}[1]{%
  \stepcounter{popsub}%
  \par\vspace{9pt}\noindent{\popxbold\fontsize{11.5}{13}\selectfont\color{popink}{\color{poporange}\thepopsec.\thepopsub}\hspace{6pt}#1}\par\nopagebreak\vspace{4pt}
}

% ============================================================
% Boîtes pédagogiques — même recette : fond blanc, bordure épaisse colorée,
% ombre dure encre, badge plein. Titre optionnel [#1].
% ============================================================
\tcbset{popbase/.style={breakable,enhanced,colback=white,boxrule=2.3pt,arc=9pt,
  shadow={3pt}{-3pt}{0pt}{popink},left=14pt,right=14pt,top=11pt,bottom=11pt,before skip=11pt,after skip=15pt,
  fontupper=\popsemi\fontsize{10.6}{14.5}\selectfont\color{popink},
  fontlower=\popsemi\fontsize{10.6}{14.5}\selectfont\color{popink}}}
% Coche dessinée (le glyphe ✓ n'existe pas dans les polices du document)
\newcommand{\popcheck}[1]{\tikz[baseline=-0.5ex]\draw[#1,line width=1.6pt,line cap=round,line join=round](-0.13,0.0)--(-0.04,-0.09)--(0.13,0.1);}
\newcommand{\popboxhead}[3]{\poppill{#1}{#2}{white}\ifx\relax#3\relax\else\hspace{6pt}{\popxbold\fontsize{10.6}{12}\selectfont\color{popink}#3}\fi\par\vspace{7pt}}

\newtcolorbox{aretenir}[1][]{popbase,colframe=popgreen,before upper={\popboxhead{À retenir}{popgreen}{#1}}}
\newtcolorbox{exemplebox}[1][]{popbase,colframe=poporange,before upper={\popboxhead{Exemple}{poporange}{#1}}}
\newtcolorbox{definition}[1][]{popbase,colframe=popblue,before upper={\popboxhead{Définition}{popblue}{#1}}}
\newtcolorbox{propriete}[1][]{popbase,colframe=poppurple,before upper={\popboxhead{Propriété}{poppurple}{#1}}}
\newtcolorbox{methode}[1][]{popbase,colframe=popgold,before upper={\popboxhead{Méthode}{popgold}{#1}}}
\newtcolorbox{attention}[1][]{popbase,colframe=poppink,before upper={\popboxhead{Attention}{poppink}{#1}}}
\newtcolorbox{experience}[1][]{popbase,colframe=popblue,colback=popblueL,before upper={\popboxhead{Expérience}{popblue}{#1}}}
% « Le savais-tu ? » : carte violette pleine (contexte, culture, Cameroun)
\newtcolorbox{savaistu}[1][]{popbase,colback=poppurple,colframe=popink,
  fontupper=\popsemi\fontsize{10.4}{14.2}\selectfont\color{white},
  before upper={\poppill{Le savais-tu\,?}{white}{poppurple}\ifx\relax#1\relax\else\hspace{6pt}{\popxbold\color{white}#1}\fi\par\vspace{7pt}}}
% Exercice résolu : énoncé en haut, \tcblower puis la solution sur fond crème
\newcounter{popexo}\newcounter{popexr}
\newtcolorbox{exoresolu}[1][]{popbase,colframe=poporange,colbacklower=popcream,
  segmentation style={popink,line width=1.2pt,dashed},
  before upper={\stepcounter{popexr}\popboxhead{Exercice résolu \thepopexr}{poporange}{#1}},
  before lower={\poppill{Solution}{popink}{popcream}\par\vspace{6pt}}}
% Exercice (non résolu en place) — la correction est dans la partie Corrigés
\newtcolorbox{exercice}[1][]{popbase,colframe=popink,
  before upper={\stepcounter{popexo}\popboxhead{Exercice \thepopexo}{popink}{#1}}}
% Corrigé
\newtcolorbox{corrige}[1][]{popbase,colframe=popgreen,colback=popgreenL,
  before upper={\popboxhead{Corrigé}{popgreen}{#1}}}
% Objectifs / prérequis / bilan
\newtcolorbox{objectifs}{popbase,colframe=popink,colback=poporangeL,
  before upper={\popboxhead{Objectifs de la leçon}{popink}{}}}
\newtcolorbox{prerequis}{popbase,colframe=popink,colback=popblueL,
  before upper={\popboxhead{Prérequis}{popblue}{}}}
\newtcolorbox{bilan}{popbase,colframe=popink,colback=white,boxrule=2.8pt,
  before upper={\popboxhead{Fiche bilan}{poporange}{L'essentiel à connaître pour l'examen}}}
% Figure encadrée avec légende
\newtcolorbox{popfigure}[1][]{enhanced,breakable=false,colback=white,colframe=popline,boxrule=1.4pt,arc=8pt,
  left=10pt,right=10pt,top=10pt,bottom=8pt,before skip=10pt,after skip=12pt,halign=center,
  lower separated=false,halign lower=center,
  fontlower=\popsemi\fontsize{9}{11.5}\selectfont\color{popmuted},
  #1}
\newcommand{\legende}[1]{\par\vspace{6pt}{\popsemi\fontsize{9}{11.5}\selectfont\color{popmuted}#1}}

% ============================================================
% Signature OnBuch+
% ============================================================
\newcommand{\poplogoinline}[1]{{\poplogo\fontsize{#1}{#1}\selectfont\color{popink}OnBuch\color{poporange}+}}
\newcommand{\signatureonbuch}{%
  \begin{tcolorbox}[enhanced,colback=popink,colframe=popink,boxrule=2.2pt,arc=10pt,
      drop shadow=poporange,left=14pt,right=14pt,top=10pt,bottom=10pt,before skip=0pt,after skip=0pt]
    \tikz[baseline=-0.7ex]\node[circle,fill=popgreen,draw=popcream,line width=1.3pt,minimum size=22pt,inner sep=0]{\popcheck{white}};%
    \hspace{9pt}\begin{minipage}[c]{0.62\linewidth}
      {\popxbold\fontsize{12}{13}\selectfont\color{popcream}Signé\ \textbullet\ {\poplogo OnBuch\color{poporange}+}}\par\vspace{2pt}
      {\raggedright\popsemi\fontsize{8}{10.5}\selectfont\color{popline}\MakeUppercase{Équipe pédagogique OnBuch+}\\\MakeUppercase{Conforme au programme officiel MINESEC}\par}
    \end{minipage}\hfill
    {\poplogo\fontsize{22}{22}\selectfont\color{popcream}OnBuch\color{poporange}+}
  \end{tcolorbox}%
}

% ============================================================
% Page de garde
% #1 titre · #2 matière · #3 niveau · #4 module · #5 n° leçon · #6 durée ·
% #7 description / objectifs (texte ou itemize)
% ============================================================
\newcommand{\pagegarde}[7]{%
  \thispagestyle{empty}
  \newgeometry{a4paper,margin=1.7cm,top=1.6cm,bottom=1.4cm}%
  \begin{tikzpicture}[remember picture,overlay]
    \fill[poporange!18] ([xshift=-3.2cm,yshift=-3.6cm]current page.north east) circle (4.6cm);
    \fill[poppurple!14] ([xshift=2.4cm,yshift=7.6cm]current page.south west) circle (3.8cm);
  \end{tikzpicture}%
  {\poplogo\fontsize{30}{30}\selectfont\color{popink}OnBuch\color{poporange}+}\hfill
  \raisebox{0.6ex}{\poppill{Cours signé \textbullet\ Premium}{popink}{popcream}}\par
  \vspace{1pt}\popsquiggle{poppurple}\par
  \vspace{8pt}
  \begin{tcolorbox}[enhanced,colback=poporange,colframe=popink,boxrule=2.8pt,arc=13pt,
      drop shadow=popink,left=18pt,right=18pt,top=12pt,bottom=13pt,after skip=10pt]
    \poppill{#2 \textbullet\ #3}{popink}{popcream}\hspace{5pt}\poppill{#4}{white}{popink}\par\vspace{8pt}
    {\popdisplay\fontsize{14}{14}\selectfont\color{popink}LEÇON #5}\par\vspace{4pt}
    {\raggedright\hyphenpenalty=10000\exhyphenpenalty=10000\popdisplay\fontsize{25}{27}\selectfont\color{popcream}\MakeUppercase{#1}\par}
    \vspace{8pt}
    {\popxbold\fontsize{9.5}{9.5}\selectfont\color{popink}\MakeUppercase{Durée conseillée : #6}}
  \end{tcolorbox}
  \begin{tcolorbox}[enhanced,colback=white,colframe=popink,boxrule=2.2pt,arc=10pt,
      drop shadow=popink,left=16pt,right=16pt,top=10pt,bottom=10pt,after skip=8pt]
    \poppill{Dans cette leçon}{poppurple}{white}\par\vspace{5pt}
    \popsemi\fontsize{9.3}{12.6}\selectfont\color{popink}#7
  \end{tcolorbox}
  \noindent\begin{minipage}[t]{0.487\linewidth}
    \begin{tcolorbox}[enhanced,colback=popgreen,colframe=popink,boxrule=2.2pt,arc=10pt,
        drop shadow=popink,left=11pt,right=11pt,top=10pt,bottom=10pt,height=3.1cm,valign=center]
      \tcbox[colback=white,colframe=popink,boxrule=1.3pt,arc=5pt,boxsep=0pt,nobeforeafter,
        left=3pt,right=3pt,top=3pt,bottom=3pt,valign=center]{\qrcode[height=1.9cm]{https://play.google.com/store/apps/details?id=com.luvvix.onbuch&hl=fr}}%
      \hspace{8pt}\begin{minipage}[c]{0.5\linewidth}
        {\popdisplay\fontsize{11}{12.5}\selectfont\color{white}\MakeUppercase{Télécharge OnBuch}}\par\vspace{4pt}
        {\raggedright\popsemi\fontsize{8.4}{11}\selectfont\color{white}Gratuit \textbullet\ Google Play\\Tuteur IA, annales, concours, résultats\par}
      \end{minipage}
    \end{tcolorbox}
  \end{minipage}\hfill
  \begin{minipage}[t]{0.487\linewidth}
    \begin{tcolorbox}[enhanced,colback=poppurple,colframe=popink,boxrule=2.2pt,arc=10pt,
        drop shadow=popink,left=11pt,right=11pt,top=10pt,bottom=10pt,height=3.1cm,valign=center]
      \tikz[baseline=-0.7ex]\node[circle,fill=white,draw=popink,line width=1.3pt,minimum size=1.6cm,inner sep=0]{\popdisplay\fontsize{24}{24}\selectfont\color{poppurple}+};%
      \hspace{8pt}\begin{minipage}[c]{0.6\linewidth}
        {\popdisplay\fontsize{11}{12.5}\selectfont\color{white}\MakeUppercase{Passe à OnBuch+}}\par\vspace{4pt}
        {\raggedright\popsemi\fontsize{8.4}{11}\selectfont\color{white}Tous les cours signés, Tuteur IA illimité, fiches PDF — onglet OnBuch+ de l'app\par}
      \end{minipage}
    \end{tcolorbox}
  \end{minipage}
  \vspace{8pt}
  \popcontenu
  \vfill
  \signatureonbuch
  \restoregeometry
  \newpage
}

% Rangée « Ce cours contient » (page de garde)
\newcommand{\poptile}[3]{% #1 couleur #2 gros texte #3 légende
  \begin{tcolorbox}[enhanced,colback=#1,colframe=popink,boxrule=2pt,arc=9pt,shadow={2.5pt}{-2.5pt}{0pt}{popink},
      left=6pt,right=6pt,top=9pt,bottom=9pt,halign=center,valign=center,height=1.75cm,width=\linewidth,nobeforeafter]
    {\popdisplay\fontsize{12.5}{13}\selectfont\color{white}\MakeUppercase{#2}}\par\vspace{3pt}
    {\centering\popsemi\fontsize{7.8}{9.5}\selectfont\color{white}#3\par}
  \end{tcolorbox}}
\newcommand{\popcontenu}{%
  \par\noindent{\popxboldls\fontsize{7.5}{7.5}\selectfont\color{popmuted}CE COURS SIGNÉ CONTIENT}\par\vspace{7pt}
  \noindent\begin{minipage}[t]{0.235\linewidth}\poptile{popblue}{Cours}{complet, illustré, pas à pas}\end{minipage}\hfill
  \begin{minipage}[t]{0.235\linewidth}\poptile{poporange}{Méthodes}{et exercices résolus}\end{minipage}\hfill
  \begin{minipage}[t]{0.235\linewidth}\poptile{poppink}{Exercices}{type examen + corrigés}\end{minipage}\hfill
  \begin{minipage}[t]{0.235\linewidth}\poptile{popgreen}{Fiche bilan}{l'essentiel à retenir}\end{minipage}\par}

% Sommaire
\newtcolorbox{sommaire}{enhanced,colback=white,colframe=popink,boxrule=2.2pt,arc=10pt,
  drop shadow=popink,left=18pt,right=18pt,top=13pt,bottom=13pt,before skip=6pt,after skip=20pt,
  before upper={\poppill{Sommaire}{poppurple}{white}\par\vspace{9pt}\popsemi\fontsize{10.6}{14.5}\selectfont\color{popink}}}

% Signature de fin de cours — poussée tout en bas de la dernière page
\newcommand{\findecours}{%
  \par\vfill\nopagebreak
  \begin{center}\popsquiggle{poporange}\end{center}\vspace{4pt}
  \signatureonbuch
}
\AtBeginDocument{\def\llbracket{\mathopen{[\mkern-3.5mu[}}\def\rrbracket{\mathclose{]\mkern-3.5mu]}}} % intervalles d'entiers (glyphe ⟦ absent de la police)
% Glyphes absents de Fira Math : redéfinis à partir de glyphes disponibles
\AtBeginDocument{%
  \def\setminus{\mathbin{\backslash}}\let\smallsetminus\setminus
  \def\vdots{\mathord{\vbox{\baselineskip3.2pt\lineskiplimit0pt\kern4pt\hbox{.}\hbox{.}\hbox{.}}}}%
  \def\ddots{\mathinner{\mkern1mu\raise6.5pt\hbox{.}\mkern2mu\raise3.6pt\hbox{.}\mkern2mu\raise0.7pt\hbox{.}\mkern1mu}}%
  \def\triangle{\mathord{\scalebox{0.9}{$\symup{\Delta}$}}}%
  \def\nabla{\mathord{\raisebox{\depth}{\scalebox{1}[-1]{$\symup{\Delta}$}}}}%
  \def\checkmark{\popcheck{popdark}}%
  \def\star{\mathbin{\ast}}\def\bigstar{\mathord{\ast}}%
  \def\diamond{\mathbin{\scalebox{0.6}{$\Diamond$}}}%
  \def\bigcup{\mathop{\mathchoice{\raisebox{-0.3em}{\scalebox{1.7}{$\cup$}}}{\scalebox{1.25}{$\cup$}}{\cup}{\cup}}\displaylimits}%
  \def\bigcap{\mathop{\mathchoice{\raisebox{-0.3em}{\scalebox{1.7}{$\cap$}}}{\scalebox{1.25}{$\cap$}}{\cap}{\cap}}\displaylimits}%
  \def\lesseqqgtr{\mathrel{\lessgtr}}\def\bigoplus{\mathop{\oplus}\displaylimits}%
}
\providecommand{\ssi}{si et seulement si} % abréviation fréquente
\AtBeginDocument{\providecommand{\wideparen}{\overparen}} % arc de cercle

%% FIGFIX-BEGIN v5 — correctifs globaux des figures (docs/onbuch-generation/tools/figfix)
\usepackage{adjustbox}\usepackage{etoolbox}\tcbuselibrary{raster}
% toute figure TikZ/pgfplots plus large que la ligne est réduite à la largeur disponible
\BeforeBeginEnvironment{tikzpicture}{\begin{adjustbox}{max width=\linewidth}}
\AfterEndEnvironment{tikzpicture}{\end{adjustbox}}
% légendes pgfplots lisibles par-dessus les courbes
\pgfplotsset{every axis/.append style={legend style={fill=white,fill opacity=0.92,text opacity=1,draw=black!15,inner sep=3pt}}}
% étiquettes d'arêtes (syntaxe edge["..."]) avec fond blanc
\tikzset{every edge quotes/.append style={fill=white,inner sep=1.2pt}}
%% FIGFIX-END
\begin{document}

\pagegarde{Nombres complexes :Approche géométrique}{Mathématiques}{Tle Tech}{Mathématiques – classe de Terminale technique}{N11}{3 séances (fiche de progression harmonisée)}{Cette leçon présente la géométrie des nombres complexes, la forme exponentielle d'Euler, la linéarisation des puissances trigonométriques et la recherche des racines n‑ièmes, en insistant sur la rédaction rigoureuse et l'application à des problèmes techniques concrets.
\vspace{4pt}
\begin{multicols}{2}\small
\begin{itemize}
\item À la fin de cette leçon, je sais représenter un nombre complexe dans le plan et en déterminer le module et l'argument principal.
\item Je sais écrire un nombre complexe sous forme exponentielle et l'utiliser pour simplifier produits, quotients et puissances.
\item Je sais appliquer les formules de Moivre pour linéariser cosⁿθ et sinⁿθ et pour exprimer cos(nθ), sin(nθ) en puissances.
\item Je sais déterminer les n racines n‑ièmes d'un nombre complexe non nul et les représenter géométriquement.
\item Je sais modéliser un signal alternatif (tension, courant) à l'aide des nombres complexes et en déduire des grandeurs physiques (impédance, puissance).
\item Je sais rédiger une démonstration complète (ex. formule de Moivre) en justifiant chaque étape.
\end{itemize}\end{multicols}}

\begin{sommaire}
\begin{multicols}{2}
\begin{enumerate}
\item 1. Représentation géométrique d'un nombre complexe
\item 2. Forme exponentielle (Euler) et calculs simplifiés
\item 3. Linéarisation (formules de Moivre) et puissances trigonométriques
\item 4. Racines n‑ièmes d'un nombre complexe non nul
\item 5. Applications techniques et rédaction rigoureuse
\item Activité d'intégration
\item Méthodes et exercices résolus
\item Exercices
\item Corrigés des exercices
\item Fiche bilan
\end{enumerate}
\end{multicols}
\end{sommaire}

\begin{objectifs}
\begin{itemize}
\item À la fin de cette leçon, je sais représenter un nombre complexe dans le plan et en déterminer le module et l'argument principal.
\item Je sais écrire un nombre complexe sous forme exponentielle et l'utiliser pour simplifier produits, quotients et puissances.
\item Je sais appliquer les formules de Moivre pour linéariser cosⁿθ et sinⁿθ et pour exprimer cos(nθ), sin(nθ) en puissances.
\item Je sais déterminer les n racines n‑ièmes d'un nombre complexe non nul et les représenter géométriquement.
\item Je sais modéliser un signal alternatif (tension, courant) à l'aide des nombres complexes et en déduire des grandeurs physiques (impédance, puissance).
\item Je sais rédiger une démonstration complète (ex. formule de Moivre) en justifiant chaque étape.
\item Je sais identifier et éviter les erreurs fréquentes sur le choix de l'argument et sur la manipulation des exposants complexes.
\end{itemize}
\end{objectifs}

\begin{prerequis}
\begin{itemize}
\item Définition d'un nombre complexe (forme algébrique) et opérations de base (addition, multiplication).
\item Notions de trigonométrie : cos, sin, formules d'addition, valeurs remarquables.
\item Calcul de module et d'argument d'un vecteur dans le plan (classe de 4ᵉ).
\item Utilisation de la notation exponentielle réelle (eˣ) et propriétés des exposants.
\item Raisonnement par récurrence (classe de 1ʳᵉ).
\end{itemize}
\end{prerequis}

\newpage

\coursec{Représentation géométrique d'un nombre complexe}

\textbf{Situation d'accroche.} Au marché de Douala, un électricien doit régler l'intensité d'un circuit alimentant des lampes LED qui clignotent à une fréquence précise. Il branche son oscilloscope : l'écran affiche une sinusoïde. Pour comprendre ce signal — son amplitude, son déphasage —, le technicien doit quitter la représentation temporelle et passer à une représentation \emph{complexe} : un point du plan, repéré par une distance à l'origine et un angle. C'est exactement ce que nous allons apprendre à faire dans cette leçon.

\textbf{Problématique.} Comment repérer un nombre complexe dans le plan ? Comment passer de son écriture algébrique $a+ib$ à une écriture géométrique (distance et angle) qui simplifie les calculs ?

\courssub{Rappel : les nombres complexes et leur écriture algébrique}

L'ensemble $\mathbb{R}$ des nombres réels ne permet pas de résoudre toutes les équations : l'équation $x^2+1=0$ n'a aucune solution réelle, car le carré d'un nombre réel est toujours positif ou nul. Les mathématiciens ont alors \emph{inventé} un nouveau nombre, noté $i$, dont le carré vaut $-1$. À partir de ce nombre, on construit tous les nombres complexes.

\begin{definition}[Nombre complexe, parties réelle et imaginaire, conjugué]
On appelle \cle{nombre complexe} tout nombre $z$ qui s'écrit sous la forme
\[ z = a + ib \]
où $a$ et $b$ sont des nombres réels et $i$ est un nombre vérifiant $i^2 = -1$.
\begin{itemize}
\item Le réel $a$ est la \cle{partie réelle} de $z$, notée $a = \mathrm{Re}(z)$.
\item Le réel $b$ est la \cle{partie imaginaire} de $z$, notée $b = \mathrm{Im}(z)$. Attention : la partie imaginaire est un nombre \textbf{réel} (c'est $b$, pas $ib$).
\item L'écriture $z = a+ib$ est la \cle{forme algébrique} de $z$.
\item Le \cle{conjugué} de $z$ est le nombre complexe $\bar{z} = a - ib$.
\end{itemize}
L'ensemble des nombres complexes est noté $\mathbb{C}$. On a $\mathbb{N} \subset \mathbb{Z} \subset \mathbb{R} \subset \mathbb{C}$.
\end{definition}

Cas particuliers importants :
\begin{itemize}
\item si $b = 0$, alors $z = a$ est un \cle{nombre réel} ;
\item si $a = 0$ et $b \neq 0$, alors $z = ib$ est un \cle{imaginaire pur}.
\end{itemize}

\begin{exemplebox}[Lecture des parties réelle et imaginaire]
Soit $z = -3 + 2i$. Alors $\mathrm{Re}(z) = -3$ et $\mathrm{Im}(z) = 2$. Son conjugué est $\bar{z} = -3 - 2i$.
Soit $w = 5i$ : c'est un imaginaire pur, $\mathrm{Re}(w) = 0$ et $\mathrm{Im}(w) = 5$.
\end{exemplebox}

\begin{attention}[La partie imaginaire ne contient pas $i$]
Une erreur très fréquente au Baccalauréat technique : écrire que la partie imaginaire de $z = -3+2i$ est $2i$. C'est faux ! La partie imaginaire est le réel $2$. De même, $\mathrm{Im}(7) = 0$ et $\mathrm{Re}(4i) = 0$.
\end{attention}

\courssub{Le plan complexe : affixe d'un point, image d'un nombre}

Tout comme un réel se représente sur une droite, un complexe $z = a+ib$ se représente par un \textbf{point du plan} : il suffit de lire $a$ en abscisse et $b$ en ordonnée.

\begin{definition}[Plan complexe, affixe, image]
Le plan est muni d'un repère orthonormé direct $(O\,;\vec{u},\vec{v})$.
\begin{itemize}
\item L'axe des abscisses $(O\,;\vec{u})$ est appelé \cle{axe réel}.
\item L'axe des ordonnées $(O\,;\vec{v})$ est appelé \cle{axe imaginaire}.
\item Au nombre complexe $z = a+ib$, on associe le point $M(a\,;b)$, appelé \cle{image} de $z$.
\item Réciproquement, $z$ est appelé l'\cle{affixe} du point $M$ (et aussi l'affixe du vecteur $\overrightarrow{OM}$).
\end{itemize}
Le plan muni de cette correspondance est appelé le \cle{plan complexe}.
\end{definition}

Ainsi, les points de l'axe réel ont pour affixes des nombres réels, et les points de l'axe imaginaire ont pour affixes des imaginaires purs. Le conjugué $\bar{z}$ a pour image le symétrique de $M$ par rapport à l'axe réel.

\begin{popfigure}
\begin{center}
\begin{tikzpicture}[scale=1.6,>=Stealth]
\draw[popline!60] (-2.4,-1.4) grid (2.6,2.2);
\draw[->,thick,popdark] (-2.4,0) -- (2.7,0) node[below right] {axe réel};
\draw[->,thick,popdark] (0,-1.4) -- (0,2.35) node[above left] {axe imaginaire};
\coordinate (O) at (0,0);
\coordinate (M) at (1.5,1.732);
\draw[thick,popblue] (O) -- (M) node[midway,above left,popblue] {$|z|$};
\draw[dashed,poporange] (M) -- (1.5,0) node[below,poporange] {$a$};
\draw[dashed,poporange] (M) -- (0,1.732) node[left,poporange] {$b$};
\draw[->,popgreen,thick] (0.8,0) arc (0:49.1:0.8) node[midway,right] {$\theta$};
\fill[popink] (M) circle (1.6pt) node[above right] {$M(a\,;b)$ affixe $z$};
\fill[popdark] (O) circle (1.3pt) node[below left] {$O$};
\draw[popdark] (1,0.06) -- (1,-0.06) node[below] {$1$};
\draw[popdark] (0.06,1) -- (-0.06,1) node[left] {$i$};
\end{tikzpicture}
\end{center}
\legende{Plan complexe : le point $M$ d'affixe $z=a+ib$, le module $|z|=OM$ et l'argument $\theta$.}
\end{popfigure}

\courssub{Module d'un nombre complexe}

Regardons la figure : le segment $[OM]$ est l'hypoténuse d'un triangle rectangle dont les côtés de l'angle droit mesurent $|a|$ et $|b|$. Le théorème de Pythagore donne immédiatement la longueur $OM$.

\begin{definition}[Module d'un nombre complexe]
Soit $z = a+ib$ un nombre complexe, d'image $M$ dans le plan complexe. Le \cle{module} de $z$ est le réel positif
\[ |z| = \sqrt{a^2+b^2} = OM. \]
\end{definition}

\begin{exemplebox}[Calculs de modules]
$|3+4i| = \sqrt{9+16} = \sqrt{25} = 5$ ; \quad $|-2i| = \sqrt{0+4} = 2$ ; \quad $|5| = 5$ ; \quad $|1+i| = \sqrt{2}$.
\end{exemplebox}

\begin{propriete}[Propriétés du module]
Pour tous nombres complexes $z$, $z_1$, $z_2$ :
\begin{itemize}
\item $|z| = 0 \iff z = 0$ ;
\item $|\bar{z}| = |z|$ et $z \times \bar{z} = |z|^2$ ;
\item $|z_1 z_2| = |z_1|\,|z_2|$ (le module d'un produit est le produit des modules) ;
\item si $z_2 \neq 0$, $\left|\dfrac{z_1}{z_2}\right| = \dfrac{|z_1|}{|z_2|}$.
\end{itemize}
\end{propriete}

\begin{exemplebox}[Vérification de $z\bar{z} = |z|^2$]
Soit $z = 1 + 2i$. Alors $\bar{z} = 1 - 2i$ et
\[ z\,\bar{z} = (1+2i)(1-2i) = 1 - (2i)^2 = 1 - 4i^2 = 1 + 4 = 5 = |z|^2, \]
car $|z|^2 = 1^2 + 2^2 = 5$. Cette identité est très utile pour rendre réel un dénominateur dans les quotients.
\end{exemplebox}

\courssub{Argument d'un nombre complexe non nul}

Le module donne la \emph{distance} de $M$ à l'origine. Pour repérer complètement $M$, il faut encore une \emph{direction} : c'est l'angle que fait la demi-droite $[OM)$ avec l'axe réel.

\begin{definition}[Argument, argument principal]
Soit $z$ un nombre complexe \textbf{non nul}, d'image $M$. On appelle \cle{argument} de $z$ toute mesure $\theta$, en radians, de l'angle orienté
\[ \theta = \left(\vec{u}\,;\overrightarrow{OM}\right). \]
On note $\theta = \arg(z)$. Un nombre complexe non nul possède une infinité d'arguments, qui diffèrent tous d'un multiple de $2\pi$ : on écrit $\arg(z) \equiv \theta \pmod{2\pi}$.
Parmi tous ces arguments, celui qui appartient à l'intervalle $]-\pi\,;\pi]$ est appelé \cle{argument principal} et est noté $\mathrm{Arg}(z)$.
\end{definition}

\begin{attention}[L'argument est défini à $2\pi$ près — et $0$ n'a pas d'argument]
Deux pièges classiques :
\begin{enumerate}
\item Le nombre $z = 0$ n'a \textbf{pas d'argument} : le point $O$ ne définit aucune direction. Ne jamais chercher $\arg(0)$.
\item Si $\theta$ est un argument de $z$, alors $\theta + 2\pi$, $\theta - 2\pi$, $\theta + 4\pi$... sont aussi des arguments. Quand l'énoncé demande \emph{l'argument} de $z$, on donne toujours la valeur \textbf{principale}, dans $]-\pi\,;\pi]$. Par exemple, $\mathrm{Arg}(i) = \dfrac{\pi}{2}$ et non $\dfrac{5\pi}{2}$.
\end{enumerate}
\end{attention}

\begin{propriete}[Argument d'un produit et d'un quotient]
Pour tous nombres complexes non nuls $z_1$ et $z_2$ :
\[ \arg(z_1 z_2) \equiv \arg(z_1) + \arg(z_2) \pmod{2\pi}, \qquad \arg\!\left(\frac{z_1}{z_2}\right) \equiv \arg(z_1) - \arg(z_2) \pmod{2\pi}. \]
En particulier, $\arg(\bar{z}) \equiv -\arg(z) \pmod{2\pi}$ et, pour tout entier naturel $n$, $\arg(z^n) \equiv n\arg(z) \pmod{2\pi}$.
\end{propriete}

\begin{savaistu}[Pourquoi l'électricien aime les arguments]
En électricité, un courant alternatif est représenté par un nombre complexe : le module correspond à l'amplitude du signal et l'argument au \emph{déphasage}. Multiplier deux complexes revient à multiplier les amplitudes et à \emph{additionner les déphasages} : c'est la propriété ci-dessus qui fait toute la puissance des calculs en électrotechnique, au Cameroun comme partout dans le monde.
\end{savaistu}

\courssub{Forme trigonométrique d'un nombre complexe}

Plaçons-nous dans le triangle rectangle de la figure. Par définition du cosinus et du sinus :
\[ \cos\theta = \frac{a}{|z|} \quad \text{et} \quad \sin\theta = \frac{b}{|z|}, \]
c'est-à-dire
\[ a = |z|\cos\theta \qquad \text{et} \qquad b = |z|\sin\theta. \]
En remplaçant dans $z = a+ib$, on obtient une nouvelle écriture de $z$.

\begin{definition}[Forme trigonométrique]
Tout nombre complexe non nul $z$, de module $r = |z|$ et d'argument $\theta$, s'écrit
\[ z = r\left(\cos\theta + i\sin\theta\right). \]
Cette écriture est la \cle{forme trigonométrique} (ou forme polaire) de $z$.
\end{definition}

\begin{methode}[Conversion : forme algébrique $\rightarrow$ forme trigonométrique]
Pour écrire $z = a+ib$ (non nul) sous forme trigonométrique :
\begin{enumerate}
\item \textbf{Calculer le module} : $r = |z| = \sqrt{a^2+b^2}$.
\item \textbf{Factoriser $r$} : $z = r\left(\dfrac{a}{r} + i\dfrac{b}{r}\right)$.
\item \textbf{Identifier l'angle} : chercher $\theta$ tel que $\cos\theta = \dfrac{a}{r}$ et $\sin\theta = \dfrac{b}{r}$, en utilisant le cercle trigonométrique et les valeurs remarquables.
\item \textbf{Choisir la valeur principale} : vérifier que $\theta \in \left]-\pi\,;\pi\right]$, puis conclure $z = r(\cos\theta + i\sin\theta)$.
\end{enumerate}
Le signe de $b$ (partie imaginaire) indique si $\theta$ est positif ou négatif ; le signe de $a$ indique si $|\theta| < \dfrac{\pi}{2}$ ou $|\theta| > \dfrac{\pi}{2}$.
\end{methode}

\begin{exoresolu}[Mise sous forme trigonométrique de $z = -1 + i\sqrt{3}$]
Énoncé. Déterminer le module et l'argument principal de $z = -1 + i\sqrt{3}$, puis écrire $z$ sous forme trigonométrique.
\tcblower
\textbf{Solution détaillée.}

\textbf{Étape 1 : le module.}
\[ r = |z| = \sqrt{(-1)^2 + (\sqrt{3})^2} = \sqrt{1+3} = \sqrt{4} = 2. \]

\textbf{Étape 2 : factorisation.}
\[ z = 2\left(-\frac{1}{2} + i\,\frac{\sqrt{3}}{2}\right). \]

\textbf{Étape 3 : identification de l'angle.} On cherche $\theta$ tel que
\[ \cos\theta = -\frac{1}{2} \quad \text{et} \quad \sin\theta = \frac{\sqrt{3}}{2}. \]
Sur le cercle trigonométrique, l'angle dont le cosinus vaut $-\dfrac{1}{2}$ et le sinus $\dfrac{\sqrt{3}}{2}$ est $\theta = \dfrac{2\pi}{3}$ (deuxième quadrant, ce qui est cohérent : $a<0$ et $b>0$).

\textbf{Étape 4 : conclusion.} Comme $\dfrac{2\pi}{3} \in \left]-\pi\,;\pi\right]$, c'est bien l'argument principal :
\[ |z| = 2, \qquad \mathrm{Arg}(z) = \frac{2\pi}{3}, \qquad z = 2\left(\cos\frac{2\pi}{3} + i\sin\frac{2\pi}{3}\right). \]
\end{exoresolu}

\begin{exoresolu}[Un argument négatif : $z = \sqrt{3} - i$]
Énoncé. Écrire $z = \sqrt{3} - i$ sous forme trigonométrique.
\tcblower
\textbf{Solution.} Module : $r = \sqrt{(\sqrt{3})^2+(-1)^2} = \sqrt{3+1} = 2$. Alors
\[ z = 2\left(\frac{\sqrt{3}}{2} - \frac{i}{2}\right), \quad \text{donc} \quad \cos\theta = \frac{\sqrt{3}}{2} \ \text{et}\ \sin\theta = -\frac{1}{2}. \]
Le sinus est négatif : l'angle est dans le quatrième quadrant, $\theta = -\dfrac{\pi}{6}$, qui appartient bien à $]-\pi\,;\pi]$. D'où
\[ z = 2\left(\cos\left(-\frac{\pi}{6}\right) + i\sin\left(-\frac{\pi}{6}\right)\right), \qquad \mathrm{Arg}(z) = -\frac{\pi}{6}. \]
\end{exoresolu}

\courssub{Cas particuliers à connaître par cœur}

\begin{center}
\begin{tabularx}{\linewidth}{|l|X|X|X|}
\hline
\rowcolor{popblueL} \textbf{Nombre $z$} & \textbf{Module $|z|$} & \textbf{Argument principal} & \textbf{Position de $M$} \\
\hline
$a$ réel, $a>0$ & $a$ & $0$ & demi-axe réel positif \\
\hline
$a$ réel, $a<0$ & $-a$ & $\pi$ & demi-axe réel négatif \\
\hline
$ib$, $b>0$ & $b$ & $\dfrac{\pi}{2}$ & demi-axe imaginaire positif \\
\hline
$ib$, $b<0$ & $-b$ & $-\dfrac{\pi}{2}$ & demi-axe imaginaire négatif \\
\hline
$0$ & $0$ & non défini & origine $O$ \\
\hline
\end{tabularx}
\end{center}

\begin{exemplebox}[Lecture directe]
$|7| = 7$ et $\mathrm{Arg}(7) = 0$ ; \quad $|-3| = 3$ et $\mathrm{Arg}(-3) = \pi$ ; \quad $|2i| = 2$ et $\mathrm{Arg}(2i) = \dfrac{\pi}{2}$ ; \quad $\mathrm{Arg}(-i) = -\dfrac{\pi}{2}$.
\end{exemplebox}

\begin{attention}[Rédaction : ne pas confondre les deux formes]
La forme algébrique $a+ib$ et la forme trigonométrique $r(\cos\theta + i\sin\theta)$ décrivent le \textbf{même} nombre. Mais attention : dans $r(\cos\theta + i\sin\theta)$, le nombre $r$ doit être \textbf{strictement positif} (c'est un module). Écrire $z = -2(\cos\theta + i\sin\theta)$ n'est pas une forme trigonométrique : il faut écrire $z = 2(\cos(\theta+\pi) + i\sin(\theta+\pi))$.
\end{attention}

\begin{exercice}[À faire]
Pour chacun des nombres suivants, calculer le module, déterminer l'argument principal, puis donner la forme trigonométrique :
\[ z_1 = 1+i, \qquad z_2 = -2i, \qquad z_3 = -\sqrt{3}+i, \qquad z_4 = 1 - i\sqrt{3}. \]
Placer ensuite les images de ces quatre nombres dans un plan complexe.
\end{exercice}

\begin{corrige}[Exercice]
$z_1 = 1+i$ : $|z_1| = \sqrt{1^2+1^2} = \sqrt{2}$, $\cos\theta = \dfrac{1}{\sqrt{2}} = \dfrac{\sqrt{2}}{2}$ et $\sin\theta = \dfrac{\sqrt{2}}{2}$, donc $\mathrm{Arg}(z_1) = \dfrac{\pi}{4}$ et $z_1 = \sqrt{2}\left(\cos\dfrac{\pi}{4}+i\sin\dfrac{\pi}{4}\right)$.

$z_2 = -2i$ : imaginaire pur, $|z_2| = 2$, $\mathrm{Arg}(z_2) = -\dfrac{\pi}{2}$, donc $z_2 = 2\left(\cos\left(-\dfrac{\pi}{2}\right)+i\sin\left(-\dfrac{\pi}{2}\right)\right)$.

$z_3 = -\sqrt{3}+i$ : $|z_3| = \sqrt{3+1} = 2$, $\cos\theta = -\dfrac{\sqrt{3}}{2}$ et $\sin\theta = \dfrac{1}{2}$, donc $\mathrm{Arg}(z_3) = \dfrac{5\pi}{6}$ et $z_3 = 2\left(\cos\dfrac{5\pi}{6}+i\sin\dfrac{5\pi}{6}\right)$.

$z_4 = 1-i\sqrt{3}$ : $|z_4| = \sqrt{1+3} = 2$, $\cos\theta = \dfrac{1}{2}$ et $\sin\theta = -\dfrac{\sqrt{3}}{2}$, donc $\mathrm{Arg}(z_4) = -\dfrac{\pi}{3}$ et $z_4 = 2\left(\cos\left(-\dfrac{\pi}{3}\right)+i\sin\left(-\dfrac{\pi}{3}\right)\right)$.
\end{corrige}

\begin{aretenir}[L'essentiel de la section]
\begin{itemize}
\item Tout complexe $z = a+ib$ est représenté par le point $M(a\,;b)$ du plan complexe ; $z$ est l'affixe de $M$.
\item Le \cle{module} $|z| = \sqrt{a^2+b^2}$ est la distance $OM$.
\item L'\cle{argument} $\theta = (\vec{u}\,;\overrightarrow{OM})$ est défini à $2\pi$ près, pour $z \neq 0$ ; la valeur dans $]-\pi\,;\pi]$ est l'argument principal.
\item Forme trigonométrique : $z = |z|(\cos\theta + i\sin\theta)$, avec $a = |z|\cos\theta$ et $b = |z|\sin\theta$.
\item Module et argument se comportent bien avec les produits : $|z_1 z_2| = |z_1||z_2|$ et $\arg(z_1z_2) \equiv \arg(z_1)+\arg(z_2) \pmod{2\pi}$.
\end{itemize}
\end{aretenir}

\coursec{Forme exponentielle (Euler) et calculs simplifiés}

Dans la section précédente, nous avons vu que tout nombre complexe non nul $z$ peut être représenté par un point $M$ du plan, repéré par son \cle{module} $r = |z|$ (la distance $OM$) et son \cle{argument} $\theta = \arg(z)$ (l'angle orienté entre l'axe des réels positifs et le vecteur $\overrightarrow{OM}$). Nous avons ainsi obtenu la \cle{forme trigonométrique} :
\[ z = r\bigl(\cos\theta + i\sin\theta\bigr). \]
Cette écriture est déjà très utile, mais elle reste lourde à manipuler : multiplier deux formes trigonométriques exige les formules d'addition de $\cos$ et $\sin$. Dans cette section, nous allons condenser l'expression $\cos\theta + i\sin\theta$ en une seule notation, $e^{i\theta}$, grâce à la célèbre \cle{formule d'Euler}. Le gain est spectaculaire : multiplier des complexes reviendra à \emph{additionner des angles}, et élever à la puissance $n$ reviendra à \emph{multiplier l'angle par $n$}. C'est exactement l'outil utilisé par l'électrotechnicien pour manipuler les impédances et les déphasages.

\courssub{La formule d'Euler}

\textbf{Intuition.} Sur le cercle trigonométrique, le point d'angle $\theta$ a pour coordonnées $(\cos\theta\,;\,\sin\theta)$, donc pour affixe $\cos\theta + i\sin\theta$. Euler a eu l'idée audacieuse de noter ce nombre comme une exponentielle, $e^{i\theta}$, car il se comporte \emph{exactement} comme une exponentielle : quand on multiplie deux tels nombres, les angles s'additionnent, tout comme les exposants s'additionnent pour les puissances ordinaires.

\begin{propriete}[Formule d'Euler (admise)]
Pour tout réel $\theta$ (exprimé en radians) :
\[ \boxed{\ e^{i\theta} = \cos\theta + i\sin\theta\ } \]
En particulier, $|e^{i\theta}| = \sqrt{\cos^2\theta + \sin^2\theta} = 1$ : le nombre $e^{i\theta}$ est toujours sur le \cle{cercle unité}.
\end{propriete}

\begin{savaistu}[D'où vient cette formule ?]
Au programme de Terminale technique, la formule d'Euler est \textbf{admise}. Pour la culture, sachez qu'elle se justifie à l'aide des \emph{séries de Maclaurin} : on peut écrire les fonctions usuelles comme des sommes infinies de puissances :
\[ e^{x} = 1 + x + \frac{x^2}{2!} + \frac{x^3}{3!} + \cdots \]
En remplaçant formellement $x$ par $i\theta$ et en regroupant les puissances paires (qui donnent $\cos\theta$) et impaires (qui donnent $i\sin\theta$), on obtient $e^{i\theta} = \cos\theta + i\sin\theta$. C'est Leonhard Euler (1707--1783) qui a dégagé cette identité, souvent citée comme « la plus belle formule des mathématiques » lorsqu'on y pose $\theta = \pi$ : $e^{i\pi} + 1 = 0$, qui relie les cinq constantes fondamentales $0$, $1$, $i$, $\pi$ et $e$.
\end{savaistu}

\begin{definition}[Forme exponentielle d'un nombre complexe]
Tout nombre complexe non nul $z$, de module $r = |z| > 0$ et d'argument $\theta = \arg(z)$, s'écrit sous \cle{forme exponentielle} :
\[ z = r\,e^{i\theta}. \]
\end{definition}

\begin{exemplebox}[Valeurs remarquables]
À l'aide des valeurs usuelles de $\cos$ et $\sin$ :
\[ e^{i0} = 1, \qquad e^{i\pi/2} = i, \qquad e^{i\pi} = -1, \qquad e^{i\pi/6} = \frac{\sqrt{3}}{2} + \frac{1}{2}i, \qquad e^{i\pi/4} = \frac{\sqrt{2}}{2} + \frac{\sqrt{2}}{2}\,i. \]
Ainsi $1 + i = \sqrt{2}\,e^{i\pi/4}$ et $\sqrt{3} + i = 2\,e^{i\pi/6}$.
\end{exemplebox}

\begin{popfigure}
\begin{center}
\begin{tikzpicture}[scale=2.6]
  % axes
  \draw[->,popmuted] (-1.35,0) -- (1.45,0) node[below right] {\small Ré};
  \draw[->,popmuted] (0,-1.25) -- (0,1.35) node[above left] {\small Im};
  % cercle unité
  \draw[popblue,thick] (0,0) circle (1);
  % points remarquables
  \foreach \a/\lab in {0/{e^{i0}=1}, 30/{e^{i\pi/6}}, 45/{e^{i\pi/4}}, 60/{e^{i\pi/3}}, 90/{e^{i\pi/2}=i}}{
    \draw[popline] (0,0) -- (\a:1);
    \fill[poporange] (\a:1) circle (0.025);
    \node[popdark] at (\a:1.18) {\small $\lab$};
  }
  % angle exemple
  \draw[popgreen,thick,->] (0.45,0) arc (0:45:0.45);
  \node[popgreen] at (22.5:0.62) {\small $\theta$};
  \node[popmuted] at (-0.08,-0.1) {\small $O$};
\end{tikzpicture}
\end{center}
\legende{Le cercle unité : pour tout réel $\theta$, le point d'affixe $e^{i\theta} = \cos\theta + i\sin\theta$ est le point du cercle de centre $O$ et de rayon $1$ repéré par l'angle $\theta$.}
\end{popfigure}

\courssub{Produit, quotient et conjugué en forme exponentielle}

Voici la raison d'être de la forme exponentielle : les opérations deviennent mécaniques.

\begin{propriete}[Produit et quotient en forme exponentielle]
Pour tous réels $\theta$ et $\theta'$, et tous réels strictement positifs $r$ et $r'$ :
\[ \bigl(r e^{i\theta}\bigr)\bigl(r' e^{i\theta'}\bigr) = rr'\,e^{i(\theta+\theta')}, \qquad
\frac{r e^{i\theta}}{r' e^{i\theta'}} = \frac{r}{r'}\,e^{i(\theta-\theta')}, \qquad
\overline{r e^{i\theta}} = r\,e^{-i\theta}. \]
Autrement dit : \textbf{les modules se multiplient (ou se divisent), les arguments s'additionnent (ou se soustraient)}.
\end{propriete}

\begin{experience}[Démonstration guidée du produit]
Démontrons que $e^{i\theta}\,e^{i\theta'} = e^{i(\theta+\theta')}$ à l'aide de la formule d'Euler et des formules d'addition trigonométriques (vues en Première).
\begin{enumerate}
  \item Écris le produit sous forme trigonométrique :
  \[ e^{i\theta}\,e^{i\theta'} = (\cos\theta + i\sin\theta)(\cos\theta' + i\sin\theta'). \]
  \item Développe en utilisant $i^2 = -1$ :
  \[ = \bigl(\cos\theta\cos\theta' - \sin\theta\sin\theta'\bigr) + i\bigl(\sin\theta\cos\theta' + \cos\theta\sin\theta'\bigr). \]
  \item Reconnais les formules d'addition : $\cos\theta\cos\theta' - \sin\theta\sin\theta' = \cos(\theta+\theta')$ et $\sin\theta\cos\theta' + \cos\theta\sin\theta' = \sin(\theta+\theta')$.
  \item Conclus : $e^{i\theta}e^{i\theta'} = \cos(\theta+\theta') + i\sin(\theta+\theta') = e^{i(\theta+\theta')}$. \hfill$\square$
\end{enumerate}
\end{experience}

\begin{attention}[Rester dans l'intervalle de l'argument principal]
Après chaque addition ou soustraction d'arguments, le résultat peut sortir de l'intervalle $]-\pi\,;\,\pi]$. Il faut alors \textbf{ramener l'argument principal} en ajoutant ou retranchant $2\pi$ autant de fois que nécessaire. Par exemple :
\[ e^{i\,7\pi/6} = e^{i(7\pi/6 - 2\pi)} = e^{-i\,5\pi/6}. \]
Écrire $\arg(z) = \dfrac{7\pi}{6}$ comme argument \emph{principal} est une erreur fréquente au Baccalauréat technique : vérifie toujours que ton angle final appartient à $]-\pi\,;\,\pi]$.
\end{attention}

\courssub{Passage d'une forme à l'autre : la méthode}

\begin{methode}[Passage forme algébrique $\rightarrow$ forme exponentielle]
Pour écrire $z = a + ib$ (non nul) sous la forme $r e^{i\theta}$ :
\begin{enumerate}
  \item Calculer le module : $r = |z| = \sqrt{a^2 + b^2}$.
  \item Factoriser $r$ : $z = r\left(\dfrac{a}{r} + i\dfrac{b}{r}\right)$.
  \item Identifier l'angle $\theta$ tel que $\cos\theta = \dfrac{a}{r}$ et $\sin\theta = \dfrac{b}{r}$, en utilisant le cercle trigonométrique et le \emph{signe} de $a$ et $b$ pour choisir le bon quadrant.
  \item Conclure : $z = r e^{i\theta}$, avec $\theta$ ramené dans $]-\pi\,;\,\pi]$.
\end{enumerate}
\textbf{Sens inverse} (exponentielle $\rightarrow$ algébrique) : développer $r e^{i\theta} = r\cos\theta + i\,r\sin\theta$ à l'aide des valeurs remarquables.
\end{methode}

\begin{exoresolu}[Mise sous forme exponentielle et produit]
\textbf{Énoncé.} On pose $z_1 = 1 + i$ et $z_2 = \sqrt{3} + i$.
\begin{enumerate}
  \item Écrire $z_1$ et $z_2$ sous forme exponentielle.
  \item En déduire la forme exponentielle du produit $z = z_1 z_2$.
\end{enumerate}
\tcblower
\textbf{Solution détaillée.}
\textbf{1.} Pour $z_1 = 1+i$ : $|z_1| = \sqrt{1^2+1^2} = \sqrt{2}$, puis
\[ z_1 = \sqrt{2}\left(\frac{\sqrt{2}}{2} + i\frac{\sqrt{2}}{2}\right) = \sqrt{2}\,e^{i\pi/4}. \]
Pour $z_2 = \sqrt{3}+i$ : $|z_2| = \sqrt{3+1} = 2$, puis
\[ z_2 = 2\left(\frac{\sqrt{3}}{2} + i\,\frac{1}{2}\right) = 2\,e^{i\pi/6}. \]
\textbf{2.} Les modules se multiplient, les arguments s'additionnent :
\[ z = \sqrt{2}\,e^{i\pi/4} \times 2\,e^{i\pi/6} = 2\sqrt{2}\,e^{i\left(\frac{\pi}{4}+\frac{\pi}{6}\right)} = 2\sqrt{2}\,e^{i\,5\pi/12}. \]
L'angle $\dfrac{5\pi}{12}$ appartient bien à $]-\pi\,;\,\pi]$ : c'est terminé.
\end{exoresolu}

\courssub{Puissances entières : la formule de Moivre}

\begin{propriete}[Formule de Moivre]
Pour tout entier naturel $n$ et tout réel $\theta$ :
\[ \bigl(r e^{i\theta}\bigr)^n = r^n e^{in\theta}, \qquad\text{en particulier}\qquad \bigl(\cos\theta + i\sin\theta\bigr)^n = \cos(n\theta) + i\sin(n\theta). \]
Cette formule s'obtient en appliquant $n$ fois la règle du produit. Elle sera exploitée dans la section 3 pour la \emph{linéarisation}.
\textbf{Remarque :} elle reste valable pour $n \in \mathbb{Z}$ (entier relatif) si $z \neq 0$.
\end{propriete}

\begin{exemplebox}[Exemple chiffré détaillé : calcul de $(\sqrt{3}+i)^4$]
On a vu que $\sqrt{3} + i = 2\,e^{i\pi/6}$. Donc :
\begin{align*}
(\sqrt{3}+i)^4 &= \bigl(2\,e^{i\pi/6}\bigr)^4 = 2^4\,e^{i\,4\pi/6} = 16\,e^{i\,2\pi/3}.
\end{align*}
Repassons à la forme algébrique : $\cos\dfrac{2\pi}{3} = -\dfrac{1}{2}$ et $\sin\dfrac{2\pi}{3} = \dfrac{\sqrt{3}}{2}$, d'où
\[ (\sqrt{3}+i)^4 = 16\left(-\frac{1}{2} + i\,\frac{\sqrt{3}}{2}\right) = -8 + 8i\sqrt{3}. \]
\textit{Remarque :} si l'on avait obtenu un angle hors de $]-\pi\,;\,\pi]$, par exemple avec $(\sqrt3+i)^8 = 256\,e^{i\,4\pi/3}$, on ramènerait l'argument principal : $e^{i\,4\pi/3} = e^{-i\,2\pi/3}$, ce qui donne $-128 - 128i\sqrt{3}$.
\end{exemplebox}

\begin{attention}[Pièges fréquents avec les puissances]
\begin{itemize}
  \item Ne pas oublier d'élever \textbf{aussi le module} à la puissance $n$ : $(2e^{i\pi/6})^4 = 16\,e^{i\,2\pi/3}$ et non $2\,e^{i\,2\pi/3}$.
  \item La forme exponentielle exige un module \emph{positif} : $-2e^{i\pi/6}$ n'est pas une forme exponentielle ; il faut écrire $-2e^{i\pi/6} = 2\,e^{i(\pi/6+\pi)} = 2\,e^{i\,7\pi/6} = 2\,e^{-i\,5\pi/6}$.
  \item Développer $(a+ib)^n$ avec le binôme de Newton pour $n = 4$ ou $5$ est long et source d'erreurs : la forme exponentielle est presque toujours plus rapide et plus sûre.
\end{itemize}
\end{attention}

\courssub{Racines n-ièmes d'un nombre complexe non nul}

\begin{definition}[Racine n-ième]
Soit $n \geq 2$ un entier. Une \cle{racine n-ième} du complexe non nul $Z$ est tout complexe $z$ tel que $z^n = Z$.
\end{definition}

\begin{propriete}[Les n racines n-ièmes (admise)]
Tout complexe non nul $Z = R\,e^{i\varphi}$ admet \textbf{exactement $n$ racines n-ièmes distinctes}, données par :
\[ z_k = R^{1/n}\, e^{\,i\frac{\varphi + 2k\pi}{n}}, \qquad k = 0,\,1,\,\dots,\,n-1. \]
Ces $n$ points sont les sommets d'un \cle{polygone régulier} à $n$ côtés inscrit dans le cercle de centre $O$ et de rayon $R^{1/n}$.
\end{propriete}

\begin{exemplebox}[Racines cubiques de $8i$]
On écrit $8i = 8\,e^{i\pi/2}$. Les trois racines cubiques sont :
\[ z_k = 8^{1/3}\, e^{\,i\frac{\pi/2 + 2k\pi}{3}} = 2\,e^{\,i\left(\frac{\pi}{6} + \frac{2k\pi}{3}\right)}, \quad k = 0,1,2, \]
soit $z_0 = 2e^{i\pi/6} = \sqrt{3} + i$, $z_1 = 2e^{i\,5\pi/6} = -\sqrt{3} + i$ et $z_2 = 2e^{-i\pi/2} = -2i$. Vérification : $z_0^3 = 8e^{i\pi/2} = 8i$. Les trois points forment un triangle équilatéral sur le cercle de rayon $2$.
\end{exemplebox}

\courssub{Applications techniques}

\begin{exoresolu}[Impédance d'un circuit en électrotechnique]
\textbf{Énoncé.} En courant alternatif, l'impédance complexe d'un dipôle s'écrit $\underline{Z} = R + jX$ (les électriciens notent $j$ le nombre $i$). Un circuit série comporte une résistance $R = \SI{30}{\ohm}$ et une réactance inductive $X = \SI{40}{\ohm}$.
\begin{enumerate}
  \item Écrire $\underline{Z}$ sous forme exponentielle (on donnera l'argument en degrés, arrondi au dixième).
  \item Deux dipôles identiques sont associés de sorte que l'impédance équivalente soit $\underline{Z}^2 / 100$. Donner son module.
\end{enumerate}
\tcblower
\textbf{Solution détaillée.}
\textbf{1.} Module : $|\underline{Z}| = \sqrt{30^2 + 40^2} = \sqrt{2500} = \SI{50}{\ohm}$. L'argument $\varphi$ vérifie $\cos\varphi = \dfrac{30}{50} = 0{,}6$ et $\sin\varphi = \dfrac{40}{50} = 0{,}8$, d'où $\varphi \approx 53{,}1^\circ$ (calculatrice en mode degré). Donc $\underline{Z} \approx 50\,e^{i\,53{,}1^\circ}$.
\textbf{2.} Par la formule de Moivre, $\underline{Z}^2 = 50^2 e^{i\,2\varphi} = 2500\,e^{i\,106{,}2^\circ}$, donc $\left|\dfrac{\underline{Z}^2}{100}\right| = \dfrac{2500}{100} = \SI{25}{\ohm}$.
\textit{Interprétation métier :} le module de l'impédance donne l'opposition du circuit au passage du courant, et l'argument donne le déphasage entre tension et courant — deux grandeurs que le technicien mesure et règle quotidiennement.
\end{exoresolu}

\begin{savaistu}[De Euler au traitement du signal]
La forme exponentielle $e^{i\theta}$ est la brique de base de la \cle{transformée de Fourier}, outil mathématique qui décompose tout signal (son, image, tension électrique) en une somme de signaux sinusoïdaux. C'est elle qui permet la compression des fichiers MP3 et JPEG, l'imagerie médicale (IRM), et l'analyse des vibrations des machines tournantes dans les ateliers. Au Cameroun, les techniciens des télécommunications et du génie électrique l'utilisent chaque jour, souvent sans même s'en rendre compte, à travers les logiciels de mesure.
\end{savaistu}

\begin{aretenir}[L'essentiel de la section]
\begin{itemize}
  \item Formule d'Euler : $e^{i\theta} = \cos\theta + i\sin\theta$ ; forme exponentielle : $z = r e^{i\theta}$ avec $r = |z|$, $\theta = \arg(z)$.
  \item Produit : les modules se multiplient, les arguments s'additionnent. Quotient : les modules se divisent, les arguments se soustraient.
  \item Moivre : $(r e^{i\theta})^n = r^n e^{in\theta}$.
  \item Racines n-ièmes de $R e^{i\varphi}$ : $z_k = R^{1/n} e^{i(\varphi + 2k\pi)/n}$, $k = 0,\dots,n-1$ (polygone régulier).
  \item Toujours ramener l'argument principal dans $]-\pi\,;\,\pi]$.
\end{itemize}
\end{aretenir}

\coursec{Linéarisation (formules de Moivre) et puissances trigonométriques}

\courssub{Formule de Moivre et démonstration par récurrence}

Dans la section précédente, nous avons vu que tout nombre complexe $z$ de module $1$ s'écrit sous forme exponentielle $z=e^{i\theta}=\cos\theta+i\sin\theta$. Élever $z$ à une puissance entière $n$ revient à multiplier l'argument par $n$ : $z^n=e^{in\theta}=\cos(n\theta)+i\sin(n\theta)$. C'est précisément le contenu de la \cle{formule de Moivre}, outil central pour passer des puissances de fonctions trigonométriques à des combinaisons linéaires de cosinus et sinus d'angles multiples.

\begin{propriete}[Formule de Moivre (démonstration exigible au Probatoire)]
Pour tout entier naturel $n$ et tout réel $\theta$,
\[
(\cos\theta+i\sin\theta)^n = \cos(n\theta)+i\sin(n\theta).
\]
\end{propriete}

\begin{experience}[Démonstration par récurrence sur $n$]
\emph{Initialisation ($n=1$) :} L'égalité est triviale : $(\cos\theta+i\sin\theta)^1 = \cos\theta+i\sin\theta = \cos(1\cdot\theta)+i\sin(1\cdot\theta)$.
\par\emph{Hérédité :} Supposons la formule vraie pour un certain $n\ge 1$, c'est-à-dire
$(\cos\theta+i\sin\theta)^n = \cos(n\theta)+i\sin(n\theta)$.
Alors
\begin{align*}
(\cos\theta+i\sin\theta)^{n+1}
&= (\cos\theta+i\sin\theta)^n(\cos\theta+i\sin\theta) \\
&= \bigl(\cos(n\theta)+i\sin(n\theta)\bigr)(\cos\theta+i\sin\theta) \quad\text{(hypothèse de récurrence)} \\
&= \cos(n\theta)\cos\theta - \sin(n\theta)\sin\theta \\
&\quad + i\bigl(\sin(n\theta)\cos\theta+\cos(n\theta)\sin\theta\bigr) \\
&= \cos\bigl((n+1)\theta\bigr) + i\sin\bigl((n+1)\theta\bigr) \quad\text{(formules d'addition)}.
\end{align*}
La propriété est donc vraie pour $n+1$. Par récurrence, elle est vraie pour tout $n\in\mathbb{N}$.
\end{experience}

\begin{savaistu}[Contexte historique]
Abraham de Moivre (1667--1754), mathématicien français exilé en Angleterre, a découvert cette formule vers 1707. Elle relie l'algèbre des complexes à la trigonométrie et préfigure la formule d'Euler $e^{i\theta}=\cos\theta+i\sin\theta$ (1748). Au Probatoire, la démonstration par récurrence est un classique : soignez la rédaction (initialisation, hérédité, conclusion).
\end{savaistu}

\courssub{Identification des parties réelle et imaginaire}

En développant $(\cos\theta+i\sin\theta)^n$ avec le binôme de Newton et en identifiant les parties réelle et imaginaire, on obtient des expressions polynomiales de $\cos(n\theta)$ et $\sin(n\theta)$ en puissances de $\cos\theta$ et $\sin\theta$. Réciproquement, ces identités permettent de \cle{linéariser} des puissances trigonométriques, c'est-à-dire de les réécrire comme combinaisons linéaires de $\cos(k\theta)$ et $\sin(k\theta)$.

\begin{propriete}[Expressions de $\cos(n\theta)$ et $\sin(n\theta)$ via le binôme de Newton]
Pour $n\in\mathbb{N}^*$,
\[
\cos(n\theta)+i\sin(n\theta)=\sum_{k=0}^{n}\binom{n}{k}\cos^{n-k}\theta\,(i\sin\theta)^k.
\]
La partie réelle (somme sur $k$ pair) donne $\cos(n\theta)$, la partie imaginaire (somme sur $k$ impair) donne $\sin(n\theta)$.
\end{propriete}

\courssub{Linéarisation des puissances trigonométriques}

L'idée directrice : écrire $\cos\theta$ et $\sin\theta$ sous forme exponentielle (formules d'Euler), élever à la puissance $n$, développer avec le binôme de Newton, puis regrouper les termes conjugués pour ne garder que des cosinus ou sinus d'angles multiples.

\begin{methode}[Linéarisation d'une puissance trigonométrique $\cos^n\theta$ ou $\sin^n\theta$]
\emph{Étape 1 :} Exprimer $\cos\theta$ et $\sin\theta$ avec $e^{i\theta}$ :
$\cos\theta=\dfrac{e^{i\theta}+e^{-i\theta}}{2}$, $\quad \sin\theta=\dfrac{e^{i\theta}-e^{-i\theta}}{2i}$.
\par\emph{Étape 2 :} Élever à la puissance $n$ et développer avec le binôme de Newton :
$(a+b)^n = \sum_{k=0}^n \binom{n}{k} a^{n-k}b^k$.
\par\emph{Étape 3 :} Regrouper les termes $e^{ik\theta}$ et $e^{-ik\theta}$ pour obtenir $2\cos(k\theta)$ ou $2i\sin(k\theta)$.
\par\emph{Étape 4 :} Simplifier les coefficients (puissances de $i$, facteurs $2$, $2i$, etc.).
\end{methode}

\begin{propriete}[Formules de linéarisation usuelles (à connaître)]
\begin{align*}
\cos^2\theta &= \tfrac12(1+\cos2\theta), &
\sin^2\theta &= \tfrac12(1-\cos2\theta), \\
\cos^3\theta &= \tfrac14(3\cos\theta+\cos3\theta), &
\sin^3\theta &= \tfrac14(3\sin\theta-\sin3\theta), \\
\cos^4\theta &= \tfrac18(3+4\cos2\theta+\cos4\theta), &
\sin^4\theta &= \tfrac18(3-4\cos2\theta+\cos4\theta).
\end{align*}
\end{propriete}

\begin{exemplebox}[Linéarisation de $\cos^3\theta$ et $\sin^4\theta$]
\emph{$\cos^3\theta$ :}
\begin{align*}
\cos^3\theta &= \left(\frac{e^{i\theta}+e^{-i\theta}}{2}\right)^3
= \frac{1}{8}\bigl(e^{3i\theta}+3e^{i\theta}+3e^{-i\theta}+e^{-3i\theta}\bigr) \\
&= \frac{1}{8}\bigl[(e^{3i\theta}+e^{-3i\theta})+3(e^{i\theta}+e^{-i\theta})\bigr] \\
&= \frac{1}{4}\bigl(\cos3\theta+3\cos\theta\bigr).
\end{align*}
\emph{$\sin^4\theta$ :}
\begin{align*}
\sin^4\theta &= \left(\frac{e^{i\theta}-e^{-i\theta}}{2i}\right)^4
= \frac{1}{16}\bigl(e^{4i\theta}-4e^{2i\theta}+6-4e^{-2i\theta}+e^{-4i\theta}\bigr) \\
&= \frac{1}{16}\bigl[(e^{4i\theta}+e^{-4i\theta})-4(e^{2i\theta}+e^{-2i\theta})+6\bigr] \\
&= \frac{1}{8}\bigl(\cos4\theta-4\cos2\theta+3\bigr)
= \frac{1}{8}\bigl(3-4\cos2\theta+\cos4\theta\bigr).
\end{align*}
\end{exemplebox}

\begin{attention}[Puissance vs multiple d'angle -- Piège fréquent au Probatoire]
Ne confondez jamais $\cos^n\theta$ (puissance $n$-ième du cosinus) avec $\cos(n\theta)$ (cosinus de l'angle multiplié par $n$). Par exemple, $\cos^2\theta\ne\cos2\theta$ ; la linéarisation montre que $\cos^2\theta=\frac12(1+\cos2\theta)$.
\par De même, $\sin^3\theta \ne \sin3\theta$ ; on a $\sin^3\theta = \frac14(3\sin\theta-\sin3\theta)$.
\par \textbf{Repère :} la linéarisation \emph{abaisse} la puissance au prix d'\emph{augmenter} la fréquence (angles multiples).
\end{attention}

\begin{exoresolu}[Linéariser $\sin^5\theta$]
\emph{Énoncé :} Exprimer $\sin^5\theta$ comme combinaison linéaire de $\sin\theta$, $\sin3\theta$, $\sin5\theta$.
\par\emph{Solution :}
\begin{align*}
\sin^5\theta &= \left(\frac{e^{i\theta}-e^{-i\theta}}{2i}\right)^5 \\
&= \frac{1}{32\,i^5}\bigl(e^{5i\theta}-5e^{3i\theta}+10e^{i\theta}-10e^{-i\theta}+5e^{-3i\theta}-e^{-5i\theta}\bigr) \\
&= \frac{1}{32\,i}\bigl(e^{5i\theta}-5e^{3i\theta}+10e^{i\theta}-10e^{-i\theta}+5e^{-3i\theta}-e^{-5i\theta}\bigr) \quad(i^5=i) \\
&= \frac{1}{32\,i}\Bigl[\bigl(e^{5i\theta}-e^{-5i\theta}\bigr)-5\bigl(e^{3i\theta}-e^{-3i\theta}\bigr)+10\bigl(e^{i\theta}-e^{-i\theta}\bigr)\Bigr] \\
&= \frac{1}{32\,i}\Bigl[2i\sin5\theta -5\cdot2i\sin3\theta +10\cdot2i\sin\theta\Bigr] \\
&= \frac{1}{16}\bigl(\sin5\theta-5\sin3\theta+10\sin\theta\bigr).
\end{align*}
Donc $\displaystyle\sin^5\theta = \frac{1}{16}\bigl(10\sin\theta-5\sin3\theta+\sin5\theta\bigr)$.
\end{exoresolu}

\begin{savaistu}[Applications techniques : électricité et intégrales]
\emph{Électricité (courant alternatif) :} La puissance instantanée $p(t)=u(t)i(t)$ fait intervenir des produits $\cos^2\omega t$ ou $\sin^2\omega t$. Leur linéarisation donne la puissance moyenne (composante continue) et les harmoniques. Exemple : $u(t)=U\sqrt2\cos\omega t$, $i(t)=I\sqrt2\cos(\omega t-\varphi)$ $\Rightarrow$ $p(t)=UI[\cos\varphi+\cos(2\omega t-\varphi)]$.
\par\emph{Intégrales :} Ces formules transforment l'intégration de puissances trigonométriques en intégrales élémentaires. Par exemple,
\[
\int\cos^4\theta\,d\theta = \int\frac{3+4\cos2\theta+\cos4\theta}{8}\,d\theta
= \frac{3\theta}{8}+\frac{\sin2\theta}{4}+\frac{\sin4\theta}{32}+C,
\]
sans formules de réduction ni intégration par parties.
\end{savaistu}

\courssub{Illustration graphique : vérification de la linéarisation de $\cos^3\theta$}

Le graphique ci-dessous superpose la courbe de $y=\cos^3\theta$ (en trait plein) et celle de sa linéarisation $y=\frac14(3\cos\theta+\cos3\theta)$ (en trait pointillé) sur l'intervalle $[0,2\pi]$. Les deux courbes sont confondues, confirmant l'égalité analytique.

\begin{popfigure}
\begin{tikzpicture}
\begin{axis}[
    width=\linewidth,
    height=7cm,
    domain=0:6.283185307179586,
    samples=200,
    trig format=rad,
    axis lines=middle,
    xlabel={$\theta$},
    ylabel={$y$},
    xtick={0,1.5707963267948966,3.141592653589793,4.71238898038469,6.283185307179586},
    xticklabels={$0$,$\frac{\pi}{2}$,$\pi$,$\frac{3\pi}{2}$,$2\pi$},
    ytick={-1,-0.5,0,0.5,1},
    ymin=-1.2,ymax=1.2,
    legend pos=south east,
    grid=major,
    grid style={popline!50},
]
\addplot[popcurve,thick] {cos(x)^3};
\addlegendentry{$\cos^3\theta$}
\addplot[popcurve,dashed,thick] {(3*cos(x)+cos(3*x))/4};
\addlegendentry{$\frac14(3\cos\theta+\cos3\theta)$}
\end{axis}
\end{tikzpicture}
\legende{Superposition de $\cos^3\theta$ et de sa linéarisation sur $[0,2\pi]$.}
\end{popfigure}

Cette visualisation montre concrètement que la linéarisation n'est pas une approximation mais une identité exacte, valable pour tout $\theta\in\mathbb{R}$. Elle illustre aussi la périodicité ($2\pi$) et la symétrie (paire) des fonctions obtenues.

\coursec{Racines n-ièmes d'un nombre complexe non nul}

Dans la section précédente, nous avons appris à \cle{élever} un nombre complexe à la puissance $n$ grâce à la formule de Moivre : si $z = \rho\,e^{i\theta}$, alors $z^n = \rho^n e^{in\theta}$. Nous allons maintenant faire le chemin \emph{en sens inverse} : étant donné un nombre complexe non nul $a$, trouver \textbf{tous} les nombres complexes $z$ tels que $z^n = a$. C'est exactement le même réflexe qu'en électricité : connaissant la puissance $n$-ième d'une grandeur, on veut remonter à la grandeur elle-même. Ce problème, appelé recherche des \cle{racines n-ièmes}, possède une réponse remarquablement élégante : il y a toujours exactement $n$ solutions, et elles forment un polygone régulier dans le plan complexe.

\courssub{Le problème : résoudre l'équation $z^n = a$}

\begin{definition}[Racine n-ième d'un nombre complexe]
Soit $a$ un nombre complexe \textbf{non nul} et $n$ un entier naturel, $n \geq 2$. On appelle \cle{racine n-ième} de $a$ tout nombre complexe $z$ vérifiant :
\[ z^n = a. \]
\end{definition}

Commençons par l'intuition. Dans $\mathbb{R}$, l'équation $x^2 = 4$ a deux solutions ($2$ et $-2$), mais l'équation $x^2 = -4$ n'en a aucune. Dans $\mathbb{C}$, la situation est bien plus riche et plus régulière : l'équation $z^n = a$ (avec $a \neq 0$) possède \textbf{toujours exactement $n$ solutions distinctes}. Par exemple, l'équation $z^2 = -4$ admet les deux solutions $2i$ et $-2i$, car $(2i)^2 = 4i^2 = -4$.

Pour trouver ces solutions, l'outil décisif est la \cle{forme exponentielle}. Écrivons le nombre $a$ sous forme exponentielle :
\[ a = \rho\, e^{i\alpha}, \qquad \text{avec } \rho = |a| > 0 \ \text{et}\ \alpha = \arg(a) \ [2\pi]. \]
Cherchons $z$ lui aussi sous forme exponentielle : $z = r\,e^{i\theta}$ avec $r > 0$. L'équation $z^n = a$ devient, grâce à la formule de Moivre :
\[ r^n e^{in\theta} = \rho\, e^{i\alpha}. \]
Deux nombres complexes non nuls écrits sous forme exponentielle sont égaux si et seulement si ils ont le \textbf{même module} et des \textbf{arguments égaux modulo $2\pi$}. On obtient donc :
\[ r^n = \rho \qquad \text{et} \qquad n\theta = \alpha + 2k\pi \quad (k \in \mathbb{Z}). \]
Comme $r$ est un réel positif, $r = \rho^{1/n}$ (la racine $n$-ième réelle usuelle de $\rho$). Et $\theta = \dfrac{\alpha + 2k\pi}{n}$. En donnant à $k$ les valeurs $0, 1, \dots, n-1$, on obtient $n$ arguments distincts modulo $2\pi$, donc $n$ solutions distinctes. Au-delà ($k = n, n+1, \dots$), on retombe sur les mêmes valeurs.

\begin{propriete}[Théorème : existence et expression des $n$ racines n-ièmes]
Soit $a = \rho\, e^{i\alpha}$ un nombre complexe non nul ($\rho > 0$) et $n \geq 2$ un entier. L'équation $z^n = a$ admet \textbf{exactement $n$ solutions distinctes} dans $\mathbb{C}$, appelées racines $n$-ièmes de $a$, données par :
\[ \boxed{\ z_k = \rho^{1/n}\, e^{\,i\frac{\alpha + 2k\pi}{n}}, \qquad k = 0, 1, \dots, n-1.\ } \]
Toutes les racines ont le \textbf{même module} $\rho^{1/n}$ ; leurs arguments forment une progression arithmétique de raison $\dfrac{2\pi}{n}$.
\end{propriete}

\begin{experience}[Démonstration guidée du théorème]
Démontrons rigoureusement le théorème (démonstration formatrice, exigible au Baccalauréat technique).

\textbf{Étape 1 : les $z_k$ sont bien solutions.} Élevons $z_k$ à la puissance $n$ en utilisant la formule de Moivre :
\begin{align*}
z_k^{\,n} &= \left(\rho^{1/n}\right)^n e^{\,i n \cdot \frac{\alpha + 2k\pi}{n}} = \rho\, e^{\,i(\alpha + 2k\pi)} = \rho\, e^{i\alpha} \cdot e^{2ik\pi} = \rho\, e^{i\alpha} = a,
\end{align*}
car $e^{2ik\pi} = 1$ pour tout entier $k$. Chaque $z_k$ est donc bien une racine $n$-ième de $a$.

\textbf{Étape 2 : les $z_k$ sont deux à deux distincts.} Pour $k$ et $k'$ dans $\{0, 1, \dots, n-1\}$ avec $k \neq k'$, les arguments $\dfrac{\alpha + 2k\pi}{n}$ et $\dfrac{\alpha + 2k'\pi}{n}$ diffèrent de $\dfrac{2(k-k')\pi}{n}$, qui n'est \emph{pas} un multiple de $2\pi$ (car $0 < |k - k'| < n$). Deux complexes de même module mais d'arguments non congrus modulo $2\pi$ sont distincts.

\textbf{Étape 3 : il n'y a pas d'autre solution.} Si $z = r e^{i\theta}$ vérifie $z^n = a$, alors $r^n = \rho$ impose $r = \rho^{1/n}$ (unicité de la racine $n$-ième réelle positive), et $n\theta \equiv \alpha \ [2\pi]$, donc $\theta = \dfrac{\alpha + 2k\pi}{n}$ pour un certain entier $k$. En réduisant $k$ modulo $n$, on retombe sur l'un des $z_k$ précédents. L'équation $z^n = a$ est une équation polynomiale de degré $n$ : elle a au plus $n$ racines. Nous en avons trouvé $n$ distinctes : ce sont donc \textbf{toutes} les solutions. $\blacksquare$
\end{experience}

\begin{attention}[Les deux pièges classiques]
\begin{itemize}
\item \textbf{Ajouter $2k\pi$ AVANT de diviser par $n$.} L'argument des racines est $\dfrac{\alpha + 2k\pi}{n}$, et non $\dfrac{\alpha}{n} + 2k\pi$. Si l'on divise d'abord puis qu'on ajoute $2k\pi$, on ne trouve qu'\textbf{une seule} solution au lieu de $n$.
\item \textbf{Ne pas oublier la dernière valeur $k = n-1$.} L'entier $k$ parcourt $\{0, 1, \dots, n-1\}$, soit $n$ valeurs. Oublier $k = n-1$ (ou s'arrêter à $k = n-2$) fait perdre une solution ; prendre $k = n$ redonne la solution $z_0$ (car $\frac{2n\pi}{n} = 2\pi$), ce qui est inutile mais pas faux.
\end{itemize}
\end{attention}

\begin{methode}[Calculer les racines n-ièmes de $a \neq 0$]
\begin{enumerate}
\item \textbf{Écrire $a$ sous forme exponentielle} : calculer $\rho = |a|$, puis un argument $\alpha$ de $a$ (placer $a$ dans le plan, utiliser $\cos\alpha = \frac{\mathrm{Re}(a)}{\rho}$ et $\sin\alpha = \frac{\mathrm{Im}(a)}{\rho}$).
\item \textbf{Prendre la racine n-ième du module} : le module commun des solutions est $r = \rho^{1/n} = \sqrt[n]{\rho}$.
\item \textbf{Diviser les arguments} : les arguments des solutions sont $\theta_k = \dfrac{\alpha + 2k\pi}{n}$ pour $k = 0, 1, \dots, n-1$.
\item \textbf{Écrire les $n$ solutions} $z_k = r\,e^{i\theta_k}$, puis, si demandé, les passer en forme algébrique $x + iy$.
\item \textbf{Vérifier} (au moins pour une solution) en élevant à la puissance $n$ avec la formule de Moivre.
\end{enumerate}
\end{methode}

\courssub{Interprétation géométrique : un polygone régulier}

\begin{propriete}[Configuration géométrique des racines n-ièmes]
Les images $M_0, M_1, \dots, M_{n-1}$ des $n$ racines $n$-ièmes de $a = \rho e^{i\alpha}$ sont les \cle{sommets d'un polygone régulier} à $n$ côtés, inscrit dans le cercle de centre $O$ et de rayon $\rho^{1/n}$. Le premier sommet $M_0$ a pour argument $\dfrac{\alpha}{n}$, et l'on passe d'un sommet au suivant par la rotation de centre $O$ et d'angle $\dfrac{2\pi}{n}$.
\end{propriete}

En effet, tous les $z_k$ ont le même module $\rho^{1/n}$ : leurs images sont donc sur le cercle de centre $O$ et de rayon $\rho^{1/n}$. De plus, $\arg(z_{k+1}) - \arg(z_k) = \dfrac{2\pi}{n}$ : les points sont régulièrement espacés sur ce cercle. C'est cette régularité qui rend la figure si belle et si utile : connaissant \textbf{une seule} racine $n$-ième de $a$, on obtient toutes les autres en la faisant tourner de $\frac{2\pi}{n}$ autour de l'origine.

\begin{popfigure}
\begin{center}
\begin{tikzpicture}[scale=1.5]
  % cercle de rayon 2
  \draw[popline, thick] (0,0) circle (2);
  % axes
  \draw[->, popmuted] (-2.5,0) -- (2.7,0) node[below right] {\footnotesize Réel};
  \draw[->, popmuted] (0,-2.5) -- (0,2.7) node[above left] {\footnotesize Imaginaire};
  % polygone régulier (5 racines de -32), angles 36,108,180,252,324
  \draw[popblue, thick] (36:2) -- (108:2) -- (180:2) -- (252:2) -- (324:2) -- cycle;
  % rayons
  \foreach \a in {36,108,180,252,324} {\draw[poporange, thin] (0,0) -- (\a:2);};
  % points et labels
  \fill[popdark] (36:2) circle (1.6pt) node[above right, popdark] {\footnotesize $z_0$};
  \fill[popdark] (108:2) circle (1.6pt) node[above left, popdark] {\footnotesize $z_1$};
  \fill[popdark] (180:2) circle (1.6pt) node[left, popdark] {\footnotesize $z_2$};
  \fill[popdark] (252:2) circle (1.6pt) node[below left, popdark] {\footnotesize $z_3$};
  \fill[popdark] (324:2) circle (1.6pt) node[below right, popdark] {\footnotesize $z_4$};
  \fill (0,0) circle (1.4pt) node[below left] {\footnotesize $O$};
  % angle 2pi/5
  \draw[popgreen, thick, ->] (36:0.9) arc (36:108:0.9);
  \node[popgreen] at (72:1.25) {\footnotesize $\frac{2\pi}{5}$};
\end{tikzpicture}
\end{center}
\legende{Les cinq racines 5-ièmes de $-32$ : sommets d'un pentagone régulier inscrit dans le cercle de centre $O$ et de rayon $\sqrt[5]{32} = 2$.}
\end{popfigure}

\courssub{Exemple chiffré détaillé : les racines 5-ièmes de $-32$}

\begin{exoresolu}[Résoudre $z^5 = -32$]
\textbf{Énoncé.} Déterminer toutes les solutions dans $\mathbb{C}$ de l'équation $z^5 = -32$, les écrire sous forme exponentielle, puis placer leurs images dans le plan complexe.
\tcblower
\textbf{Solution détaillée.}

\textbf{Étape 1 : forme exponentielle de $a = -32$.} On a $\rho = |-32| = 32$. Le nombre $-32$ est un réel négatif, donc $\alpha = \arg(-32) = \pi \ [2\pi]$. Ainsi :
\[ -32 = 32\, e^{i\pi}. \]

\textbf{Étape 2 : racine 5-ième du module.} $r = 32^{1/5} = \sqrt[5]{32} = 2$, car $2^5 = 32$.

\textbf{Étape 3 : arguments des solutions.} Pour $k = 0, 1, 2, 3, 4$ :
\[ \theta_k = \frac{\pi + 2k\pi}{5}. \]

\textbf{Étape 4 : les cinq solutions.}
\begin{align*}
z_0 &= 2\,e^{i\pi/5}, & z_1 &= 2\,e^{i3\pi/5}, & z_2 &= 2\,e^{i\pi} = -2, \\
z_3 &= 2\,e^{i7\pi/5} = 2\,e^{-i3\pi/5}, & z_4 &= 2\,e^{i9\pi/5} = 2\,e^{-i\pi/5}.
\end{align*}
On remarque que $z_2 = -2$ est l'unique solution \textbf{réelle} (ce qui est cohérent : $(-2)^5 = -32$). Les quatre autres solutions sont non réelles, conjuguées deux à deux : $z_1 = \overline{z_3}$ et $z_0 = \overline{z_4}$.

\textbf{Étape 5 : vérification.} Par la formule de Moivre :
\[ z_1^{\,5} = 2^5\, e^{i \cdot 5 \cdot \frac{3\pi}{5}} = 32\, e^{i3\pi} = 32 \times (-1) = -32. \checkmark \]

\textbf{Étape 6 : représentation.} Les cinq images sont sur le cercle de centre $O$ et de rayon $2$, aux angles $\frac{\pi}{5} = 36^\circ$, $\frac{3\pi}{5} = 108^\circ$, $\pi = 180^\circ$, $\frac{7\pi}{5} = 252^\circ$, $\frac{9\pi}{5} = 324^\circ$ : ce sont les sommets d'un pentagone régulier (voir la figure ci-dessus).
\end{exoresolu}

\courssub{Cas particulier fondamental : les racines n-ièmes de l'unité}

\begin{definition}[Racines n-ièmes de l'unité]
On appelle \cle{racines n-ièmes de l'unité} les solutions de l'équation $z^n = 1$. Ici $\rho = 1$ et $\alpha = 0$, donc ce sont les $n$ nombres :
\[ \omega_k = e^{\,i\frac{2k\pi}{n}}, \qquad k = 0, 1, \dots, n-1. \]
On note souvent $\omega = e^{2i\pi/n}$ ; alors les racines sont $1, \omega, \omega^2, \dots, \omega^{n-1}$.
\end{definition}

Géométriquement, les racines $n$-ièmes de l'unité sont les sommets du \textbf{polygone régulier à $n$ côtés inscrit dans le cercle trigonométrique} (cercle de centre $O$ et de rayon $1$), dont l'un des sommets est le point d'affixe $1$. Par exemple :
\begin{itemize}
\item $n = 2$ : racines $1$ et $-1$ ;
\item $n = 3$ : racines $1$, $j = e^{2i\pi/3} = -\frac{1}{2} + i\frac{\sqrt{3}}{2}$ et $j^2 = \overline{j}$ (triangle équilatéral) ;
\item $n = 4$ : racines $1, i, -1, -i$ (carré).
\end{itemize}

\begin{propriete}[Somme des racines n-ièmes de l'unité]
Pour tout entier $n \geq 2$, la somme des $n$ racines $n$-ièmes de l'unité est nulle :
\[ 1 + \omega + \omega^2 + \dots + \omega^{n-1} = 0, \qquad \text{où } \omega = e^{2i\pi/n}. \]
\textbf{Démonstration.} C'est la somme des $n$ premiers termes d'une suite géométrique de raison $\omega \neq 1$ (car $n \geq 2$) :
\[ 1 + \omega + \dots + \omega^{n-1} = \frac{1 - \omega^n}{1 - \omega} = \frac{1 - 1}{1 - \omega} = 0, \]
puisque $\omega^n = \left(e^{2i\pi/n}\right)^n = e^{2i\pi} = 1$. $\blacksquare$

Géométriquement, cela signifie que l'\emph{isobarycentre} des sommets d'un polygone régulier centré en $O$ est le point $O$ lui-même : les sommets se « compensent » parfaitement.
\end{propriete}

\begin{exemplebox}[Racines cubiques de l'unité et le nombre $j$]
Les racines cubiques de l'unité sont $1$, $j = e^{2i\pi/3}$ et $j^2 = e^{4i\pi/3}$. La propriété précédente donne la célèbre identité :
\[ 1 + j + j^2 = 0. \]
Cette relation est très utilisée en \textbf{électrotechnique} : dans un système triphasé équilibré, les trois tensions déphasées de $\frac{2\pi}{3}$ (soit $120^\circ$) ont une somme nulle à tout instant, ce qui explique l'absence de courant dans le neutre.
\end{exemplebox}

\begin{exercice}[À vous de jouer]
\begin{enumerate}
\item Résoudre dans $\mathbb{C}$ l'équation $z^3 = -8$. Donner les solutions sous forme exponentielle puis algébrique.
\item Résoudre $z^4 = 16$ et représenter les solutions dans le plan complexe. Quelle figure obtient-on ?
\item Montrer que si $z_0$ est une racine $n$-ième de $a$, alors toutes les racines $n$-ièmes de $a$ s'écrivent $z_0 \, \omega^k$ avec $\omega = e^{2i\pi/n}$ et $k = 0, \dots, n-1$.
\end{enumerate}
\end{exercice}

\begin{corrige}[Exercice]
\textbf{1.} $-8 = 8\,e^{i\pi}$, donc $z_k = 2\,e^{i(\pi + 2k\pi)/3}$ pour $k = 0, 1, 2$ : $z_0 = 2e^{i\pi/3} = 1 + i\sqrt{3}$, $z_1 = 2e^{i\pi} = -2$, $z_2 = 2e^{i5\pi/3} = 1 - i\sqrt{3}$.

\textbf{2.} $16 = 16\,e^{i0}$, donc $z_k = 2\,e^{ik\pi/2}$ pour $k = 0, 1, 2, 3$ : $z_0 = 2$, $z_1 = 2i$, $z_2 = -2$, $z_3 = -2i$. Les quatre images forment un \textbf{carré} inscrit dans le cercle de rayon $2$.

\textbf{3.} Si $z_0^n = a$, alors $(z_0 \omega^k)^n = z_0^n (\omega^k)^n = a \times (\omega^n)^k = a \times 1 = a$. Les $n$ nombres $z_0\omega^k$ sont distincts (multiplier par $z_0 \neq 0$ est injectif) et solutions : ce sont donc les $n$ racines de $a$.
\end{corrige}

\begin{attention}[Rédaction et conditions d'application]
\begin{itemize}
\item La formule $z_k = \rho^{1/n} e^{i(\alpha + 2k\pi)/n}$ exige $a \neq 0$ : le nombre $0$ n'a qu'\textbf{une} racine $n$-ième, qui est $0$.
\item $\rho^{1/n}$ désigne la racine $n$-ième \textbf{réelle positive} de $\rho$ : c'est un réel, pas un complexe à chercher.
\item Pensez à \textbf{conclure géométriquement} quand l'énoncé le demande : « les images sont les sommets d'un polygone régulier à $n$ côtés inscrit dans le cercle de centre $O$ et de rayon $\sqrt[n]{\rho}$ ».
\end{itemize}
\end{attention}

\begin{savaistu}[Gauss et les polygones constructibles]
Les racines de l'unité sont au cœur d'un problème célèbre vieux de plus de deux mille ans : quels polygones réguliers peut-on construire \textbf{à la règle et au compas} ? Les Grecs savaient construire le triangle équilatéral, le carré et le pentagone régulier. En 1796, à seulement 19 ans, Carl Friedrich \textbf{Gauss} prouva que le polygone régulier à $17$ côtés est constructible, en montrant que les racines 17-ièmes de l'unité s'expriment à l'aide de racines carrées. Ce résultat, dont Gauss était si fier qu'il souhaita qu'un polygone à 17 côtés figure sur sa pierre tombale (le sculpteur refusa, jugeant la figure indiscernable d'un cercle !), marque la naissance des mathématiques modernes. Aujourd'hui encore, les racines de l'unité sont omniprésentes : traitement du signal, transformée de Fourier rapide (utilisée dans le téléphone mobile et la compression d'images), et systèmes triphasés des installations électriques.
\end{savaistu}

\coursec{Applications techniques et rédaction rigoureuse}

Dans la section précédente, nous avons appris à extraire les racines $n$-ièmes d'un nombre complexe non nul grâce à la forme exponentielle. Nous allons maintenant voir \emph{pourquoi} cette forme exponentielle est si puissante : elle permet de modéliser et de résoudre des problèmes concrets d'électricité, omniprésents dans les métiers techniques (électrotechnique, maintenance, automatique). C'est aussi l'occasion d'apprendre à \cle{rédiger} une démonstration complète et rigoureuse, compétence exigée au Probatoire et au Baccalauréat technique.

\courssub{Modélisation d'un signal sinusoïdal par un nombre complexe}

\textbf{Intuition.} En électricité, les tensions et les courants alternatifs sont des signaux sinusoïdaux : ils oscillent dans le temps. Or, manipuler des cosinus et des sinus est fastidieux. L'idée géniale des ingénieurs : remplacer le signal réel par un nombre complexe tournant, dont le signal réel est simplement la \cle{partie réelle}.

\begin{definition}[Signal sinusoïdal et représentation complexe]
Un signal sinusoïdal de tension s'écrit :
\[ v(t) = V_0 \cos(\omega t + \varphi) \]
où $V_0$ est l'\cle{amplitude} (en volts), $\omega$ la \cle{pulsation} (en rad/s, avec $\omega = 2\pi f$, $f$ étant la fréquence en hertz) et $\varphi$ la \cle{phase à l'origine} (en radians).

On lui associe le nombre complexe dépendant du temps :
\[ \underline{V}(t) = V_0\, e^{i(\omega t + \varphi)} \quad \text{tel que} \quad v(t) = \operatorname{Re}\big(\underline{V}(t)\big). \]
Le nombre complexe fixe $\underline{V} = V_0 e^{i\varphi}$ est appelé \cle{phaseur} (ou amplitude complexe) du signal.
\end{definition}

Pourquoi cela fonctionne-t-il ? Par la formule d'Euler : $e^{i\theta} = \cos\theta + i\sin\theta$, donc $\operatorname{Re}\big(V_0 e^{i(\omega t+\varphi)}\big) = V_0\cos(\omega t + \varphi)$. Le point d'affixe $\underline{V}(t)$ tourne sur le cercle de centre $O$ et de rayon $V_0$ à la vitesse angulaire $\omega$, et sa projection sur l'axe des réels décrit exactement le signal $v(t)$.

\begin{attention}[Grandeurs instantanées et grandeurs complexes]
Ne jamais confondre :
\begin{itemize}
\item la grandeur \textbf{instantanée} $v(t) = V_0\cos(\omega t + \varphi)$, qui est un \textbf{nombre réel} dépendant du temps ;
\item le \textbf{phaseur} $\underline{V} = V_0 e^{i\varphi}$, qui est un \textbf{nombre complexe} constant (le temps a disparu !).
\end{itemize}
On passe de l'une à l'autre par $v(t) = \operatorname{Re}\big(\underline{V}\, e^{i\omega t}\big)$. Écrire $v(t) = \underline{V}$ est une erreur grave : un réel ne peut pas être égal à un complexe non réel.
\end{attention}

\courssub{Impédance complexe et loi d'Ohm complexe}

\textbf{Intuition.} Dans un circuit en courant continu, la résistance $R$ s'oppose au passage du courant : $U = RI$. En courant alternatif, une bobine (inductance $L$) et un condensateur (capacité $C$) s'opposent aussi au courant, mais avec un \emph{déphasage}. Les nombres complexes capturent à la fois l'opposition (module) et le déphasage (argument).

\begin{definition}[Impédance complexe]
L'\cle{impédance complexe} d'un dipôle s'écrit :
\[ \underline{Z} = R + iX \]
où $R \geq 0$ est la \cle{résistance} (partie réelle) et $X$ la \cle{réactance} (partie imaginaire), toutes deux en ohms ($\Omega$).

Pour les trois dipôles élémentaires, à la pulsation $\omega$ :
\begin{itemize}
\item résistor : $\underline{Z}_R = R$ ;
\item bobine : $\underline{Z}_L = iL\omega$ ;
\item condensateur : $\underline{Z}_C = \dfrac{1}{iC\omega} = -\dfrac{i}{C\omega}$.
\end{itemize}
\end{definition}

\begin{propriete}[Loi d'Ohm complexe]
Pour un dipôle parcouru par un courant de phaseur $\underline{I}$ et soumis à une tension de phaseur $\underline{U}$ :
\[ \underline{U} = \underline{Z}\,\underline{I}. \]
En série, les impédances s'additionnent : $\underline{Z}_{\text{total}} = \underline{Z}_1 + \underline{Z}_2 + \cdots$
\end{propriete}

Le module $|\underline{Z}| = \sqrt{R^2 + X^2}$ donne le rapport des amplitudes $U_0/I_0$, et l'argument de $\underline{Z}$ donne le déphasage de la tension par rapport au courant.

\courssub{Puissances active et réactive}

\begin{definition}[Puissances en régime sinusoïdal]
À partir des phaseurs $\underline{U}$ et $\underline{I}$ (valeurs efficaces), on forme la \cle{puissance complexe} :
\[ \underline{S} = \underline{U}\,\overline{\underline{I}} \]
où $\overline{\underline{I}}$ est le conjugué de $\underline{I}$. Alors :
\begin{itemize}
\item la \cle{puissance active} $P = \operatorname{Re}\big(\underline{U}\,\overline{\underline{I}}\big)$, en watts (W) : c'est la puissance réellement consommée (chaleur, travail mécanique) ;
\item la \cle{puissance réactive} $Q = \operatorname{Im}\big(\underline{U}\,\overline{\underline{I}}\big)$, en vars : puissance échangée par les bobines et condensateurs, sans être consommée.
\end{itemize}
\end{definition}

\begin{attention}[Le conjugué est obligatoire]
C'est bien $\underline{U}\,\overline{\underline{I}}$ et non $\underline{U}\,\underline{I}$. Le conjugué garantit que les phases se \emph{soustraient} : si $\underline{U} = U e^{i\varphi_u}$ et $\underline{I} = I e^{i\varphi_i}$, alors
\[ \underline{U}\,\overline{\underline{I}} = U e^{i\varphi_u} \times I e^{-i\varphi_i} = UI\, e^{i(\varphi_u - \varphi_i)}, \]
et l'on retrouve les formules classiques $P = UI\cos(\varphi_u - \varphi_i)$ et $Q = UI\sin(\varphi_u - \varphi_i)$.
\end{attention}

\courssub{Circuit RLC série : courant et résonance}

Considérons un circuit série comportant un résistor $R$, une bobine $L$ et un condensateur $C$, alimenté par un générateur de tension sinusoïdale.

\begin{popfigure}
\begin{center}
\begin{tikzpicture}[scale=1.0, >=Stealth]
% fils et composants
\draw[thick, popdark] (0,0) -- (0,2.25);
\draw[thick, popdark] (0,2.25) -- (1,2.25);
% résistor R
\draw[thick, popblue] (1,2) rectangle (2.6,2.5);
\draw[thick, popdark] (2.6,2.25) -- (3.2,2.25);
\node[popblue, above] at (1.8,2.5) {$R$};
% bobine L
\draw[thick, poporange] (3.2,2.25) arc[start angle=180, end angle=0, radius=0.3];
\draw[thick, poporange] (3.8,2.25) arc[start angle=180, end angle=0, radius=0.3];
\draw[thick, poporange] (4.4,2.25) arc[start angle=180, end angle=0, radius=0.3];
\draw[thick, popdark] (5.0,2.25) -- (5.6,2.25);
\node[poporange, above] at (4.1,2.85) {$L$};
% condensateur C
\draw[thick, popgreen] (5.6,2.25) -- (6.1,2.25);
\draw[thick, popgreen] (6.1,1.95) -- (6.1,2.55);
\draw[thick, popgreen] (6.35,1.95) -- (6.35,2.55);
\draw[thick, popdark] (6.35,2.25) -- (7,2.25);
\node[popgreen, above] at (6.22,2.6) {$C$};
\draw[thick, popdark] (7,2.25) -- (7,0);
\draw[thick, popdark] (7,0) -- (1.2,0);
% générateur
\draw[thick, poppurple] (0.6,0) circle (0.6);
\node[poppurple] at (0.6,0) {$\sim$};
\node[poppurple, left] at (-0.1,0) {$\underline{U}$};
% courant
\draw[->, thick, poppink] (1.2,1.55) -- (2.2,1.55);
\node[poppink, below] at (1.7,1.5) {$\underline{I}$};
\end{tikzpicture}
\end{center}
\legende{Circuit RLC série alimenté par un générateur de tension sinusoïdale.}
\end{popfigure}

L'impédance totale du circuit série est :
\[ \underline{Z} = R + iL\omega - \frac{i}{C\omega} = R + i\left(L\omega - \frac{1}{C\omega}\right). \]

\begin{propriete}[Résonance du circuit RLC série]
Le module du courant $|\underline{I}| = \dfrac{|\underline{U}|}{|\underline{Z}|}$ est maximal lorsque la réactance totale s'annule :
\[ L\omega - \frac{1}{C\omega} = 0 \quad \Longleftrightarrow \quad \omega_0 = \frac{1}{\sqrt{LC}}. \]
Cette pulsation $\omega_0$ est appelée \cle{pulsation de résonance}. À la résonance, $\underline{Z} = R$ : le circuit se comporte comme une simple résistance.
\end{propriete}

\begin{experience}[Démonstration guidée de la pulsation de résonance]
\textbf{Hypothèses :} circuit série $R$, $L$, $C$ parcouru par un courant sinusoïdal de pulsation $\omega > 0$.

\textbf{Étape 1.} Les dipôles étant en série, leurs impédances s'additionnent :
\[ \underline{Z} = \underline{Z}_R + \underline{Z}_L + \underline{Z}_C = R + iL\omega + \frac{1}{iC\omega}. \]

\textbf{Étape 2.} Or $\dfrac{1}{i} = -i$ (car $i \times (-i) = -i^2 = 1$). Donc :
\[ \underline{Z} = R + i\left(L\omega - \frac{1}{C\omega}\right). \]

\textbf{Étape 3.} Le module vaut $|\underline{Z}| = \sqrt{R^2 + \left(L\omega - \dfrac{1}{C\omega}\right)^2}$. Comme $|\underline{I}| = \dfrac{|\underline{U}|}{|\underline{Z}|}$ avec $|\underline{U}|$ fixé, le courant est maximal lorsque $|\underline{Z}|$ est minimal, c'est-à-dire lorsque la quantité sous la racine est minimale, donc lorsque le carré est nul :
\[ L\omega - \frac{1}{C\omega} = 0 \quad \Longleftrightarrow \quad LC\omega^2 = 1 \quad \Longleftrightarrow \quad \omega^2 = \frac{1}{LC}. \]

\textbf{Étape 4.} Comme $\omega > 0$, on conclut :
\[ \omega_0 = \frac{1}{\sqrt{LC}}. \qquad \blacksquare \]
\end{experience}

\courssub{Exemple complet : étude d'un circuit concret}

\begin{exoresolu}[Circuit RLC série en atelier]
Un circuit série est constitué d'un résistor $R = \SI{10}{\Omega}$, d'une bobine $L = \SI{0,1}{H}$ et d'un condensateur de capacité $C = 10^{-4}$ F (soit $100\ \mu$F). Il est alimenté par une tension de phaseur $\underline{U} = 230\,e^{i \times 0} = 230$ (valeur efficace, phase nulle) à la fréquence $f = \SI{50}{Hz}$.

Déterminer : 1) l'impédance complexe $\underline{Z}$ ; 2) le courant $\underline{I}$ ; 3) les puissances $P$ et $Q$ ; 4) la fréquence de résonance.

\tcblower
\textbf{Solution détaillée.}

\textbf{Étape 1 : pulsation.} $\omega = 2\pi f = 2\pi \times 50 = 100\pi \approx \num{314,16}$ rad/s.

\textbf{Étape 2 : impédances partielles.}
\begin{align*}
\underline{Z}_L &= iL\omega = i \times 0{,}1 \times 314{,}16 \approx 31{,}42\,i \ (\Omega)\\
\underline{Z}_C &= -\frac{i}{C\omega} = -\frac{i}{10^{-4} \times 314{,}16} \approx -31{,}83\,i \ (\Omega)
\end{align*}

\textbf{Étape 3 : impédance totale.} Les dipôles sont en série, donc :
\[ \underline{Z} = 10 + i(31{,}42 - 31{,}83) = 10 - 0{,}41\,i \ (\Omega). \]
Module : $|\underline{Z}| = \sqrt{10^2 + 0{,}41^2} \approx \num{10,01}\ \Omega$.

\textbf{Étape 4 : courant.} Par la loi d'Ohm complexe, en multipliant par le conjugué du dénominateur :
\[ \underline{I} = \frac{\underline{U}}{\underline{Z}} = \frac{230}{10 - 0{,}41\,i} = \frac{230\,(10 + 0{,}41\,i)}{10^2 + 0{,}41^2} \approx 22{,}96 + 0{,}94\,i \ (\text{A}). \]
L'intensité efficace est $|\underline{I}| \approx \sqrt{22{,}96^2 + 0{,}94^2} \approx \num{23,0}$ A.

\textbf{Étape 5 : puissances.}
\[ \underline{S} = \underline{U}\,\overline{\underline{I}} = 230\,(22{,}96 - 0{,}94\,i) \approx 5281 - 216\,i. \]
Donc $P = \operatorname{Re}(\underline{S}) \approx \SI{5281}{W}$ et $Q = \operatorname{Im}(\underline{S}) \approx \SI{-216}{var}$.

\textbf{Étape 6 : résonance.}
\[ \omega_0 = \frac{1}{\sqrt{LC}} = \frac{1}{\sqrt{0{,}1 \times 10^{-4}}} = \frac{1}{\sqrt{10^{-5}}} \approx \SI{316,2}{rad/s}, \]
soit $f_0 = \dfrac{\omega_0}{2\pi} \approx \SI{50,3}{Hz}$.

\textbf{Interprétation.} La fréquence du réseau (50 Hz) est très proche de la résonance : la réactance totale ($-0{,}41\ \Omega$) est presque nulle, le courant est presque maximal et le circuit est quasiment résistif ($Q$ proche de 0). C'est pourquoi $P$ est élevée.
\end{exoresolu}

\courssub{Méthode générale de modélisation et conseils de rédaction}

\begin{methode}[Modéliser un problème physique par les nombres complexes]
\begin{enumerate}
\item \textbf{Identifier les grandeurs} : repérer dans l'énoncé amplitudes, fréquences, phases, composants ($R$, $L$, $C$).
\item \textbf{Choisir la forme adaptée} : forme algébrique $a + ib$ pour additionner (impédances en série) ; forme exponentielle $re^{i\theta}$ pour multiplier ou diviser (loi d'Ohm, rapports).
\item \textbf{Calculer dans $\mathbb{C}$} : appliquer les règles de calcul sur les complexes (conjugué, module, argument).
\item \textbf{Revenir au réel} : extraire partie réelle (signal, puissance active), partie imaginaire (puissance réactive), module (amplitude) ou argument (déphasage).
\item \textbf{Interpréter physiquement} : conclure avec les unités et le sens concret (le courant est en avance, le circuit résonne, etc.).
\end{enumerate}
\end{methode}

\begin{methode}[Rédiger une démonstration complète (exemple : formule de Moivre)]
Une démonstration exigible doit comporter quatre parties :
\begin{enumerate}
\item \textbf{Énoncé} : « Montrons que pour tout réel $\theta$ et tout entier naturel $n$, $(\cos\theta + i\sin\theta)^n = \cos(n\theta) + i\sin(n\theta)$. »
\item \textbf{Hypothèses et outil} : « D'après la formule d'Euler, $\cos\theta + i\sin\theta = e^{i\theta}$. »
\item \textbf{Étapes du calcul} : « Par conséquent, $(\cos\theta + i\sin\theta)^n = (e^{i\theta})^n$. Or $(e^{i\theta})^n = e^{in\theta}$ (propriété de l'exponentielle complexe). Donc $(e^{i\theta})^n = \cos(n\theta) + i\sin(n\theta)$, toujours d'après la formule d'Euler. »
\item \textbf{Conclusion} : « On a donc établi la formule de Moivre : $(\cos\theta + i\sin\theta)^n = \cos(n\theta) + i\sin(n\theta)$. »
\end{enumerate}
\end{methode}

\begin{attention}[Conseils de rédaction pour l'examen]
\begin{itemize}
\item Utiliser des \cle{connecteurs logiques} : \emph{donc}, \emph{or}, \emph{par conséquent}, \emph{ainsi}, \emph{c'est-à-dire}. Chaque ligne doit être justifiée par la précédente.
\item \textbf{Numéroter les équations} importantes : $(1)$, $(2)$... pour pouvoir écrire « d'après $(1)$ ».
\item Annoncer ce que l'on fait : « Calculons le module de $\underline{Z}$ » avant de calculer.
\item Encadrer ou souligner le résultat final, avec son \textbf{unité}.
\item Ne jamais écrire une égalité entre un réel et un complexe non réel.
\end{itemize}
\end{attention}

\begin{savaistu}[De Euler à Laplace : les complexes au cœur de l'automatique]
La forme exponentielle $e^{i\omega t}$ que tu utilises pour les signaux sinusoïdaux possède une généralisation puissante : la \cle{transformée de Laplace}, où la pulsation imaginaire $i\omega$ est remplacée par une variable complexe complète $s = \sigma + i\omega$. La partie réelle $\sigma$ décrit l'amortissement d'un système (comment les oscillations s'atténuent), et la partie imaginaire $\omega$ leur fréquence. Cet outil, étudié après le baccalauréat, est à la base de l'automatique : régulation des machines, pilotage des robots, stabilisation des drones. Les nombres complexes, jugés « imaginaires » à leur invention, pilotent aujourd'hui les usines, y compris les unités de transformation agroalimentaire et les centrales électriques du Cameroun.
\end{savaistu}

\begin{exercice}[Exercice de synthèse]
Un atelier est alimenté par une tension de phaseur $\underline{U} = 220$ V (phase nulle) à $f = \SI{50}{Hz}$. Un moteur est modélisé par une impédance $\underline{Z} = 20 + 15\,i\ (\Omega)$.
\begin{enumerate}
\item Calculer le module et l'argument de $\underline{Z}$ (arrondir l'argument à $0{,}01$ rad).
\item En déduire le courant $\underline{I}$ sous forme exponentielle, puis l'intensité efficace.
\item Calculer les puissances active $P$ et réactive $Q$.
\item Le moteur est-il plutôt inductif ou capacitif ? Justifier par le signe de $Q$.
\end{enumerate}
\end{exercice}

\begin{corrige}[Exercice de synthèse]
\begin{enumerate}
\item $|\underline{Z}| = \sqrt{20^2 + 15^2} = \sqrt{625} = 25\ \Omega$. Comme $\underline{Z}$ a ses parties réelle et imaginaire positives, $\arg(\underline{Z}) = \arctan(15/20) = \arctan(0{,}75) \approx 0{,}64$ rad. Donc $\underline{Z} = 25\,e^{0{,}64\,i}$.
\item $\underline{I} = \dfrac{\underline{U}}{\underline{Z}} = \dfrac{220}{25}\,e^{-0{,}64\,i} = 8{,}8\,e^{-0{,}64\,i}$ A. Intensité efficace : $|\underline{I}| = \SI{8,8}{A}$.
\item $\underline{S} = \underline{U}\,\overline{\underline{I}} = 220 \times 8{,}8\,e^{0{,}64\,i} = 1936\,e^{0{,}64\,i}$. Donc $P = 1936\cos(0{,}64) \approx \SI{1550}{W}$ et $Q = 1936\sin(0{,}64) \approx \SI{1160}{var}$.
\item $Q > 0$ : la réactance $X = 15\ \Omega$ est positive, donc le moteur est \textbf{inductif} (il contient des bobinages), ce qui est cohérent avec la réalité d'un moteur électrique.
\end{enumerate}
\end{corrige}

\begin{aretenir}[L'essentiel de la section]
\begin{itemize}
\item Un signal sinusoïdal $v(t) = V_0\cos(\omega t + \varphi)$ est la partie réelle de $V_0 e^{i(\omega t + \varphi)}$ ; le phaseur $\underline{V} = V_0 e^{i\varphi}$ le représente sans le temps.
\item L'impédance complexe $\underline{Z} = R + iX$ généralise la résistance ; la loi d'Ohm devient $\underline{U} = \underline{Z}\,\underline{I}$.
\item $P = \operatorname{Re}\big(\underline{U}\,\overline{\underline{I}}\big)$ (puissance active) et $Q = \operatorname{Im}\big(\underline{U}\,\overline{\underline{I}}\big)$ (puissance réactive).
\item Le circuit RLC série résonne à la pulsation $\omega_0 = \dfrac{1}{\sqrt{LC}}$.
\item Une bonne démonstration : énoncé, hypothèses, étapes justifiées avec connecteurs logiques, conclusion.
\end{itemize}
\end{aretenir}

\newpage

\coursec{Activité d'intégration}

\coursec{Leçon N11 : Nombres complexes — Approche géométrique}

\courssub{Objectifs de l'activité}
\begin{itemize}
    \item Mobiliser l'ensemble des savoirs de la leçon (représentation géométrique, forme exponentielle, linéarisation, racines $n$-ièmes) dans une situation professionnelle unique.
    \item Modéliser un circuit électrique réel (RLC série) sous régime sinusoïdal forcé à l'aide des nombres complexes.
    \item Calculer des grandeurs physiques (impédance, courant, puissance active) et interpréter les résultats géométriquement.
    \item Appliquer les formules de Moivre (linéarisation) pour analyser les harmoniques d'un signal non linéaire.
    \item Déterminer les racines $n$-ièmes d'un nombre complexe non nul et en donner l'interprétation géométrique (régularité polygonale).
    \item Rédiger un compte-rendu technique structuré, justifiant chaque étape par des propriétés mathématiques ou des lois physiques.
\end{itemize}

\courssub{Situation professionnelle : Maintenance préventive poste HTA/BT \textsc{Socapalm} Kribi}
Dans le cadre de la maintenance préventive du poste de transformation \textbf{HTA/BT} de l'usine \textsc{Socapalm} à \textbf{Kribi}, un technicien supérieur doit analyser la qualité de l'énergie fournie à un atelier de soudure. La tension aux bornes d'un récepteur critique (un groupe de redressement pour poste de soudure) est mesurée à l'aide d'un analyseur de réseau. L'expression instantanée de cette tension est :
\[
u(t) = 230\sqrt{2} \cos(100\pi t + \tfrac{\pi}{6}) \quad \text{en Volts (V)}.
\]
La fréquence du réseau est donc $f = 50\,\text{Hz}$ ($\omega = 100\pi\,\text{rad/s}$). La valeur efficace (RMS) de la tension est $U = 230\,\text{V}$.

Le récepteur est modélisé par un circuit \textbf{RLC série} dont les paramètres nominaux sont :
\begin{itemize}
    \item Résistance $R = 5\,\Omega$ (représentant les pertes cuivre et fer),
    \item Inductance $L = 20\,\text{mH} = 2 \times 10^{-2}\,\text{H}$,
    \item Capacité $C = 50\,\mu\text{F} = 5 \times 10^{-5}\,\text{F}$ (batterie de condensateurs de découplage).
\end{itemize}

L'ingénieur de maintenance demande au technicien de réaliser une étude complète pour :
\begin{enumerate}
    \item Vérifier le dimensionnement du circuit (impédance, courant, puissance active).
    \item Anticiper les perturbations harmoniques dues au redressement (étude du cube de la tension).
    \item Comprendre la symétrie géométrique liée aux racines cubiques de l'impédance au cube (lien avec la stabilité du système).
\end{enumerate}

\courssub{Consignes}
Travailler par binôme. Rédiger les réponses sur une feuille de travail en justifiant chaque calcul. Présenter les résultats finaux avec leurs unités.

\begin{enumerate}
    \item \textbf{Représentation phasorielle.} Écrire le phasor tension $\underline{U}$ (valeur efficace) sous forme exponentielle. Représenter ce phasor dans le plan complexe.
    \item \textbf{Analyse du circuit RLC.}
    \begin{enumerate}
        \item Calculer les réactances $X_L$ et $X_C$. En déduire l'impédance complexe $\underline{Z}$ sous forme algébrique $a + \mathrm{j}b$ et sous forme exponentielle $\rho e^{\mathrm{j}\varphi}$.
        \item Déterminer le phasor courant $\underline{I}$ (valeur efficace) sous forme exponentielle. Donner la valeur efficace $I$ et le déphasage $\varphi_I$ du courant par rapport à la tension.
        \item Calculer la puissance active $P$ absorbée par le récepteur. Préciser la nature du circuit (inductif ou capacitif).
    \end{enumerate}
    \item \textbf{Linéarisation et harmoniques.} Le redressement génère des non-linéarités. On cherche à exprimer $u^3(t)$.
    \begin{enumerate}
        \item Linéariser $\cos^3(\theta)$ en combinaison linéaire de cosinus d'angles multiples (formules de Moivre / linéarisation).
        \item En déduire l'expression de $u^3(t)$. Identifier l'amplitude et la pulsation de la composante de fréquence triple (harmonique d'ordre 3).
    \end{enumerate}
    \item \textbf{Racines cubiques de l'impédance au cube.}
    \begin{enumerate}
        \item Calculer $\underline{Z}^3$ sous forme exponentielle.
        \item Déterminer les trois racines cubiques de $\underline{Z}^3$, c'est-à-dire les solutions de l'équation $z^3 = \underline{Z}^3$. Les exprimer sous forme exponentielle.
        \item Représenter ces trois racines dans le plan complexe. Quelle figure géométrique forment-elles ? Quel est le lien avec $\underline{Z}$ ?
    \end{enumerate}
\end{enumerate}

\courssub{Ressources mathématiques et physiques mobilisables}
\begin{aretenir}[Formulaire de référence]
\begin{itemize}
    \item \textbf{Phasor (valeur efficace) :} $\underline{U} = U e^{\mathrm{j}\varphi_u}$ avec $U = \frac{U_{\text{max}}}{\sqrt{2}}$.
    \item \textbf{Réactances :} $X_L = \omega L$, $X_C = \frac{1}{\omega C}$.
    \item \textbf{Impédance série :} $\underline{Z} = R + \mathrm{j}(X_L - X_C)$. Module $Z = |\underline{Z}| = \sqrt{R^2 + (X_L - X_C)^2}$. Argument $\varphi_Z = \arg(\underline{Z}) = \arctan\left(\frac{X_L - X_C}{R}\right)$ (choisi dans $]-\pi,\pi]$).
    \item \textbf{Loi d'Ohm complexe :} $\underline{I} = \frac{\underline{U}}{\underline{Z}}$. Module $I = \frac{U}{Z}$. Argument $\varphi_I = \varphi_U - \varphi_Z$.
    \item \textbf{Puissance active :} $P = U I \cos(\varphi_U - \varphi_I) = R I^2 = \frac{U^2 R}{Z^2}$.
    \item \textbf{Formules d'Euler / Moivre :} $\cos\theta = \frac{e^{\mathrm{j}\theta}+e^{-\mathrm{j}\theta}}{2}$, $(\cos\theta + \mathrm{j}\sin\theta)^n = \cos(n\theta) + \mathrm{j}\sin(n\theta)$.
    \item \textbf{Linéarisation de $\cos^3\theta$ :} $\cos^3\theta = \frac{1}{4}(3\cos\theta + \cos 3\theta)$.
    \item \textbf{Racines $n$-ièmes :} Les solutions de $z^n = \rho e^{\mathrm{j}\varphi}$ ($\rho>0$) sont $z_k = \sqrt[n]{\rho} \, e^{\mathrm{j}\left(\frac{\varphi + 2k\pi}{n}\right)}$, $k = 0, 1, \dots, n-1$. Elles sont les sommets d'un polygone régulier à $n$ côtés centré à l'origine.
\end{itemize}
\end{aretenir}

\courssub{Démarche de résolution et corrigé commenté}

\begin{exoresolu}{Question 1 : Représentation phasorielle de la tension}
La tension instantanée est $u(t) = 230\sqrt{2} \cos(100\pi t + \frac{\pi}{6})$.
La valeur efficace (RMS) est $U = 230\,\text{V}$. Le déphasage initial est $\varphi_U = \frac{\pi}{6}$.
Le phasor tension (convention valeur efficace) s'écrit sous forme exponentielle :
\[
\boxed{\underline{U} = 230 \, e^{\mathrm{j}\frac{\pi}{6}} \ \text{V}}.
\]
\emph{Représentation géométrique :} Dans le plan complexe (référentiel $(\vec{1}, \vec{\mathrm{j}})$), le point $U$ d'affixe $\underline{U}$ a pour coordonnées polaires $(230; \frac{\pi}{6})$. Il se situe dans le premier quadrant, sur le cercle de rayon 230, formant un angle de $30^\circ$ avec l'axe des réels positifs.

\begin{popfigure}
\begin{tikzpicture}[scale=0.012, >=Stealth]
    % Axes
    \draw[->, popmuted] (-260,0) -- (260,0) node[right] {$\text{Re}$};
    \draw[->, popmuted] (0,-260) -- (0,260) node[above] {$\text{Im}$};
    % Cercle
    \draw[popblue!50, dashed] (0,0) circle (230);
    % Phasor U
    \draw[popcurve, thick, ->] (0,0) -- (199.2, 115) node[midway, above right] {$230$};
    % Angle
    \draw[poporange, thick] (30,0) arc (0:30:30) node[midway, right] {$\frac{\pi}{6}$};
    % Point U
    \fill[popblue] (199.2, 115) circle (6) node[above right] {$U(\underline{U})$};
    % Coordinates
    \draw[dashed, popmuted] (199.2, 115) -- (199.2, 0) node[below] {$199$};
    \draw[dashed, popmuted] (199.2, 115) -- (0, 115) node[left] {$115$};
\end{tikzpicture}
\legende{Représentation du phasor tension $\underline{U} = 230 e^{\mathrm{j}\pi/6}$ V dans le plan complexe.}
\end{popfigure}
\end{exoresolu}

\begin{exoresolu}{Question 2 : Analyse du circuit RLC série}
\emph{Données :} $\omega = 100\pi \approx 314,16\,\text{rad/s}$, $R = 5\,\Omega$, $L = 0,02\,\text{H}$, $C = 5 \times 10^{-5}\,\text{F}$.

\begin{enumerate}
    \item \textbf{Réactances et Impédance.}
    \begin{align*}
        X_L &= \omega L = 100\pi \times 0,02 = 2\pi \approx \mathbf{6,283\,\Omega}, \\
        X_C &= \frac{1}{\omega C} = \frac{1}{100\pi \times 5 \times 10^{-5}} = \frac{1}{5\pi \times 10^{-3}} = \frac{200}{\pi} \approx \mathbf{63,662\,\Omega}.
    \end{align*}
    Comme $X_C > X_L$, le circuit est à prédominance \textbf{capacitive}.
    L'impédance complexe :
    \[
    \underline{Z} = R + \mathrm{j}(X_L - X_C) = 5 + \mathrm{j}\left(2\pi - \tfrac{200}{\pi}\right) = 5 - \mathrm{j}\left(\tfrac{200}{\pi} - 2\pi\right).
    \]
    Valeur numérique de la partie imaginaire : $X_L - X_C \approx 6,283 - 63,662 = -57,379\,\Omega$.
    Forme algébrique : $\boxed{\underline{Z} \approx 5 - \mathrm{j}\,57,38 \ \Omega}$.
    
    Module de l'impédance :
    \[
    Z = |\underline{Z}| = \sqrt{R^2 + (X_L - X_C)^2} = \sqrt{5^2 + \left(\tfrac{200}{\pi} - 2\pi\right)^2} \approx \sqrt{25 + 3292,3} \approx \sqrt{3317,3} \approx \mathbf{57,60\,\Omega}.
    \]
    Argument de l'impédance (déphasage tension/courant aux bornes de $Z$) :
    \[
    \varphi_Z = \arg(\underline{Z}) = \arctan\left(\frac{X_L - X_C}{R}\right) = \arctan\left(\frac{-57,379}{5}\right) = \arctan(-11,476) \approx \mathbf{-1,484\,\text{rad}} \ (\approx -85,0^\circ).
    \]
    Forme exponentielle :
    \[
    \boxed{\underline{Z} \approx 57,60 \, e^{-\mathrm{j}\,1,484} \ \Omega}.
    \]

    \item \textbf{Phasor courant et valeurs efficaces.}
    Loi d'Ohm complexe : $\underline{I} = \frac{\underline{U}}{\underline{Z}} = \frac{230 e^{\mathrm{j}\pi/6}}{57,60 e^{-\mathrm{j}\,1,484}}$.
    Module du courant :
    \[
    I = \frac{U}{Z} = \frac{230}{57,60} \approx \mathbf{3,993\,\text{A}} \approx 4,0\,\text{A}.
    \]
    Argument du courant (phase du courant par rapport à l'origine des temps) :
    \[
    \varphi_I = \varphi_U - \varphi_Z = \frac{\pi}{6} - (-1,484) = 0,524 + 1,484 = \mathbf{2,008\,\text{rad}} \ (\approx 115,0^\circ).
    \]
    Phasor courant sous forme exponentielle :
    \[
    \boxed{\underline{I} \approx 3,99 \, e^{\mathrm{j}\,2,01} \ \text{A}}.
    \]
    \emph{Interprétation :} Le courant est en avance de phase sur la tension de $\varphi_I - \varphi_U = -\varphi_Z \approx 1,484\,\text{rad}$ ($85^\circ$), ce qui confirme le comportement \textbf{capacitif} du circuit ($X_C \gg X_L$).

    \item \textbf{Puissance active.}
    \[
    P = R I^2 = 5 \times (3,993)^2 \approx 5 \times 15,94 \approx \mathbf{79,7\,\text{W}}.
    \]
    Vérification : $P = U I \cos(\varphi_U - \varphi_I) = 230 \times 3,993 \times \cos(-1,484) \approx 918,4 \times 0,087 \approx 79,9\,\text{W}$ (écarts d'arrondi).
    La puissance active absorbée est d'environ \textbf{80 W}. Le facteur de puissance est $\cos(\varphi_U - \varphi_I) \approx 0,087$ (très faible, circuit très réactif).
\end{enumerate}
\end{exoresolu}

\begin{exoresolu}{Question 3 : Linéarisation et harmoniques d'ordre 3}
On cherche à exprimer $u^3(t) = \left(230\sqrt{2} \cos(100\pi t + \tfrac{\pi}{6})\right)^3 = (230\sqrt{2})^3 \cos^3(\omega t + \tfrac{\pi}{6})$.
Posons $\theta = \omega t + \tfrac{\pi}{6}$.

\begin{enumerate}
    \item \textbf{Linéarisation de $\cos^3\theta$.}
    Méthode via formules d'Euler :
    \[
    \cos\theta = \frac{e^{\mathrm{j}\theta} + e^{-\mathrm{j}\theta}}{2} \implies \cos^3\theta = \frac{1}{8}(e^{\mathrm{j}\theta} + e^{-\mathrm{j}\theta})^3.
    \]
    Développement du binôme (Newton) :
    \[
    (e^{\mathrm{j}\theta} + e^{-\mathrm{j}\theta})^3 = e^{\mathrm{j}3\theta} + 3 e^{\mathrm{j}\theta} + 3 e^{-\mathrm{j}\theta} + e^{-\mathrm{j}3\theta} = (e^{\mathrm{j}3\theta} + e^{-\mathrm{j}3\theta}) + 3(e^{\mathrm{j}\theta} + e^{-\mathrm{j}\theta}).
    \]
    Donc :
    \[
    \cos^3\theta = \frac{1}{8}\left[ 2\cos(3\theta) + 6\cos\theta \right] = \frac{1}{4}\cos(3\theta) + \frac{3}{4}\cos\theta.
    \]
    \textbf{Formule de linéarisation :} $\boxed{\cos^3\theta = \frac{1}{4}(3\cos\theta + \cos 3\theta)}$.

    \item \textbf{Expression de $u^3(t)$ et harmonique d'ordre 3.}
    \[
    u^3(t) = (230\sqrt{2})^3 \left[ \frac{3}{4}\cos(\omega t + \tfrac{\pi}{6}) + \frac{1}{4}\cos(3\omega t + \tfrac{\pi}{2}) \right].
    \]
    Calcul du coefficient : $(230\sqrt{2})^3 = 230^3 \times 2\sqrt{2} = 12\,167\,000 \times 2,828 \approx 34\,410\,000 \approx 3,44 \times 10^7\,\text{V}^3$.
    \[
    u^3(t) \approx 2,58 \times 10^7 \cos(100\pi t + \tfrac{\pi}{6}) + 8,60 \times 10^6 \cos(300\pi t + \tfrac{\pi}{2}).
    \]
    \textbf{Composante de fréquence triple (harmonique 3) :}
    \begin{itemize}
        \item Pulsation : $\omega_3 = 3\omega = 300\pi\,\text{rad/s}$ ($f_3 = 150\,\text{Hz}$).
        \item Amplitude (valeur de crête) : $\frac{1}{4}(230\sqrt{2})^3 \approx \mathbf{8,6 \times 10^6 \ \text{V}^3}$.
        \item Phase : $\frac{\pi}{2}$ (déphasage de $90^\circ$ par rapport au cosinus fondamental).
    \end{itemize}
    \emph{Note technique :} Cette composante à 150 Hz, bien que ne représentant pas une tension physique directe (unité V$^3$), modélise la distorsion de tension générée par la non-linéarité du redresseur. Elle justifie l'installation de filtres harmoniques passifs (LC) accordés sur 150 Hz.
\end{enumerate}
\end{exoresolu}

\begin{exoresolu}{Question 4 : Racines cubiques de $\underline{Z}^3$}
\begin{enumerate}
    \item \textbf{Calcul de $\underline{Z}^3$.}
    On utilise la forme exponentielle de $\underline{Z} \approx 57,60 \, e^{-\mathrm{j}\,1,484}$.
    Par la formule de Moivre : $\underline{Z}^3 = Z^3 e^{\mathrm{j}\,3\varphi_Z}$.
    \[
    Z^3 \approx (57,60)^3 \approx 191\,100 \approx 1,91 \times 10^5 \ \Omega^3.
    \]
    \[
    3\varphi_Z \approx 3 \times (-1,484) = -4,452\,\text{rad}.
    \]
    On ramène l'argument principal dans $]-\pi, \pi]$ : $-4,452 + 2\pi \approx -4,452 + 6,283 = \mathbf{1,831\,\text{rad}}$.
    \[
    \boxed{\underline{Z}^3 \approx 1,91 \times 10^5 \, e^{\mathrm{j}\,1,831} \ \Omega^3}.
    \]

    \item \textbf{Résolution de $z^3 = \underline{Z}^3$.}
    Soit $w = \underline{Z}^3 = \rho e^{\mathrm{j}\varphi}$ avec $\rho = Z^3 \approx 1,91 \times 10^5$ et $\varphi \approx 1,831\,\text{rad}$.
    Les trois racines cubiques sont données par ($k = 0, 1, 2$) :
    \[
    z_k = \sqrt[3]{\rho} \, e^{\mathrm{j}\left(\frac{\varphi + 2k\pi}{3}\right)} = Z \, e^{\mathrm{j}\left(\frac{\varphi + 2k\pi}{3}\right)}.
    \]
    Comme $\sqrt[3]{\rho} = Z \approx 57,60$.
    \begin{itemize}
        \item $k=0$ : $z_0 = 57,60 \, e^{\mathrm{j}\frac{1,831}{3}} = 57,60 \, e^{\mathrm{j}\,0,610}$.
        \item $k=1$ : $z_1 = 57,60 \, e^{\mathrm{j}\frac{1,831 + 2\pi}{3}} = 57,60 \, e^{\mathrm{j}\frac{1,831 + 6,283}{3}} = 57,60 \, e^{\mathrm{j}\,2,705}$.
        \item $k=2$ : $z_2 = 57,60 \, e^{\mathrm{j}\frac{1,831 + 4\pi}{3}} = 57,60 \, e^{\mathrm{j}\frac{1,831 + 12,566}{3}} = 57,60 \, e^{\mathrm{j}\,4,799}$.
              Argument principal : $4,799 - 2\pi \approx -1,484\,\text{rad}$.
    \end{itemize}
    Les trois racines sont :
    \[
    \boxed{z_0 \approx 57,60 e^{\mathrm{j}\,0,610},\quad z_1 \approx 57,60 e^{\mathrm{j}\,2,705},\quad z_2 \approx 57,60 e^{-\mathrm{j}\,1,484}}.
    \]
    \emph{Remarque :} $z_2$ correspond exactement à l'impédance initiale $\underline{Z}$ (argument $-1,484$). $z_0$ et $z_1$ sont les deux autres solutions complexes.

    \item \textbf{Représentation géométrique et interprétation.}
    Dans le plan complexe, les points $Z_0, Z_1, Z_2$ d'affixes $z_0, z_1, z_2$ ont tous le même module $Z \approx 57,60$. Ils sont situés sur le cercle de centre $O$ (origine) et de rayon $Z$.
    Leurs arguments sont $\alpha_0 \approx 0,610$, $\alpha_1 \approx 2,705$, $\alpha_2 \approx -1,484$ (ou $4,799$).
    Les différences d'arguments successives sont :
    \[
    \alpha_1 - \alpha_0 \approx 2,095 \approx \frac{2\pi}{3},\quad \alpha_2 - \alpha_1 \approx -4,189 \equiv \frac{2\pi}{3}\ (\text{mod } 2\pi),\quad \alpha_0 - \alpha_2 \approx 2,094 \approx \frac{2\pi}{3}.
    \]
    Les trois points forment les sommets d'un \textbf{triangle équilatéral} inscrit dans le cercle de rayon $Z$, centré à l'origine.

\begin{popfigure}
\begin{tikzpicture}[scale=0.025, >=Stealth]
    % Axes
    \draw[->, popmuted] (-70,0) -- (70,0) node[right] {$\text{Re}$};
    \draw[->, popmuted] (0,-70) -- (0,70) node[above] {$\text{Im}$};
    % Cercle
    \draw[popblue!50, dashed] (0,0) circle (57.6);
    % Racines
    \coordinate (O) at (0,0);
    \coordinate (Z0) at (34.95:57.6);
    \coordinate (Z1) at (155.0:57.6);
    \coordinate (Z2) at (-85.03:57.6);
    % Triangle
    \draw[popgreen, thick] (Z0) -- (Z1) -- (Z2) -- cycle;
    % Points
    \fill[popblue] (Z0) circle (2.5) node[above right] {$z_0$};
    \fill[popblue] (Z1) circle (2.5) node[above left] {$z_1$};
    \fill[poporange] (Z2) circle (2.5) node[below right] {$z_2 = \underline{Z}$};
    % Rayons
    \draw[popmuted, thin] (O) -- (Z0);
    \draw[popmuted, thin] (O) -- (Z1);
    \draw[popmuted, thin] (O) -- (Z2);
    % Angles
    \draw[poporange, thick] (10,0) arc (0:0.610:10) node[midway, right] {$\alpha_0$};
    \draw[poporange, thick] (10,0) arc (0:2.705:15) node[midway, left] {$\alpha_1$};
    \draw[poporange, thick] (10,0) arc (0:-1.484:10) node[midway, below] {$\alpha_2$};
    % Centre
    \fill[popdark] (O) circle (1.5) node[below left] {$O$};
\end{tikzpicture}
\legende{Racines cubiques de $\underline{Z}^3$ : sommets d'un triangle équilatéral. $z_2$ coïncide avec l'impédance physique $\underline{Z}$.}
\end{popfigure}

    \emph{Interprétation physique :} L'impédance physique $\underline{Z}$ n'est qu'une des trois solutions mathématiques de l'équation $z^3 = \underline{Z}^3$. Les deux autres ($z_0, z_1$) correspondent à des impédances « virtuelles » qui, mises au cube, donnent la même valeur complexe. Géométriquement, elles sont obtenues par rotations successives de $\frac{2\pi}{3}$ ($120^\circ$) autour de l'origine. Cette symétrie de rotation d'ordre 3 est la signature géométrique de l'équation $z^3 = c$.
\end{enumerate}
\end{exoresolu}

\courssub{Bilan et synthèse de la leçon}
\begin{savaistu}[Ce qu'il faut retenir de cette activité d'intégration]
Cette activité a permis de mettre en évidence la puissance unificatrice de l'outil « nombres complexes » en électrotechnique :
\begin{itemize}
    \item \textbf{Modélisation unifiée :} La forme exponentielle a transformé les équations différentielles du circuit RLC en algèbre complexe simple (loi d'Ohm généralisée), donnant directement module et phase du courant.
    \item \textbf{Diagnostic énergétique :} Le calcul de la puissance active $P \approx 80\,\text{W}$ pour un courant de $4\,\text{A}$ sous $230\,\text{V}$ révèle un facteur de puissance très dégradé ($\cos\varphi \approx 0,09$), signature d'un circuit fortement capacitif ($X_C \gg X_L$). Cela impose une correction du facteur de puissance (ajout d'inductances) pour éviter les surintensités inutiles sur le transformateur de l'usine.
    \item \textbf{Analyse spectrale par linéarisation :} L'utilisation des formules de Moivre pour linéariser $\cos^3\theta$ a permis d'isoler analytiquement l'harmonique d'ordre 3 (150 Hz) générée par la non-linéarité du redresseur. C'est la base mathématique du dimensionnement des filtres harmoniques passifs.
    \item \textbf{Géométrie des racines :} La recherche des racines cubiques de $\underline{Z}^3$ a illustré le théorème fondamental : les $n$ racines $n$-ièmes d'un nombre complexe non nul sont les sommets d'un polygone régulier à $n$ côtés centré à l'origine. Ici, un triangle équilatéral. L'impédance physique n'est qu'un des sommets ; les deux autres, bien que non physiques pour ce circuit, existent mathématiquement par symétrie de rotation de $120^\circ$.
    \item \textbf{Compétences transverses :} L'activité a exigé une rédaction rigoureuse (unités, forme exponentielle/algébrique, arrondis cohérents), le passage fluide entre domaine temporel ($u(t)$) et domaine fréquentiel (phasors), et l'interprétation physique des résultats mathématiques (nature capacitive, harmoniques, symétrie).
\end{itemize}
L'élève termine cette leçon capable de traiter un problème technique complet : de la mesure oscilloscopique au dimensionnement de filtre, en passant par l'analyse de la stabilité géométrique du système.
\end{savaistu}

\newpage

\coursec{Méthodes et exercices résolus}

\courssub{Méthode 1 : Placer un point et lire l'affixe dans le plan complexe}

\begin{definition}[Plan complexe, image et affixe]
Le plan est muni d'un repère orthonormé direct $(O\,;\vec{u},\vec{v})$. À tout nombre complexe $z = a + \mathrm{i}b$ (avec $a$ et $b$ réels), on associe :
\begin{itemize}
\item le point $M$ de coordonnées $(a\,;b)$, appelé \cle{image} de $z$ ;
\item le vecteur $\vec{w}$ de coordonnées $(a\,;b)$, appelé \cle{vecteur image} de $z$.
\end{itemize}
Réciproquement, $z$ est appelé l'\cle{affixe} du point $M$ (ou du vecteur $\vec{w}$). L'axe $(O\,;\vec{u})$ est l'\cle{axe des réels}, l'axe $(O\,;\vec{v})$ est l'\cle{axe des imaginaires}.
\end{definition}

\begin{methode}[Représenter géométriquement un nombre complexe]
\begin{enumerate}
\item Identifier la \cle{partie réelle} $a$ et la \cle{partie imaginaire} $b$ de $z = a + \mathrm{i}b$.
\item Dans le repère orthonormé $(O\,;\vec{u},\vec{v})$, placer le point $M$ d'abscisse $a$ et d'ordonnée $b$ : $M$ est l'\cle{image} de $z$, et $z$ est l'\cle{affixe} de $M$.
\item Pour lire l'affixe d'un point donné : projeter le point sur l'axe des réels (on obtient $a$) puis sur l'axe des imaginaires (on obtient $b$), puis écrire $z = a + \mathrm{i}b$.
\item Rappels utiles : l'affixe du vecteur $\overrightarrow{AB}$ est $z_B - z_A$, et le milieu de $[AB]$ a pour affixe $\dfrac{z_A+z_B}{2}$.
\end{enumerate}
\end{methode}

\begin{exoresolu}[Lecture et placement d'affixes]
Le plan est muni d'un repère orthonormé $(O\,;\vec{u},\vec{v})$. On donne les points $A$, $B$, $C$ d'affixes respectives $z_A = 2 + \mathrm{i}$, $z_B = -1 + 2\mathrm{i}$, $z_C = 3 - 2\mathrm{i}$.
\begin{enumerate}
\item Placer les points $A$, $B$, $C$.
\item Déterminer l'affixe du vecteur $\overrightarrow{AB}$ et celle du milieu $I$ de $[AC]$.
\end{enumerate}
\tcblower
\textbf{1.} Le point $A$ a pour coordonnées $(2\,;1)$, $B$ a pour coordonnées $(-1\,;2)$ et $C$ a pour coordonnées $(3\,;-2)$. On place chaque point en reportant la partie réelle sur l'axe horizontal et la partie imaginaire sur l'axe vertical.

\textbf{2.} L'affixe du vecteur $\overrightarrow{AB}$ est :
\[ z_{\overrightarrow{AB}} = z_B - z_A = (-1 + 2\mathrm{i}) - (2 + \mathrm{i}) = -1 - 2 + (2-1)\mathrm{i} = -3 + \mathrm{i}. \]
L'affixe du milieu $I$ de $[AC]$ est :
\[ z_I = \frac{z_A + z_C}{2} = \frac{(2+\mathrm{i}) + (3 - 2\mathrm{i})}{2} = \frac{5 - \mathrm{i}}{2} = \frac{5}{2} - \frac{1}{2}\mathrm{i}. \]
\textbf{Conclusion :} $z_{\overrightarrow{AB}} = -3 + \mathrm{i}$ et $z_I = \dfrac{5}{2} - \dfrac{1}{2}\mathrm{i}$.
\end{exoresolu}

\begin{popfigure}
\begin{center}
\begin{tikzpicture}[scale=1.1]
\draw[popline] (-2.4,0) -- (4.2,0);
\draw[popline] (0,-2.8) -- (0,2.8);
\draw[->,thick,popdark] (0,0) -- (1,0) node[below left] {$\vec{u}$};
\draw[->,thick,popdark] (0,0) -- (0,1) node[below left] {$\vec{v}$};
\node[below left, popdark] at (0,0) {$O$};
\foreach \x in {-2,-1,1,2,3,4} \draw[popline] (\x,0.06) -- (\x,-0.06) node[below, scale=0.8] {\x};
\foreach \y in {-2,-1,1,2} \draw[popline] (0.06,\y) -- (-0.06,\y) node[left, scale=0.8] {\y};
\fill[popblue] (2,1) circle (2.2pt) node[above right, popblue] {$A\ (2+\mathrm{i})$};
\fill[poporange] (-1,2) circle (2.2pt) node[above left, poporange] {$B\ (-1+2\mathrm{i})$};
\fill[popgreen] (3,-2) circle (2.2pt) node[below right, popgreen!60!black] {$C\ (3-2\mathrm{i})$};
\fill[poppurple] (2.5,-0.5) circle (2pt) node[right, poppurple] {$I$};
\draw[->,thick,popink] (2,1) -- (-1,2) node[midway, above, sloped, scale=0.85] {$\overrightarrow{AB}$};
\end{tikzpicture}
\end{center}
\legende{Images des affixes $z_A$, $z_B$, $z_C$ et vecteur $\overrightarrow{AB}$ dans le plan complexe.}
\end{popfigure}

\courssub{Méthode 2 : Passer de la forme algébrique à la forme trigonométrique}

\begin{definition}[Module et argument]
Soit $z = a + \mathrm{i}b$ un nombre complexe et $M$ son image dans le plan complexe.
\begin{itemize}
\item Le \cle{module} de $z$ est le réel positif $|z| = \sqrt{a^2+b^2} = OM$.
\item Si $z \neq 0$, un \cle{argument} de $z$ est une mesure, en radians, de l'angle orienté $(\vec{u}\,;\overrightarrow{OM})$. On note $\arg(z) = \theta\ [2\pi]$ : un complexe non nul possède une infinité d'arguments, qui diffèrent d'un multiple de $2\pi$.
\end{itemize}
\end{definition}

\begin{methode}[Calculer module et argument]
Pour $z = a + \mathrm{i}b$ \textbf{non nul} :
\begin{enumerate}
\item Calculer le \cle{module} : $|z| = \sqrt{a^2 + b^2}$.
\item Déterminer un \cle{argument} $\theta$ grâce au système :
\[ \cos\theta = \frac{a}{|z|} \qquad \text{et} \qquad \sin\theta = \frac{b}{|z|}. \]
\item Identifier $\theta$ à l'aide du cercle trigonométrique (valeurs remarquables : $0$, $\frac{\pi}{6}$, $\frac{\pi}{4}$, $\frac{\pi}{3}$, $\frac{\pi}{2}$, $\pi$\ldots). Les DEUX conditions (cosinus ET sinus) doivent être vérifiées.
\item Écrire la forme trigonométrique : $z = |z|\,(\cos\theta + \mathrm{i}\sin\theta)$.
\end{enumerate}
\end{methode}

\begin{exoresolu}[Forme trigonométrique de $z = 1 + \mathrm{i}\sqrt{3}$]
Soit $z = 1 + \mathrm{i}\sqrt{3}$.
\begin{enumerate}
\item Calculer le module de $z$.
\item Déterminer un argument de $z$, puis écrire $z$ sous forme trigonométrique.
\end{enumerate}
\tcblower
\textbf{1.} On a $a = 1$ et $b = \sqrt{3}$, donc :
\[ |z| = \sqrt{1^2 + (\sqrt{3})^2} = \sqrt{1 + 3} = \sqrt{4} = 2. \]

\textbf{2.} On cherche $\theta$ tel que :
\[ \cos\theta = \frac{a}{|z|} = \frac{1}{2} \qquad \text{et} \qquad \sin\theta = \frac{b}{|z|} = \frac{\sqrt{3}}{2}. \]
Sur le cercle trigonométrique, le seul angle de $]-\pi\,;\pi]$ vérifiant ces deux conditions est $\theta = \dfrac{\pi}{3}$.

La forme trigonométrique est donc :
\[ z = 2\left(\cos\frac{\pi}{3} + \mathrm{i}\sin\frac{\pi}{3}\right). \]
\textbf{Conclusion :} $|z| = 2$ et $\arg(z) = \dfrac{\pi}{3}\ [2\pi]$.
\end{exoresolu}

\begin{attention}[Cas particuliers à connaître]
\begin{itemize}
\item Si $z$ est un réel strictement positif, $\arg(z) = 0\ [2\pi]$ ; si $z$ est un réel strictement négatif, $\arg(z) = \pi\ [2\pi]$.
\item Si $z = \mathrm{i}b$ avec $b > 0$, $\arg(z) = \dfrac{\pi}{2}\ [2\pi]$ ; avec $b < 0$, $\arg(z) = -\dfrac{\pi}{2}\ [2\pi]$.
\item Le nombre $0$ n'a \textbf{pas} d'argument : la forme trigonométrique n'existe que pour $z \neq 0$.
\end{itemize}
\end{attention}

\courssub{Méthode 3 : Utiliser la forme exponentielle pour multiplier, diviser, élever à une puissance}

\begin{definition}[Forme exponentielle]
Pour tout réel $\theta$, on pose $\mathrm{e}^{\mathrm{i}\theta} = \cos\theta + \mathrm{i}\sin\theta$ (\cle{notation d'Euler}). Tout nombre complexe non nul $z$ de module $r$ et d'argument $\theta$ s'écrit :
\[ z = r\,\mathrm{e}^{\mathrm{i}\theta} \qquad \text{(\cle{forme exponentielle})}. \]
\end{definition}

\begin{propriete}[Règles de calcul avec la forme exponentielle]
Pour tous réels $\theta$, $\theta'$, tous réels strictement positifs $r$, $r'$ et tout entier $n$ :
\[ r\mathrm{e}^{\mathrm{i}\theta} \times r'\mathrm{e}^{\mathrm{i}\theta'} = rr'\,\mathrm{e}^{\mathrm{i}(\theta+\theta')}, \qquad \frac{r\mathrm{e}^{\mathrm{i}\theta}}{r'\mathrm{e}^{\mathrm{i}\theta'}} = \frac{r}{r'}\,\mathrm{e}^{\mathrm{i}(\theta-\theta')}, \qquad \left(r\mathrm{e}^{\mathrm{i}\theta}\right)^n = r^n\,\mathrm{e}^{\mathrm{i}n\theta}. \]
En particulier : $|zz'| = |z|\,|z'|$, $\arg(zz') = \arg(z) + \arg(z')\ [2\pi]$, et $\left|\dfrac{z}{z'}\right| = \dfrac{|z|}{|z'|}$, $\arg\left(\dfrac{z}{z'}\right) = \arg(z) - \arg(z')\ [2\pi]$.
\end{propriete}

\begin{methode}[Forme exponentielle et calculs]
\begin{enumerate}
\item Mettre chaque nombre sous la forme $r\,\mathrm{e}^{\mathrm{i}\theta}$ (module + argument).
\item Appliquer les règles de calcul ci-dessus, identiques à celles des puissances.
\item Réduire l'argument obtenu dans $]-\pi\,;\pi]$ si demandé, puis revenir éventuellement à la forme algébrique avec les valeurs remarquables.
\end{enumerate}
\end{methode}

\begin{exoresolu}[Produit, quotient et puissance avec la forme exponentielle]
On donne $z_1 = 2\,\mathrm{e}^{\mathrm{i}\frac{\pi}{6}}$ et $z_2 = 3\,\mathrm{e}^{\mathrm{i}\frac{\pi}{4}}$.
\begin{enumerate}
\item Déterminer la forme exponentielle de $z_1 z_2$ et de $\dfrac{z_2}{z_1}$.
\item Déterminer la forme algébrique de $z_1^3$.
\end{enumerate}
\tcblower
\textbf{1.} Pour le produit, on multiplie les modules et on additionne les arguments :
\[ z_1 z_2 = 2 \times 3 \times \mathrm{e}^{\mathrm{i}\left(\frac{\pi}{6}+\frac{\pi}{4}\right)} = 6\,\mathrm{e}^{\mathrm{i}\frac{2\pi+3\pi}{12}} = 6\,\mathrm{e}^{\mathrm{i}\frac{5\pi}{12}}. \]
Pour le quotient, on divise les modules et on soustrait les arguments :
\[ \frac{z_2}{z_1} = \frac{3}{2}\,\mathrm{e}^{\mathrm{i}\left(\frac{\pi}{4}-\frac{\pi}{6}\right)} = \frac{3}{2}\,\mathrm{e}^{\mathrm{i}\frac{3\pi-2\pi}{12}} = \frac{3}{2}\,\mathrm{e}^{\mathrm{i}\frac{\pi}{12}}. \]

\textbf{2.} D'après la formule des puissances (formule de Moivre) :
\[ z_1^3 = 2^3\,\mathrm{e}^{\mathrm{i}\,3\times\frac{\pi}{6}} = 8\,\mathrm{e}^{\mathrm{i}\frac{\pi}{2}}. \]
Or $\mathrm{e}^{\mathrm{i}\frac{\pi}{2}} = \cos\dfrac{\pi}{2} + \mathrm{i}\sin\dfrac{\pi}{2} = 0 + \mathrm{i} = \mathrm{i}$.
Donc $z_1^3 = 8\mathrm{i}$.

\textbf{Conclusion :} $z_1 z_2 = 6\,\mathrm{e}^{\mathrm{i}\frac{5\pi}{12}}$, $\dfrac{z_2}{z_1} = \dfrac{3}{2}\,\mathrm{e}^{\mathrm{i}\frac{\pi}{12}}$ et $z_1^3 = 8\mathrm{i}$.
\end{exoresolu}

\begin{aretenir}[Formule de Moivre]
Pour tout réel $\theta$ et tout entier naturel $n$ :
\[ \left(\cos\theta + \mathrm{i}\sin\theta\right)^n = \cos(n\theta) + \mathrm{i}\sin(n\theta). \]
C'est une conséquence directe de $\left(\mathrm{e}^{\mathrm{i}\theta}\right)^n = \mathrm{e}^{\mathrm{i}n\theta}$. Elle sert à exprimer $\cos(n\theta)$ et $\sin(n\theta)$ en fonction de $\cos\theta$ et $\sin\theta$.
\end{aretenir}

\courssub{Méthode 4 : Linéariser une puissance de cosinus ou de sinus}

\begin{propriete}[Formules d'Euler]
Pour tout réel $\theta$ :
\[ \cos\theta = \frac{\mathrm{e}^{\mathrm{i}\theta}+\mathrm{e}^{-\mathrm{i}\theta}}{2} \qquad \text{et} \qquad \sin\theta = \frac{\mathrm{e}^{\mathrm{i}\theta}-\mathrm{e}^{-\mathrm{i}\theta}}{2\mathrm{i}}. \]
Ces formules se déduisent de $\mathrm{e}^{\mathrm{i}\theta} = \cos\theta + \mathrm{i}\sin\theta$ et $\mathrm{e}^{-\mathrm{i}\theta} = \cos\theta - \mathrm{i}\sin\theta$ par addition et soustraction. La démonstration est exigible au Probatoire.
\end{propriete}

\begin{methode}[Linéarisation avec les formules d'Euler]
\emph{Linéariser} $\cos^n x$ ou $\sin^n x$ signifie l'écrire comme une somme de termes en $\cos(kx)$ et $\sin(kx)$, sans puissance.
\begin{enumerate}
\item Écrire les \cle{formules d'Euler} : $\cos x = \dfrac{\mathrm{e}^{\mathrm{i}x}+\mathrm{e}^{-\mathrm{i}x}}{2}$ et $\sin x = \dfrac{\mathrm{e}^{\mathrm{i}x}-\mathrm{e}^{-\mathrm{i}x}}{2\mathrm{i}}$.
\item Élever à la puissance voulue et développer avec la formule du binôme (pour les petites puissances : $(a+b)^2 = a^2+2ab+b^2$, $(a+b)^3 = a^3+3a^2b+3ab^2+b^3$).
\item Regrouper les termes deux à deux : $\mathrm{e}^{\mathrm{i}kx} + \mathrm{e}^{-\mathrm{i}kx} = 2\cos(kx)$ et $\mathrm{e}^{\mathrm{i}kx} - \mathrm{e}^{-\mathrm{i}kx} = 2\mathrm{i}\sin(kx)$.
\item Simplifier pour obtenir une somme de $\cos(kx)$ et $\sin(kx)$ : c'est la \cle{linéarisation}.
\end{enumerate}
\end{methode}

\begin{exoresolu}[Linéariser $\cos^2 x$ et $\sin^3 x$]
\begin{enumerate}
\item Démontrer que $\cos^2 x = \dfrac{1 + \cos(2x)}{2}$.
\item Linéariser $\sin^3 x$.
\end{enumerate}
\tcblower
\textbf{1.} D'après les formules d'Euler, $\cos x = \dfrac{\mathrm{e}^{\mathrm{i}x}+\mathrm{e}^{-\mathrm{i}x}}{2}$. Donc :
\begin{align*}
\cos^2 x &= \left(\frac{\mathrm{e}^{\mathrm{i}x}+\mathrm{e}^{-\mathrm{i}x}}{2}\right)^2 = \frac{\mathrm{e}^{2\mathrm{i}x} + 2\,\mathrm{e}^{\mathrm{i}x}\mathrm{e}^{-\mathrm{i}x} + \mathrm{e}^{-2\mathrm{i}x}}{4} \\
&= \frac{\mathrm{e}^{2\mathrm{i}x} + 2 + \mathrm{e}^{-2\mathrm{i}x}}{4} = \frac{1}{2}\times\frac{\mathrm{e}^{2\mathrm{i}x}+\mathrm{e}^{-2\mathrm{i}x}}{2} + \frac{2}{4} = \frac{\cos(2x)}{2} + \frac{1}{2}.
\end{align*}
D'où $\cos^2 x = \dfrac{1+\cos(2x)}{2}$.

\textbf{2.} De même, $\sin x = \dfrac{\mathrm{e}^{\mathrm{i}x}-\mathrm{e}^{-\mathrm{i}x}}{2\mathrm{i}}$, donc, en remarquant que $(2\mathrm{i})^3 = 8\mathrm{i}^3 = -8\mathrm{i}$ :
\begin{align*}
\sin^3 x &= \left(\frac{\mathrm{e}^{\mathrm{i}x}-\mathrm{e}^{-\mathrm{i}x}}{2\mathrm{i}}\right)^3 = \frac{\mathrm{e}^{3\mathrm{i}x} - 3\mathrm{e}^{\mathrm{i}x} + 3\mathrm{e}^{-\mathrm{i}x} - \mathrm{e}^{-3\mathrm{i}x}}{-8\mathrm{i}} \\
&= \frac{\left(\mathrm{e}^{3\mathrm{i}x}-\mathrm{e}^{-3\mathrm{i}x}\right) - 3\left(\mathrm{e}^{\mathrm{i}x}-\mathrm{e}^{-\mathrm{i}x}\right)}{-8\mathrm{i}} = \frac{2\mathrm{i}\sin(3x) - 6\mathrm{i}\sin x}{-8\mathrm{i}}.
\end{align*}
En simplifiant par $2\mathrm{i}$ : $\sin^3 x = \dfrac{-\sin(3x) + 3\sin x}{4} = \dfrac{3\sin x - \sin(3x)}{4}$.

\textbf{Conclusion :} $\cos^2 x = \dfrac{1+\cos(2x)}{2}$ et $\sin^3 x = \dfrac{3\sin x - \sin(3x)}{4}$.
\end{exoresolu}

\begin{attention}[Pièges de la linéarisation]
\begin{itemize}
\item Ne pas oublier d'élever aussi le \textbf{dénominateur} à la puissance : $(2\mathrm{i})^3 = -8\mathrm{i}$ et non $8\mathrm{i}$, car $\mathrm{i}^3 = -\mathrm{i}$.
\item Dans le développement de $(a-b)^3 = a^3 - 3a^2b + 3ab^2 - b^3$, les signes alternent : une erreur de signe fausse tout le résultat.
\item Vérifier à la fin que le résultat est bien \textbf{réel} : il ne doit rester aucun $\mathrm{i}$ dans une linéarisation de $\cos^n x$ ou $\sin^n x$.
\end{itemize}
\end{attention}

\courssub{Méthode 5 : Déterminer les racines n-ièmes d'un nombre complexe non nul}

\begin{definition}[Racine n-ième]
Soit $Z$ un nombre complexe non nul et $n$ un entier naturel, $n \geq 2$. On appelle \cle{racine n-ième} de $Z$ tout nombre complexe $z$ tel que $z^n = Z$.
\end{definition}

\begin{propriete}[Nombre et expression des racines n-ièmes]
Tout nombre complexe non nul $Z = R\,\mathrm{e}^{\mathrm{i}\varphi}$ admet exactement $n$ racines n-ièmes distinctes, données par :
\[ z_k = \sqrt[n]{R}\ \mathrm{e}^{\mathrm{i}\frac{\varphi + 2k\pi}{n}}, \qquad k \in \{0\,;1\,;\ldots\,;n-1\}. \]
Leurs images sont les sommets d'un \cle{polygone régulier} à $n$ côtés, inscrit dans le cercle de centre $O$ et de rayon $\sqrt[n]{R}$.
\end{propriete}

\begin{methode}[Racines n-ièmes]
Pour résoudre $z^n = Z$ avec $Z \neq 0$ :
\begin{enumerate}
\item Écrire $Z$ sous forme exponentielle : $Z = R\,\mathrm{e}^{\mathrm{i}\varphi}$.
\item Poser $z = r\,\mathrm{e}^{\mathrm{i}\theta}$. L'équation devient $r^n \mathrm{e}^{\mathrm{i}n\theta} = R\,\mathrm{e}^{\mathrm{i}\varphi}$.
\item Identifier : $r = \sqrt[n]{R}$ (racine réelle positive) et $n\theta = \varphi + 2k\pi$, soit $\theta = \dfrac{\varphi + 2k\pi}{n}$.
\item Donner les $n$ solutions en prenant $k = 0, 1, \ldots, n-1$.
\item Conclure par l'ensemble $S$ des solutions et, si demandé, placer les points images (polygone régulier).
\end{enumerate}
\end{methode}

\begin{exoresolu}[Racines cubiques de $-8$]
Résoudre dans $\mathbb{C}$ l'équation $z^3 = -8$. Donner les solutions sous forme exponentielle puis sous forme algébrique.
\tcblower
\textbf{Étape 1 :} écrivons $-8$ sous forme exponentielle. On a $|-8| = 8$ et $\arg(-8) = \pi\ [2\pi]$, donc $-8 = 8\,\mathrm{e}^{\mathrm{i}\pi}$.

\textbf{Étape 2 :} posons $z = r\,\mathrm{e}^{\mathrm{i}\theta}$. Alors $z^3 = r^3\mathrm{e}^{3\mathrm{i}\theta} = 8\,\mathrm{e}^{\mathrm{i}\pi}$.

\textbf{Étape 3 :} par identification, $r^3 = 8$ donc $r = \sqrt[3]{8} = 2$, et $3\theta = \pi + 2k\pi$ donc $\theta = \dfrac{\pi + 2k\pi}{3}$.

\textbf{Étape 4 :} les trois solutions s'obtiennent pour $k = 0, 1, 2$ :
\begin{align*}
z_0 &= 2\,\mathrm{e}^{\mathrm{i}\frac{\pi}{3}} = 2\left(\cos\frac{\pi}{3}+\mathrm{i}\sin\frac{\pi}{3}\right) = 2\left(\frac{1}{2}+\mathrm{i}\frac{\sqrt{3}}{2}\right) = 1 + \mathrm{i}\sqrt{3}, \\
z_1 &= 2\,\mathrm{e}^{\mathrm{i}\frac{\pi+2\pi}{3}} = 2\,\mathrm{e}^{\mathrm{i}\pi} = 2(\cos\pi + \mathrm{i}\sin\pi) = -2, \\
z_2 &= 2\,\mathrm{e}^{\mathrm{i}\frac{\pi+4\pi}{3}} = 2\,\mathrm{e}^{\mathrm{i}\frac{5\pi}{3}} = 2\left(\cos\frac{5\pi}{3}+\mathrm{i}\sin\frac{5\pi}{3}\right) = 2\left(\frac{1}{2}-\mathrm{i}\frac{\sqrt{3}}{2}\right) = 1 - \mathrm{i}\sqrt{3}.
\end{align*}

\textbf{Vérification :} $(-2)^3 = -8$, ce qui confirme que $z_1 = -2$ est bien solution.

\textbf{Conclusion :} l'équation $z^3 = -8$ admet trois solutions : $S = \left\{-2\ ;\ 1+\mathrm{i}\sqrt{3}\ ;\ 1-\mathrm{i}\sqrt{3}\right\}$. Les trois points images forment un triangle équilatéral inscrit dans le cercle de centre $O$ et de rayon $2$.
\end{exoresolu}

\begin{popfigure}
\begin{center}
\begin{tikzpicture}[scale=1.15]
\draw[popline] (-2.7,0) -- (2.9,0);
\draw[popline] (0,-2.4) -- (0,2.4);
\draw[popblue, thick] (0,0) circle (2);
\fill[poporange] (1,1.732) circle (2.2pt) node[above right, poporange] {$z_0 = 1+\mathrm{i}\sqrt{3}$};
\fill[popgreen] (-2,0) circle (2.2pt) node[below left, popgreen!60!black] {$z_1 = -2$};
\fill[poppurple] (1,-1.732) circle (2.2pt) node[below right, poppurple] {$z_2 = 1-\mathrm{i}\sqrt{3}$};
\draw[popink, thick] (1,1.732) -- (-2,0) -- (1,-1.732) -- cycle;
\draw[->,popmuted] (0.7,0) arc (0:60:0.7) node[midway, right, scale=0.8] {$\frac{\pi}{3}$};
\node[below left, popdark] at (0,0) {$O$};
\end{tikzpicture}
\end{center}
\legende{Les trois racines cubiques de $-8$ : sommets d'un triangle équilatéral inscrit dans le cercle de rayon $2$.}
\end{popfigure}

\courssub{Méthode 6 : Appliquer les nombres complexes à une situation technique}

\begin{methode}[Modéliser une situation concrète]
\begin{enumerate}
\item Lire attentivement la situation (électricité, mécanique, topographie) et identifier ce qui joue le rôle de module (intensité, distance, amplitude) et d'argument (déphasage, angle de rotation).
\item Traduire les données en nombres complexes : affixe, forme exponentielle.
\item Choisir l'opération adaptée : une rotation de centre $O$ et d'angle $\theta$ se traduit par une \textbf{multiplication} par $\mathrm{e}^{\mathrm{i}\theta}$ ; une translation se traduit par une \textbf{addition}.
\item Effectuer les calculs rigoureusement, puis \textbf{interpréter le résultat} dans le contexte (unité, phrase de conclusion concrète).
\end{enumerate}
\end{methode}

\begin{exoresolu}[Rotation d'une pièce sur une machine à commande numérique]
Sur une fraiseuse à commande numérique, la position d'un outil est repérée dans le plan complexe par le point $M$ d'affixe $z_M = 3 + \mathrm{i}\sqrt{3}$ (unité : cm). La pièce subit une rotation de centre $O$ et d'angle $\dfrac{\pi}{6}$.
\begin{enumerate}
\item Écrire $z_M$ sous forme exponentielle.
\item Déterminer l'affixe $z_{M'}$ du point $M'$, image de $M$ par cette rotation, sous forme exponentielle puis algébrique.
\end{enumerate}
\tcblower
\textbf{1.} Module : $|z_M| = \sqrt{3^2 + (\sqrt{3})^2} = \sqrt{9+3} = \sqrt{12} = 2\sqrt{3}$.
Argument : on cherche $\theta$ tel que $\cos\theta = \dfrac{3}{2\sqrt{3}} = \dfrac{3\sqrt{3}}{6} = \dfrac{\sqrt{3}}{2}$ et $\sin\theta = \dfrac{\sqrt{3}}{2\sqrt{3}} = \dfrac{1}{2}$. D'où $\theta = \dfrac{\pi}{6}$.
Donc $z_M = 2\sqrt{3}\,\mathrm{e}^{\mathrm{i}\frac{\pi}{6}}$.

\textbf{2.} Une rotation de centre $O$ et d'angle $\dfrac{\pi}{6}$ se traduit par la multiplication par $\mathrm{e}^{\mathrm{i}\frac{\pi}{6}}$ :
\[ z_{M'} = z_M \times \mathrm{e}^{\mathrm{i}\frac{\pi}{6}} = 2\sqrt{3}\,\mathrm{e}^{\mathrm{i}\left(\frac{\pi}{6}+\frac{\pi}{6}\right)} = 2\sqrt{3}\,\mathrm{e}^{\mathrm{i}\frac{\pi}{3}}. \]
Forme algébrique :
\[ z_{M'} = 2\sqrt{3}\left(\cos\frac{\pi}{3}+\mathrm{i}\sin\frac{\pi}{3}\right) = 2\sqrt{3}\left(\frac{1}{2}+\mathrm{i}\frac{\sqrt{3}}{2}\right) = \sqrt{3} + 3\mathrm{i}. \]
\textbf{Conclusion :} après rotation, l'outil se trouve au point $M'$ de coordonnées $(\sqrt{3}\,;3)$, soit environ $(1{,}73\ \text{cm}\,;3\ \text{cm})$. La distance à l'origine est conservée ($|z_{M'}| = 2\sqrt{3} = |z_M|$), ce qui est cohérent avec une rotation.
\end{exoresolu}

\begin{savaistu}[Les complexes, outil quotidien de l'électrotechnicien]
En électricité, un courant alternatif est représenté par un nombre complexe : son module donne l'intensité efficace et son argument le \emph{déphasage} par rapport à la tension. L'\emph{impédance} d'un circuit (résistance $R$, bobine, condensateur) s'écrit $Z = R + \mathrm{i}X$ : multiplier ou diviser des impédances revient à appliquer les règles vues dans cette leçon. C'est pour cela que la forme exponentielle, introduite par le mathématicien suisse Leonhard Euler au XVIII\textsuperscript{e} siècle, est devenue le langage universel des techniciens et des ingénieurs.
\end{savaistu}

\begin{attention}[Les 5 erreurs qui coûtent des points au Probatoire]
\begin{enumerate}
\item \textbf{Oublier la condition $z \neq 0$ pour l'argument.} Le nombre $0$ n'a pas d'argument : ne jamais écrire $\arg(0)$ ni chercher une forme trigonométrique de $0$.
\item \textbf{Déterminer l'argument avec le cosinus seul.} $\cos\theta = \frac{1}{2}$ correspond à $\theta = \frac{\pi}{3}$ \emph{ou} $\theta = -\frac{\pi}{3}$ : il faut vérifier le signe du sinus pour choisir. Les deux conditions sont indispensables.
\item \textbf{Confondre module et argument dans les opérations.} Pour un produit : on \emph{multiplie} les modules et on \emph{additionne} les arguments. Écrire $|z z'| = |z| + |z'|$ est faux.
\item \textbf{Oublier des racines dans $z^n = Z$.} Il y a exactement $n$ solutions, obtenues avec $k = 0, 1, \ldots, n-1$. S'arrêter à la solution évidente (par exemple $z = -2$ pour $z^3 = -8$) fait perdre la moitié des points.
\item \textbf{Mal rédiger la conclusion.} Une équation se termine par l'ensemble des solutions $S = \{\ldots\}$ ; un problème concret se termine par une phrase avec l'unité. Toute réponse non justifiée (calculs absents, formule non citée) est pénalisée.
\end{enumerate}
\end{attention}

\begin{exercice}[Entraînement au Probatoire]
\begin{enumerate}
\item Mettre sous forme trigonométrique puis exponentielle : $z = -1 + \mathrm{i}$ et $z' = \sqrt{3} - \mathrm{i}$.
\item On donne $a = 2\,\mathrm{e}^{\mathrm{i}\frac{\pi}{3}}$ et $b = \mathrm{e}^{-\mathrm{i}\frac{\pi}{4}}$. Déterminer la forme exponentielle de $ab$, de $\dfrac{a}{b}$ et la forme algébrique de $b^4$.
\item Linéariser $\sin^2 x$ et $\cos^3 x$.
\item Résoudre dans $\mathbb{C}$ l'équation $z^4 = 16$ et représenter les solutions dans le plan complexe.
\end{enumerate}
\end{exercice}

\begin{corrige}[Exercice : Entraînement au Probatoire]
\textbf{1.} Pour $z = -1 + \mathrm{i}$ : $|z| = \sqrt{1+1} = \sqrt{2}$ ; $\cos\theta = -\dfrac{1}{\sqrt{2}} = -\dfrac{\sqrt{2}}{2}$ et $\sin\theta = \dfrac{\sqrt{2}}{2}$, donc $\theta = \dfrac{3\pi}{4}$. Ainsi $z = \sqrt{2}\left(\cos\dfrac{3\pi}{4} + \mathrm{i}\sin\dfrac{3\pi}{4}\right) = \sqrt{2}\,\mathrm{e}^{\mathrm{i}\frac{3\pi}{4}}$.
Pour $z' = \sqrt{3} - \mathrm{i}$ : $|z'| = \sqrt{3+1} = 2$ ; $\cos\theta' = \dfrac{\sqrt{3}}{2}$ et $\sin\theta' = -\dfrac{1}{2}$, donc $\theta' = -\dfrac{\pi}{6}$. Ainsi $z' = 2\left(\cos\dfrac{\pi}{6} - \mathrm{i}\sin\dfrac{\pi}{6}\right) = 2\,\mathrm{e}^{-\mathrm{i}\frac{\pi}{6}}$.

\textbf{2.} $ab = 2\,\mathrm{e}^{\mathrm{i}\left(\frac{\pi}{3}-\frac{\pi}{4}\right)} = 2\,\mathrm{e}^{\mathrm{i}\frac{\pi}{12}}$ ; $\dfrac{a}{b} = 2\,\mathrm{e}^{\mathrm{i}\left(\frac{\pi}{3}+\frac{\pi}{4}\right)} = 2\,\mathrm{e}^{\mathrm{i}\frac{7\pi}{12}}$ ; $b^4 = \mathrm{e}^{-\mathrm{i}\pi} = \cos(-\pi) + \mathrm{i}\sin(-\pi) = -1$.

\textbf{3.} $\sin^2 x = \left(\dfrac{\mathrm{e}^{\mathrm{i}x}-\mathrm{e}^{-\mathrm{i}x}}{2\mathrm{i}}\right)^2 = \dfrac{\mathrm{e}^{2\mathrm{i}x} - 2 + \mathrm{e}^{-2\mathrm{i}x}}{-4} = \dfrac{2 - 2\cos(2x)}{4} = \dfrac{1-\cos(2x)}{2}$.
$\cos^3 x = \left(\dfrac{\mathrm{e}^{\mathrm{i}x}+\mathrm{e}^{-\mathrm{i}x}}{2}\right)^3 = \dfrac{\mathrm{e}^{3\mathrm{i}x} + 3\mathrm{e}^{\mathrm{i}x} + 3\mathrm{e}^{-\mathrm{i}x} + \mathrm{e}^{-3\mathrm{i}x}}{8} = \dfrac{2\cos(3x) + 6\cos x}{8} = \dfrac{3\cos x + \cos(3x)}{4}$.

\textbf{4.} $16 = 16\,\mathrm{e}^{\mathrm{i}\,0}$, donc $z_k = \sqrt[4]{16}\,\mathrm{e}^{\mathrm{i}\frac{2k\pi}{4}} = 2\,\mathrm{e}^{\mathrm{i}\frac{k\pi}{2}}$ pour $k = 0, 1, 2, 3$ : $z_0 = 2$, $z_1 = 2\mathrm{i}$, $z_2 = -2$, $z_3 = -2\mathrm{i}$. Donc $S = \{2\,;\,2\mathrm{i}\,;\,-2\,;\,-2\mathrm{i}\}$ : les images forment un carré inscrit dans le cercle de centre $O$ et de rayon $2$.
\end{corrige}

\newpage

\coursec{Exercices}

\coursec{Nombres complexes : Approche géométrique}

\courssub{Je vérifie mes connaissances}

\begin{exercice}[QCM : forme exponentielle]
Pour chaque question, une seule réponse est exacte. Entourer la bonne réponse.
\begin{enumerate}
\item Le module de $z = -3 + 3i$ est :
\begin{enumerate}
\item $3$ \quad \item $3\sqrt{2}$ \quad \item $6$ \quad \item $9$
\end{enumerate}
\item Un argument de $z = 2e^{i\pi/6}$ est :
\begin{enumerate}
\item $2$ \quad \item $\dfrac{\pi}{6}$ \quad \item $-\dfrac{\pi}{6}$ \quad \item $\dfrac{5\pi}{6}$
\end{enumerate}
\item Si $z = e^{i\pi/3}$, alors $z^3$ vaut :
\begin{enumerate}
\item $1$ \quad \item $-1$ \quad \item $i$ \quad \item $-i$
\end{enumerate}
\item Le nombre de racines quatrièmes distinctes d'un nombre complexe non nul est :
\begin{enumerate}
\item $2$ \quad \item $3$ \quad \item $4$ \quad \item une infinité
\end{enumerate}
\end{enumerate}
\end{exercice}

\begin{exercice}[Vrai ou faux, en justifiant]
Dire si chaque affirmation est vraie ou fausse, et justifier la réponse.
\begin{enumerate}
\item Tout nombre complexe non nul admet exactement un argument.
\item Si $z = re^{i\theta}$ avec $r > 0$, alors $|z| = r$.
\item Les racines $n$-ièmes d'un nombre complexe non nul sont les sommets d'un polygone régulier à $n$ côtés.
\item Pour tout réel $\theta$, $\cos^2\theta = \dfrac{1 + \cos(2\theta)}{2}$.
\end{enumerate}
\end{exercice}

\begin{exercice}[Définitions et vocabulaire]
\begin{enumerate}
\item Donner la définition de la forme exponentielle d'un nombre complexe non nul.
\item Qu'appelle-t-on affixe d'un point dans le plan complexe ? Image d'un nombre complexe ?
\item Énoncer la formule de Moivre.
\item Compléter : les racines $n$-ièmes de $z = re^{i\theta}$ sont les nombres $z_k = \ldots$ pour $k \in \{\ldots\}$.
\end{enumerate}
\end{exercice}

\begin{exercice}[Calculs directs]
\begin{enumerate}
\item Déterminer le module et un argument de $z_1 = 1 + i$ et de $z_2 = -\sqrt{3} + i$.
\item Écrire sous forme exponentielle : $z_3 = 5i$ ; $z_4 = -4$ ; $z_5 = 2e^{-i\pi/4} \times 3e^{i\pi/2}$.
\item Calculer $(1 + i)^4$ en utilisant la forme exponentielle.
\item Donner les deux racines carrées de $z = 4e^{i\pi/3}$.
\end{enumerate}
\end{exercice}

\courssub{Je m'entraîne}

\begin{exercice}[Conversion et représentation]
On considère le nombre complexe $z = -2 + 2i$.
\begin{enumerate}
\item Calculer le module de $z$.
\item Déterminer un argument de $z$.
\item En déduire la forme exponentielle de $z$.
\item Placer dans le plan complexe muni d'un repère orthonormé direct $(O\,;\vec{u},\vec{v})$ le point $M$ d'affixe $z$ (unité : \SI{2}{cm}).
\end{enumerate}
\end{exercice}

\begin{popfigure}
\begin{tikzpicture}[scale=0.7, >=Stealth]
\draw[popline, thin] (-3.5,0) -- (3.5,0) node[right] {$\vec{u}$};
\draw[popline, thin] (0,-3.5) -- (0,3.5) node[above] {$\vec{v}$};
\foreach \x in {-3,-2,-1,1,2,3} \draw (\x,0.1) -- (\x,-0.1) node[below, font=\tiny] {\x};
\foreach \y in {-3,-2,-1,1,2,3} \draw (0.1,\y) -- (-0.1,\y) node[left, font=\tiny] {\y};
\draw[popblue, thick] (0,0) -- (-2,2);
\fill[popblue] (-2,2) circle (2.5pt) node[above left] {$M(-2+2i)$};
\draw[poporange, dashed] (-2,0) node[below, font=\tiny] {$-2$} -- (-2,2);
\draw[poporange, dashed] (0,2) node[left, font=\tiny] {$2$} -- (-2,2);
\draw[popmuted, ->] (0.5,0) arc (0:135:0.5) node[midway, right, font=\small] {$\frac{3\pi}{4}$};
\end{tikzpicture}
\legende{Représentation de $z = -2+2i = 2\sqrt{2}e^{i3\pi/4}$ (unité 2 cm).}
\end{popfigure}

\begin{exercice}[Linéarisation de $\cos^4\theta$]
\begin{enumerate}
\item À l'aide de la formule d'Euler $\cos\theta = \dfrac{e^{i\theta} + e^{-i\theta}}{2}$, montrer que :
\[ \cos^4\theta = \frac{1}{8}\cos(4\theta) + \frac{1}{2}\cos(2\theta) + \frac{3}{8}. \]
\item Vérifier cette égalité pour $\theta = \dfrac{\pi}{3}$.
\item En déduire la valeur exacte de $\cos^4\dfrac{\pi}{8} + \cos^4\dfrac{3\pi}{8}$. (On pourra utiliser $\cos\dfrac{3\pi}{8} = \sin\dfrac{\pi}{8}$.)
\end{enumerate}
\end{exercice}

\begin{exercice}[Racines quatrièmes]
On considère le nombre complexe $Z = 16e^{i\pi/3}$.
\begin{enumerate}
\item Déterminer les quatre racines quatrièmes de $Z$ sous forme exponentielle.
\item Représenter leurs images dans le plan complexe. Quelle figure géométrique obtient-on ?
\item Montrer que $-8\sqrt{3} + 8i = 16e^{i5\pi/6}$.
\item En déduire les solutions dans $\mathbb{C}$ de l'équation $z^4 = -8\sqrt{3} + 8i$.
\end{enumerate}
\end{exercice}

\begin{popfigure}
\begin{tikzpicture}[scale=1.2, >=Stealth]
\draw[popline, thin] (-2.5,0) -- (2.5,0) node[right] {$\vec{u}$};
\draw[popline, thin] (0,-2.5) -- (0,2.5) node[above] {$\vec{v}$};
\foreach \x in {-2,-1,1,2} \draw (\x,0.05) -- (\x,-0.05) node[below, font=\tiny] {\x};
\foreach \y in {-2,-1,1,2} \draw (0.05,\y) -- (-0.05,\y) node[left, font=\tiny] {\y};
\draw[popblueL] (0,0) circle (2);
% Roots: 2 * exp(i(pi/12 + k*pi/2))
\foreach \k/\angle in {0/15, 1/105, 2/195, 3/285} {
    \pgfmathsetmacro{\x}{2*cos(\angle)}
    \pgfmathsetmacro{\y}{2*sin(\angle)}
    \fill[popblue] (\x,\y) circle (2pt);
    \draw[popblue, thick] (0,0) -- (\x,\y);
    \node[popblue, font=\tiny] at (1.2*\x,1.2*\y) {$z_{\k}$};
}
\draw[poporange, thick] (1.932,0.518) -- (-0.518,1.932) -- (-1.932,-0.518) -- (0.518,-1.932) -- cycle;
\end{tikzpicture}
\legende{Racines quatrièmes de $16e^{i\pi/3}$ : sommets d'un carré de centre $O$ et rayon $2$.}
\end{popfigure}

\begin{exercice}[Formule de Moivre par récurrence]
On veut démontrer que pour tout entier $n \geq 1$ et tout réel $\theta$ :
\[ (\cos\theta + i\sin\theta)^n = \cos(n\theta) + i\sin(n\theta). \]
\begin{enumerate}
\item Vérifier la propriété pour $n = 1$.
\item On suppose la propriété vraie au rang $n$ (hypothèse de récurrence). En développant $(\cos\theta + i\sin\theta)^{n+1} = (\cos\theta + i\sin\theta)^n(\cos\theta + i\sin\theta)$ et en utilisant les formules d'addition, montrer qu'elle est vraie au rang $n+1$.
\item Conclure.
\item Application : exprimer $\cos(3\theta)$ en fonction de $\cos\theta$.
\end{enumerate}
\end{exercice}

\courssub{Je me prépare au Probatoire}

\begin{exercice}[Étude complète d'une configuration du plan complexe]
Le plan est muni d'un repère orthonormé direct $(O\,;\vec{u},\vec{v})$ d'unité graphique \SI{2}{cm}. On considère les points $A$, $B$ et $C$ d'affixes respectives :
\[ z_A = 2e^{i\pi/6}, \qquad z_B = 2e^{i5\pi/6}, \qquad z_C = z_A \times z_B. \]
\begin{enumerate}
\item Écrire $z_A$ et $z_B$ sous forme algébrique.
\item Déterminer le module et un argument de $z_C$, puis écrire $z_C$ sous forme algébrique.
\item Placer les points $A$, $B$ et $C$ dans le repère.
\item Démontrer que le triangle $OAB$ est isocèle et déterminer une mesure de l'angle $\widehat{AOB}$.
\item On pose $z_D = \dfrac{z_B}{z_A}$.
\begin{enumerate}
\item Déterminer la forme exponentielle de $z_D$.
\item En déduire la nature exacte du triangle $OAD$ en justifiant.
\end{enumerate}
\item Calculer $(z_A)^{12}$. Que peut-on en conclure pour le point image de $(z_A)^{12}$ ?
\end{enumerate}
\end{exercice}

\begin{popfigure}
\begin{tikzpicture}[scale=0.7, >=Stealth]
\draw[popline, thin] (-5,0) -- (5,0) node[right] {$\vec{u}$};
\draw[popline, thin] (0,-5) -- (0,5) node[above] {$\vec{v}$};
\foreach \x in {-4,-2,2,4} \draw (\x,0.1) -- (\x,-0.1) node[below, font=\tiny] {\x};
\foreach \y in {-4,-2,2,4} \draw (0.1,\y) -- (-0.1,\y) node[left, font=\tiny] {\y};
% Points
\coordinate (O) at (0,0);
\coordinate (A) at (1.732,1); % sqrt3, 1
\coordinate (B) at (-1.732,1);
\coordinate (C) at (-4,0);
\coordinate (D) at (-0.5,0.866); % cos120, sin120
\draw[popblue, thick] (O) -- (A) node[midway, above right, font=\small] {$2$};
\draw[popblue, thick] (O) -- (B) node[midway, above left, font=\small] {$2$};
\draw[poporange, thick] (A) -- (B);
\draw[popgreen, thick] (O) -- (C) node[midway, below, font=\small] {$4$};
\draw[poppurple, thick] (O) -- (D) node[midway, left, font=\small] {$1$};
\fill[popblue] (A) circle (2.5pt) node[above right] {$A(\sqrt{3}+i)$};
\fill[popblue] (B) circle (2.5pt) node[above left] {$B(-\sqrt{3}+i)$};
\fill[popgreen] (C) circle (2.5pt) node[below] {$C(-4)$};
\fill[poppurple] (D) circle (2.5pt) node[above left] {$D(-\frac{1}{2}+i\frac{\sqrt{3}}{2})$};
\fill[popmuted] (O) circle (2pt) node[below left] {$O$};
% Angle AOB
\draw[popmuted, ->] (0.5,0) arc (0:120:0.5) node[midway, right, font=\small] {$120^\circ$};
% Right angle OAD
\draw[poporange] (0.3,0) -- (0.3,0.3) -- (0,0.3);
\end{tikzpicture}
\legende{Configuration des points $A, B, C, D$ (unité 2 cm). Triangle $OAB$ isocèle en $O$, triangle $OAD$ rectangle en $O$.}
\end{popfigure}

\begin{exercice}[Problème d'électricité : moteur asynchrone triphasé]
Un atelier de Douala utilise un moteur asynchrone alimenté par un réseau triphasé de tension simple efficace $U = \SI{220}{V}$ et de fréquence $f = \SI{50}{Hz}$. En électrotechnique, on associe à toute grandeur sinusoïdale un nombre complexe appelé \emph{phaseur} : la tension d'une phase est modélisée par $\underline{U} = 220\,e^{i\,0} = 220$ (en volts), et l'impédance d'une phase du moteur est $\underline{Z} = 8 + 6i$ (en ohms). Le courant est donné par la loi d'Ohm complexe : $\underline{I} = \dfrac{\underline{U}}{\underline{Z}}$.
\begin{enumerate}
\item Déterminer le module et un argument de $\underline{Z}$ (l'argument sera arrondi au centième de radian).
\item Écrire $\underline{Z}$ sous forme exponentielle.
\item En déduire la forme exponentielle de $\underline{I}$, puis sa forme algébrique (arrondir au dixième).
\item Quelle est la valeur efficace $I$ du courant absorbé par phase ?
\item Le déphasage $\varphi$ du courant par rapport à la tension est $\varphi = \arg(\underline{U}) - \arg(\underline{I})$. Calculer $\cos\varphi$ (facteur de puissance).
\item La puissance active absorbée par le moteur est $P = 3UI\cos\varphi$. Calculer $P$ en watts.
\item Le propriétaire de l'atelier souhaite réduire sa facture : expliquer pourquoi un facteur de puissance proche de $1$ est souhaitable.
\end{enumerate}
\end{exercice}

\begin{exercice}[Racines de l'unité et polygone régulier]
On considère dans $\mathbb{C}$ l'équation $(E) : z^6 = 1$.
\begin{enumerate}
\item Déterminer les six solutions de $(E)$ sous forme exponentielle. On les note $z_k$ pour $k \in \{0, 1, 2, 3, 4, 5\}$ avec des arguments croissants dans $[0\,;2\pi[$.
\item Placer les images $M_k$ des $z_k$ dans le plan complexe muni d'un repère orthonormé $(O\,;\vec{u},\vec{v})$ (unité : \SI{3}{cm}). Quelle figure obtient-on en joignant les points $M_0, M_1, \ldots, M_5$ ?
\item On pose $j = e^{i2\pi/3}$. Montrer que $j$ est solution de $(E)$ et que $1 + j + j^2 = 0$.
\item Écrire $j$ sous forme algébrique et vérifier que $j^2 = \bar{j}$.
\item Un technicien en mécanique usine une pièce cylindrique sur laquelle doivent être percés $6$ trous régulièrement espacés sur un cercle de rayon $R = \SI{5}{cm}$. À l'aide des résultats précédents, calculer la distance exacte entre les centres de deux trous consécutifs.
\end{enumerate}
\end{exercice}

\begin{popfigure}
\begin{tikzpicture}[scale=1.0, >=Stealth]
\draw[popline, thin] (-3.5,0) -- (3.5,0) node[right] {$\vec{u}$};
\draw[popline, thin] (0,-3.5) -- (0,3.5) node[above] {$\vec{v}$};
\foreach \x in {-3,-2,-1,1,2,3} \draw (\x,0.05) -- (\x,-0.05) node[below, font=\tiny] {\x};
\foreach \y in {-3,-2,-1,1,2,3} \draw (0.05,\y) -- (-0.05,\y) node[left, font=\tiny] {\y};
\draw[popblueL] (0,0) circle (3);
\foreach \k/\angle in {0/0, 1/60, 2/120, 3/180, 4/240, 5/300} {
    \pgfmathsetmacro{\x}{3*cos(\angle)}
    \pgfmathsetmacro{\y}{3*sin(\angle)}
    \fill[popblue] (\x,\y) circle (2.5pt);
    \node[popblue, font=\small] at (1.15*\x,1.15*\y) {$M_{\k}$};
    \draw[popblue, thin] (0,0) -- (\x,\y);
}
\draw[poporange, thick] (3,0) -- (1.5,2.598) -- (-1.5,2.598) -- (-3,0) -- (-1.5,-2.598) -- (1.5,-2.598) -- cycle;
\end{tikzpicture}
\legende{Racines $6$-èmes de l'unité : hexagone régulier inscrit dans le cercle unité (échelle 3 cm).}
\end{popfigure}

\newpage

\coursec{Corrigés des exercices}

\coursec{Corrigés détaillés — Nombres complexes : Approche géométrique}

\courssub{Représentation géométrique et forme exponentielle}

\begin{corrige}[Exercice 1 — QCM : forme exponentielle]
\begin{enumerate}
\item Pour $z = -3 + 3i$, le module est $|z| = \sqrt{(-3)^2 + 3^2} = \sqrt{18} = 3\sqrt{2}$. \textbf{Réponse b.}
\item Dans l'écriture $z = 2e^{i\pi/6}$, le module est $2$ et un argument est $\dfrac{\pi}{6}$. \textbf{Réponse b.}
\item $z^3 = \left(e^{i\pi/3}\right)^3 = e^{i\pi} = -1$. \textbf{Réponse b.}
\item Tout nombre complexe non nul admet exactement $n$ racines $n$-ièmes distinctes ; pour $n = 4$, il y en a $4$. \textbf{Réponse c.}
\end{enumerate}
\end{corrige}

\begin{corrige}[Exercice 2 — Vrai ou faux, en justifiant]
\begin{enumerate}
\item \textbf{Faux.} Un nombre complexe non nul admet une infinité d'arguments : si $\theta$ est un argument de $z$, alors $\theta + 2k\pi$ ($k \in \mathbb{Z}$) en est un aussi. On dit que l'argument est défini \emph{modulo $2\pi$}.
\item \textbf{Vrai.} $|re^{i\theta}| = r \times |e^{i\theta}| = r \times 1 = r$, car $r > 0$ et $|e^{i\theta}| = \sqrt{\cos^2\theta + \sin^2\theta} = 1$.
\item \textbf{Vrai.} Les racines $n$-ièmes de $z = re^{i\theta}$ ont toutes le même module $\sqrt[n]{r}$ (elles sont sur un même cercle de centre $O$) et leurs arguments diffèrent de $\dfrac{2\pi}{n}$ : leurs images sont donc les sommets d'un polygone régulier à $n$ côtés.
\item \textbf{Vrai.} C'est une formule de linéarisation (ou duplication) : $\cos(2\theta) = 2\cos^2\theta - 1$, d'où $\cos^2\theta = \dfrac{1 + \cos(2\theta)}{2}$.
\end{enumerate}
\end{corrige}

\begin{corrige}[Exercice 3 — Définitions et vocabulaire]
\begin{enumerate}
\item Soit $z$ un nombre complexe non nul de module $r > 0$ et d'argument $\theta$. La \cle{forme exponentielle} de $z$ est l'écriture $z = re^{i\theta}$, où $e^{i\theta} = \cos\theta + i\sin\theta$ (formule d'Euler).
\item Dans le plan muni d'un repère orthonormé direct $(O\,;\vec{u},\vec{v})$, l'\cle{affixe} du point $M(x\,;y)$ est le nombre complexe $z = x + iy$. Réciproquement, le point $M(x\,;y)$ est l'\cle{image} du nombre complexe $z = x + iy$.
\item \textbf{Formule de Moivre} : pour tout réel $\theta$ et tout entier naturel $n$,
\[ (\cos\theta + i\sin\theta)^n = \cos(n\theta) + i\sin(n\theta). \]
\item Les racines $n$-ièmes de $z = re^{i\theta}$ ($r>0$) sont les nombres
\[ z_k = \sqrt[n]{r}\, e^{i\left(\frac{\theta}{n} + \frac{2k\pi}{n}\right)} \quad \text{pour } k \in \{0, 1, 2, \ldots, n-1\}. \]
\end{enumerate}
\end{corrige}

\begin{corrige}[Exercice 4 — Calculs directs]
\begin{enumerate}
\item Pour $z_1 = 1 + i$ : $|z_1| = \sqrt{1^2 + 1^2} = \sqrt{2}$. On écrit $z_1 = \sqrt{2}\left(\dfrac{\sqrt{2}}{2} + i\dfrac{\sqrt{2}}{2}\right)$, d'où $\cos\theta_1 = \dfrac{\sqrt{2}}{2}$ et $\sin\theta_1 = \dfrac{\sqrt{2}}{2}$, donc $\theta_1 = \dfrac{\pi}{4}$ (premier quadrant).

Pour $z_2 = -\sqrt{3} + i$ : $|z_2| = \sqrt{3 + 1} = 2$. On écrit $z_2 = 2\left(-\dfrac{\sqrt{3}}{2} + i\dfrac{1}{2}\right)$, d'où $\cos\theta_2 = -\dfrac{\sqrt{3}}{2}$ et $\sin\theta_2 = \dfrac{1}{2}$, donc $\theta_2 = \dfrac{5\pi}{6}$ (deuxième quadrant).
\item $z_3 = 5i$ : module $5$, argument $\dfrac{\pi}{2}$, donc $z_3 = 5e^{i\pi/2}$.

$z_4 = -4$ : module $4$, argument $\pi$, donc $z_4 = 4e^{i\pi}$.

$z_5 = 2e^{-i\pi/4} \times 3e^{i\pi/2} = 6\,e^{i\left(-\frac{\pi}{4} + \frac{\pi}{2}\right)} = 6e^{i\pi/4}$ (les modules se multiplient, les arguments s'additionnent modulo $2\pi$).
\item $1 + i = \sqrt{2}\,e^{i\pi/4}$, donc
\[ (1+i)^4 = \left(\sqrt{2}\right)^4 e^{i4 \times \pi/4} = 4\,e^{i\pi} = -4. \]
Vérification : $(1+i)^2 = 2i$, donc $(1+i)^4 = (2i)^2 = -4$. Cohérent.
\item Les racines carrées de $z = 4e^{i\pi/3}$ sont $z_k = \sqrt{4}\,e^{i\left(\frac{\pi}{6} + k\pi\right)}$ pour $k \in \{0,1\}$ :
\[ z_0 = 2e^{i\pi/6} = \sqrt{3} + i \qquad \text{et} \qquad z_1 = 2e^{i7\pi/6} = -\sqrt{3} - i. \]
\end{enumerate}
\end{corrige}

\begin{corrige}[Exercice 5 — Conversion et représentation graphique]
\begin{enumerate}
\item $|z| = \sqrt{(-2)^2 + 2^2} = \sqrt{8} = 2\sqrt{2}$.
\item On factorise le module : $z = 2\sqrt{2}\left(-\dfrac{\sqrt{2}}{2} + i\dfrac{\sqrt{2}}{2}\right)$. On cherche $\theta$ tel que $\cos\theta = -\dfrac{\sqrt{2}}{2}$ et $\sin\theta = \dfrac{\sqrt{2}}{2}$ : $\theta = \dfrac{3\pi}{4}$.
\item Forme exponentielle : $z = 2\sqrt{2}\,e^{i3\pi/4}$.
\item Le point $M$ a pour coordonnées $(-2\,;2)$. Avec une unité de \SI{2}{cm}, on reporte \SI{4}{cm} vers la gauche sur l'axe des réels et \SI{4}{cm} vers le haut sur l'axe des imaginaires.
\begin{popfigure}
\begin{tikzpicture}[scale=0.5, >=Stealth]
\draw[popline, ->] (-6,0) -- (2,0) node[right] {$\mathbb{R}$};
\draw[popline, ->] (0,-2) -- (0,6) node[above] {$\mathbb{I}$};
\foreach \x in {-4,-2,2} \draw (\x,0.1) -- (\x,-0.1) node[below] {\small $\x$};
\foreach \y in {2,4} \draw (0.1,\y) -- (-0.1,\y) node[left] {\small $\y$};
\draw[popblue, thick, ->] (0,0) -- (-4,4) node[midway, above left] {$z$};
\fill[popblue] (-4,4) circle (3pt) node[above right] {$M(-2;2)$};
\draw[dashed, popmuted] (-4,0) node[below] {$-2$} -- (-4,4) -- (0,4) node[left] {$2$};
\draw[poporange, |<->|] (0,0) -- (-4,4) node[midway, above left, sloped] {$OM = 2\sqrt{2}$};
\end{tikzpicture}
\legende{Représentation de $z = -2+2i$ (unité \SI{2}{cm})}
\end{popfigure}
On vérifie graphiquement que $OM = 2\sqrt{2} \approx \num{2,83}$ unités, soit environ \SI{5,7}{cm} sur la figure.
\end{enumerate}
\end{corrige}

\courssub{Linéarisation, formule de Moivre et puissances}

\begin{corrige}[Exercice 6 — Linéarisation de $\cos^4\theta$]
\begin{enumerate}
\item D'après la formule d'Euler, $\cos\theta = \dfrac{e^{i\theta} + e^{-i\theta}}{2}$. Donc :
\begin{align*}
\cos^4\theta &= \frac{1}{16}\left(e^{i\theta} + e^{-i\theta}\right)^4 \\
&= \frac{1}{16}\left(e^{i4\theta} + 4e^{i2\theta} + 6 + 4e^{-i2\theta} + e^{-i4\theta}\right) \\
&= \frac{1}{16}\left[\left(e^{i4\theta} + e^{-i4\theta}\right) + 4\left(e^{i2\theta} + e^{-i2\theta}\right) + 6\right] \\
&= \frac{1}{16}\left[2\cos(4\theta) + 8\cos(2\theta) + 6\right] \\
&= \frac{1}{8}\cos(4\theta) + \frac{1}{2}\cos(2\theta) + \frac{3}{8}.
\end{align*}
On a utilisé le binôme de Newton $(a+b)^4 = a^4 + 4a^3b + 6a^2b^2 + 4ab^3 + b^4$ et la formule d'Euler $\cos x = \dfrac{e^{ix} + e^{-ix}}{2}$.
\item Pour $\theta = \dfrac{\pi}{3}$ : $\cos\dfrac{\pi}{3} = \dfrac{1}{2}$, donc $\cos^4\dfrac{\pi}{3} = \dfrac{1}{16}$.

Membre de droite : $\dfrac{1}{8}\cos\dfrac{4\pi}{3} + \dfrac{1}{2}\cos\dfrac{2\pi}{3} + \dfrac{3}{8} = \dfrac{1}{8}\times\left(-\dfrac{1}{2}\right) + \dfrac{1}{2}\times\left(-\dfrac{1}{2}\right) + \dfrac{3}{8} = -\dfrac{1}{16} - \dfrac{4}{16} + \dfrac{6}{16} = \dfrac{1}{16}$. L'égalité est vérifiée.
\item Comme $\cos\dfrac{3\pi}{8} = \sin\dfrac{\pi}{8}$, on a $\cos^4\dfrac{3\pi}{8} = \sin^4\dfrac{\pi}{8}$. Posons $S = \cos^4\dfrac{\pi}{8} + \sin^4\dfrac{\pi}{8}$. Alors :
\[ S = \left(\cos^2\tfrac{\pi}{8} + \sin^2\tfrac{\pi}{8}\right)^2 - 2\cos^2\tfrac{\pi}{8}\sin^2\tfrac{\pi}{8} = 1 - \frac{1}{2}\sin^2\tfrac{\pi}{4} = 1 - \frac{1}{2}\times\frac{1}{2} = \frac{3}{4}. \]
Vérification par la formule du 1. appliquée à $\theta = \dfrac{\pi}{8}$ et $\theta = \dfrac{3\pi}{8}$ :
\[ S = \frac{1}{8}\left[\cos\tfrac{\pi}{2} + \cos\tfrac{3\pi}{2}\right] + \frac{1}{2}\left[\cos\tfrac{\pi}{4} + \cos\tfrac{3\pi}{4}\right] + \frac{3}{4} = 0 + \frac{1}{2}\times 0 + \frac{3}{4} = \frac{3}{4}. \]
\end{enumerate}
\end{corrige}

\begin{corrige}[Exercice 8 — Formule de Moivre par récurrence]
\begin{enumerate}
\item \textbf{Initialisation} : pour $n = 1$, $(\cos\theta + i\sin\theta)^1 = \cos\theta + i\sin\theta = \cos(1\times\theta) + i\sin(1\times\theta)$. La propriété est vraie au rang $1$.
\item \textbf{Hérédité} : supposons la propriété vraie au rang $n \geq 1$, c'est-à-dire $(\cos\theta + i\sin\theta)^n = \cos(n\theta) + i\sin(n\theta)$. Alors :
\begin{align*}
(\cos\theta + i\sin\theta)^{n+1} &= (\cos\theta + i\sin\theta)^n(\cos\theta + i\sin\theta) \\
&= [\cos(n\theta) + i\sin(n\theta)][\cos\theta + i\sin\theta] \\
&= \cos(n\theta)\cos\theta - \sin(n\theta)\sin\theta \\
&\quad + i[\sin(n\theta)\cos\theta + \cos(n\theta)\sin\theta] \\
&= \cos(n\theta + \theta) + i\sin(n\theta + \theta) \\
&= \cos((n+1)\theta) + i\sin((n+1)\theta),
\end{align*}
en utilisant les formules d'addition $\cos(a+b) = \cos a\cos b - \sin a\sin b$ et $\sin(a+b) = \sin a\cos b + \cos a\sin b$. La propriété est donc vraie au rang $n+1$.
\item \textbf{Conclusion} : la propriété est vraie au rang $1$ et héréditaire ; par récurrence, elle est vraie pour tout entier $n \geq 1$.
\item D'après Moivre, $\cos(3\theta) + i\sin(3\theta) = (\cos\theta + i\sin\theta)^3$. En développant :
\[ (\cos\theta + i\sin\theta)^3 = \cos^3\theta + 3i\cos^2\theta\sin\theta - 3\cos\theta\sin^2\theta - i\sin^3\theta. \]
En identifiant les parties réelles : $\cos(3\theta) = \cos^3\theta - 3\cos\theta\sin^2\theta = \cos^3\theta - 3\cos\theta(1 - \cos^2\theta)$, d'où :
\[ \boxed{\cos(3\theta) = 4\cos^3\theta - 3\cos\theta.} \]
\end{enumerate}
\end{corrige}

\courssub{Racines n-ièmes et configurations géométriques}

\begin{corrige}[Exercice 7 — Racines quatrièmes]
\begin{enumerate}
\item Les racines quatrièmes de $Z = 16e^{i\pi/3}$ sont :
\[ z_k = 16^{1/4}\,e^{i\left(\frac{\pi}{12} + \frac{2k\pi}{4}\right)} = 2\,e^{i\left(\frac{\pi}{12} + \frac{k\pi}{2}\right)}, \quad k \in \{0,1,2,3\}. \]
Soit : $z_0 = 2e^{i\pi/12}$, $z_1 = 2e^{i7\pi/12}$, $z_2 = 2e^{i13\pi/12}$, $z_3 = 2e^{i19\pi/12}$.
\item Les quatre images sont sur le cercle de centre $O$ et de rayon $2$, et leurs arguments diffèrent de $\dfrac{\pi}{2}$ : elles forment un \textbf{carré} de centre $O$ inscrit dans ce cercle.
\begin{popfigure}
\begin{tikzpicture}[scale=0.7, >=Stealth]
\draw[popline, ->] (-3,0) -- (3,0) node[right] {$\mathbb{R}$};
\draw[popline, ->] (0,-3) -- (0,3) node[above] {$\mathbb{I}$};
\draw[popblueL, thick] (0,0) circle (2);
\foreach \k/\angle in {0/15, 1/105, 2/195, 3/285} {
  \draw[poporange, thick, ->] (0,0) -- (\angle:2);
  \fill[poporange] (\angle:2) circle (2pt) node[shift={(\angle:0.4)}] {$z_{\k}$};
}
\draw[dashed, popmuted] (0,0) -- (15:2) -- (105:2) -- (195:2) -- (285:2) -- cycle;
\end{tikzpicture}
\legende{Racines quatrièmes de $16e^{i\pi/3}$ : sommets d'un carré}
\end{popfigure}
\item $|-8\sqrt{3} + 8i| = \sqrt{192 + 64} = \sqrt{256} = 16$. On factorise : $-8\sqrt{3} + 8i = 16\left(-\dfrac{\sqrt{3}}{2} + i\dfrac{1}{2}\right)$. Or $\cos\dfrac{5\pi}{6} = -\dfrac{\sqrt{3}}{2}$ et $\sin\dfrac{5\pi}{6} = \dfrac{1}{2}$. Donc $-8\sqrt{3} + 8i = 16e^{i5\pi/6}$.
\item L'équation $z^4 = 16e^{i5\pi/6}$ a pour solutions :
\[ z_k = 2\,e^{i\left(\frac{5\pi}{24} + \frac{k\pi}{2}\right)}, \quad k \in \{0,1,2,3\}, \]
c'est-à-dire $z_0 = 2e^{i5\pi/24}$, $z_1 = 2e^{i17\pi/24}$, $z_2 = 2e^{i29\pi/24}$, $z_3 = 2e^{i41\pi/24}$.
\end{enumerate}
\end{corrige}

\begin{corrige}[Exercice 9 — Étude complète d'une configuration du plan complexe]
\begin{enumerate}
\item $z_A = 2\left(\cos\dfrac{\pi}{6} + i\sin\dfrac{\pi}{6}\right) = 2\left(\dfrac{\sqrt{3}}{2} + i\dfrac{1}{2}\right) = \sqrt{3} + i$.

$z_B = 2\left(\cos\dfrac{5\pi}{6} + i\sin\dfrac{5\pi}{6}\right) = 2\left(-\dfrac{\sqrt{3}}{2} + i\dfrac{1}{2}\right) = -\sqrt{3} + i$.
\item $z_C = z_A \times z_B = 2e^{i\pi/6} \times 2e^{i5\pi/6} = 4e^{i\pi}$. Donc $|z_C| = 4$ et $\arg(z_C) = \pi$ ; forme algébrique : $z_C = -4$.

Vérification par le calcul direct : $(\sqrt{3}+i)(-\sqrt{3}+i) = -3 + i\sqrt{3} - i\sqrt{3} + i^2 = -4$. Cohérent.
\item $A(\sqrt{3}\,;1) \approx (\num{1,73}\,;1)$, $B(-\sqrt{3}\,;1)$, $C(-4\,;0)$. Avec l'unité de \SI{2}{cm}, on multiplie chaque coordonnée par $2$ pour le report graphique.
\begin{popfigure}
\begin{tikzpicture}[scale=0.6, >=Stealth]
\draw[popline, ->] (-9,0) -- (5,0) node[right] {$\mathbb{R}$};
\draw[popline, ->] (0,-3) -- (0,5) node[above] {$\mathbb{I}$};
\foreach \x in {-8,-6,-4,-2,2,4} \draw (\x,0.1) -- (\x,-0.1) node[below] {\small $\x$};
\foreach \y in {2,4} \draw (0.1,\y) -- (-0.1,\y) node[left] {\small $\y$};
% Points
\coordinate (O) at (0,0);
\coordinate (A) at (3.46,2); % sqrt3*2, 1*2
\coordinate (B) at (-3.46,2);
\coordinate (C) at (-8,0);
\coordinate (D) at (-1,1.73); % -0.5*2, sqrt3/2*2
\fill[popblue] (A) circle (3pt) node[above right] {$A$};
\fill[popblue] (B) circle (3pt) node[above left] {$B$};
\fill[popblue] (C) circle (3pt) node[below left] {$C$};
\fill[poporange] (D) circle (3pt) node[above right] {$D$};
\fill[popdark] (O) circle (3pt) node[below left] {$O$};
% Segments
\draw[popblue, thick] (O) -- (A) -- (B) -- cycle;
\draw[poporange, thick, dashed] (O) -- (D) -- (A);
\draw[popmuted, thin] (O) -- (C);
% Right angle mark
\draw[poporange] (0.2,0) -- (0.2,0.2) -- (0,0.2);
\end{tikzpicture}
\legende{Configuration des points $O, A, B, C, D$ (unité \SI{2}{cm})}
\end{popfigure}
\item $OA = |z_A| = 2$ et $OB = |z_B| = 2$, donc $OA = OB$ : le triangle $OAB$ est \textbf{isocèle en $O$}.

Une mesure de $\widehat{AOB}$ est :
\[ \arg\left(\frac{z_B}{z_A}\right) = \arg(z_B) - \arg(z_A) = \frac{5\pi}{6} - \frac{\pi}{6} = \frac{4\pi}{6} = \frac{2\pi}{3}. \]
Donc $\widehat{AOB} = \dfrac{2\pi}{3}$ rad, soit $120\degree$.
\item \begin{enumerate}
\item $z_D = \dfrac{z_B}{z_A} = \dfrac{2e^{i5\pi/6}}{2e^{i\pi/6}} = e^{i2\pi/3}$.
\item $|z_D| = 1$, donc $OD = 1$. D'autre part $OA = 2$ et $AD = |z_D - z_A|$. Calculons : $z_D = -\dfrac{1}{2} + i\dfrac{\sqrt{3}}{2}$, donc $z_D - z_A = -\dfrac{1}{2} - \sqrt{3} + i\left(\dfrac{\sqrt{3}}{2} - 1\right)$, et
\[ AD^2 = \left(\frac{1}{2} + \sqrt{3}\right)^2 + \left(\frac{\sqrt{3}}{2} - 1\right)^2 = \frac{1}{4} + \sqrt{3} + 3 + \frac{3}{4} - \sqrt{3} + 1 = 5. \]
On a $OA^2 + OD^2 = 4 + 1 = 5 = AD^2$ : par la réciproque du théorème de Pythagore, le triangle $OAD$ est \textbf{rectangle en $O$}.
\end{enumerate}
\item $(z_A)^{12} = 2^{12}e^{i12\pi/6} = 4096\,e^{i2\pi} = 4096$. Le point image de $(z_A)^{12}$ est le point de l'axe des réels d'abscisse $4096$ : c'est un réel positif, situé sur la demi-droite $[O\,;\vec{u})$.
\end{enumerate}

\textbf{Barème indicatif (sur 10)} : formes algébriques de $z_A, z_B$ : 1,5 pt ; module et argument de $z_C$ : 1,5 pt ; placement : 1 pt ; triangle isocèle et angle : 2 pts ; forme exponentielle de $z_D$ : 1 pt ; nature du triangle $OAD$ : 2 pts ; $(z_A)^{12}$ : 1 pt.

\textbf{Points de vigilance} : citer les propriétés utilisées (« le module d'un produit est le produit des modules », « un argument d'un quotient est la différence des arguments ») ; pour la nature du triangle, la justification par la réciproque de Pythagore doit être rédigée complètement ; ne pas oublier que $\arg(z_B/z_A)$ donne une mesure de l'angle orienté $(\overrightarrow{OA}\,;\overrightarrow{OB})$.
\end{corrige}

\begin{corrige}[Exercice 11 — Racines de l'unité et polygone régulier]
\begin{enumerate}
\item Les solutions de $z^6 = 1 = e^{i\,0}$ sont les racines $6$-èmes de l'unité :
\[ z_k = e^{i\frac{2k\pi}{6}} = e^{i\frac{k\pi}{3}}, \quad k \in \{0,1,2,3,4,5\}. \]
Soit : $z_0 = 1$, $z_1 = e^{i\pi/3}$, $z_2 = e^{i2\pi/3}$, $z_3 = e^{i\pi} = -1$, $z_4 = e^{i4\pi/3}$, $z_5 = e^{i5\pi/3}$.
\item Les points $M_k$ sont sur le cercle de centre $O$ et de rayon $1$ (soit \SI{3}{cm} sur la figure), régulièrement espacés d'un angle de $\dfrac{\pi}{3}$. En les joignant, on obtient un \textbf{hexagone régulier} inscrit dans le cercle unité.
\begin{popfigure}
\begin{tikzpicture}[scale=1, >=Stealth]
\draw[popblueL, thick] (0,0) circle (3);
\draw[popline, ->] (-3.5,0) -- (3.5,0) node[right] {$\mathbb{R}$};
\draw[popline, ->] (0,-3.5) -- (0,3.5) node[above] {$\mathbb{I}$};
\foreach \k/\angle in {0/0, 1/60, 2/120, 3/180, 4/240, 5/300} {
  \draw[poporange, thick, ->] (0,0) -- (\angle:3);
  \fill[poporange] (\angle:3) circle (2.5pt) node[shift={(\angle:0.5)}] {$M_{\k}$};
}
\draw[poporange, thick] (0:3) -- (60:3) -- (120:3) -- (180:3) -- (240:3) -- (300:3) -- cycle;
\end{tikzpicture}
\legende{Racines $6$-èmes de l'unité : hexagone régulier (unité \SI{3}{cm})}
\end{popfigure}
\item $j = e^{i2\pi/3} = z_2$, et $j^6 = e^{i \times 6 \times 2\pi/3} = e^{i4\pi} = 1$, donc $j$ est bien solution de $(E)$.

Pour la somme : $1 + j + j^2$ est la somme des trois premiers termes d'une suite géométrique de raison $j \neq 1$ :
\[ 1 + j + j^2 = \frac{1 - j^3}{1 - j}. \]
Or $j^3 = e^{i2\pi} = 1$, donc $1 + j + j^2 = \dfrac{1-1}{1-j} = 0$.
\item $j = \cos\dfrac{2\pi}{3} + i\sin\dfrac{2\pi}{3} = -\dfrac{1}{2} + i\dfrac{\sqrt{3}}{2}$.

$j^2 = e^{i4\pi/3} = -\dfrac{1}{2} - i\dfrac{\sqrt{3}}{2}$ et $\bar{j} = -\dfrac{1}{2} - i\dfrac{\sqrt{3}}{2}$ : on a bien $j^2 = \bar{j}$.
\item Les six trous sont les sommets d'un hexagone régulier inscrit dans un cercle de rayon $R = \SI{5}{cm}$. La distance entre deux sommets consécutifs est le côté de l'hexagone :
\[ d = |z_{k+1} - z_k| \times R = 2R\sin\frac{\pi}{6} = 2 \times 5 \times \frac{1}{2} = \SI{5}{cm}. \]
Propriété remarquable : le côté d'un hexagone régulier est égal au rayon du cercle circonscrit. Le technicien doit donc espacer les centres des trous de $\SI{5}{cm}$ exactement.
\end{enumerate}

\textbf{Points de vigilance} : pour la somme géométrique, vérifier que la raison est différente de $1$ avant d'appliquer la formule ; pour la distance, justifier que l'écart angulaire entre deux racines consécutives est $\dfrac{2\pi}{6} = \dfrac{\pi}{3}$ et utiliser la corde $= 2R\sin\left(\dfrac{\text{angle}}{2}\right)$.
\end{corrige}

\courssub{Application technique : Électricité}

\begin{corrige}[Exercice 10 — Problème d'électricité : moteur asynchrone triphasé]
\begin{enumerate}
\item $|\underline{Z}| = \sqrt{8^2 + 6^2} = \sqrt{100} = \SI{10}{\Omega}$. Un argument $\theta$ vérifie $\cos\theta = \dfrac{8}{10} = 0{,}8$ et $\sin\theta = \dfrac{6}{10} = 0{,}6$, donc $\theta \approx \num{0,64}$ rad (soit environ $36{,}9\degree$).
\item $\underline{Z} = 10\,e^{i\,0{,}64}$ (en ohms).
\item $\underline{I} = \dfrac{\underline{U}}{\underline{Z}} = \dfrac{220\,e^{i\,0}}{10\,e^{i\,0{,}64}} = 22\,e^{-i\,0{,}64}$ (en ampères).

Forme algébrique : $\cos(0{,}64) \approx 0{,}80$ et $\sin(0{,}64) \approx 0{,}60$, donc $\underline{I} \approx 22(0{,}80 - 0{,}60i) \approx 17{,}6 - 13{,}2i$ (en A).

Vérification directe : $\underline{I} = \dfrac{220}{8+6i} = \dfrac{220(8-6i)}{100} = 2{,}2(8-6i) = 17{,}6 - 13{,}2i$. Cohérent.
\item La valeur efficace du courant est le module du phaseur : $I = |\underline{I}| = \SI{22}{A}$.
\item $\varphi = \arg(\underline{U}) - \arg(\underline{I}) = 0 - (-0{,}64) = 0{,}64$ rad. Donc $\cos\varphi = \cos(0{,}64) \approx 0{,}8$.

Remarque : on retrouve directement $\cos\varphi = \dfrac{\mathrm{Re}(\underline{Z})}{|\underline{Z}|} = \dfrac{8}{10} = 0{,}8$.
\item $P = 3UI\cos\varphi = 3 \times 220 \times 22 \times 0{,}8 = \SI{11616}{W}$, soit environ $\SI{11,6}{kW}$.
\item À puissance utile (active) fixée, un facteur de puissance faible impose un courant efficace $I$ plus grand, puisque $P = 3UI\cos\varphi$. Or un courant plus élevé augmente les pertes par effet Joule dans les câbles et le fournisseur facture souvent des pénalités pour la puissance réactive consommée. Un $\cos\varphi$ proche de $1$ (obtenu par exemple avec des batteries de condensateurs) réduit donc le courant appelé, les pertes et la facture d'électricité.
\end{enumerate}

\textbf{Points de vigilance} : bien distinguer la valeur efficace (module du phaseur) de la forme algébrique ; le déphasage se lit sur les arguments, pas sur les parties réelles ; arrondir uniquement à la fin des calculs.
\end{corrige}

\newpage

\coursec{Fiche bilan}

\begin{popfigure}
\centering
\begin{center}\tcbox[colback=popdark,colframe=popdark,coltext=white,arc=9pt,boxsep=4pt,left=6pt,right=6pt]{\bfseries\small Nombres complexes Approche géométrique}\end{center}
\vspace{-2pt}
\begin{tcbraster}[raster columns=3,raster equal height=rows,raster column skip=6pt,raster row skip=6pt,enhanced,blankest]
\begin{tcolorbox}[enhanced,colback=popblue,colframe=popblue,coltext=white,boxrule=0.9pt,arc=3pt,left=4pt,right=4pt,top=3pt,bottom=3pt,boxsep=1pt,halign=left]\bfseries\scriptsize 1. Représentation géométrique Point, Vecteur, Module, Argument\end{tcolorbox}
\begin{tcolorbox}[enhanced,colback=poporange,colframe=poporange,coltext=white,boxrule=0.9pt,arc=3pt,left=4pt,right=4pt,top=3pt,bottom=3pt,boxsep=1pt,halign=left]\bfseries\scriptsize 2. Forme exponentielle $z = \rho e^{i\theta}$ Produit, Quotient, Puissance\end{tcolorbox}
\begin{tcolorbox}[enhanced,colback=popgreen,colframe=popgreen,coltext=white,boxrule=0.9pt,arc=3pt,left=4pt,right=4pt,top=3pt,bottom=3pt,boxsep=1pt,halign=left]\bfseries\scriptsize 3. Linéarisation Formules de Moivre $\cos^n\theta$, $\sin^n\theta$ $\cos(n\theta)$, $\sin(n\theta)$\end{tcolorbox}
\begin{tcolorbox}[enhanced,colback=poppurple,colframe=poppurple,coltext=white,boxrule=0.9pt,arc=3pt,left=4pt,right=4pt,top=3pt,bottom=3pt,boxsep=1pt,halign=left]\bfseries\scriptsize 4. Racines $n$-ièmes $z_k = \sqrt[n]{\rho}\, e^{i\left(\frac{\theta}{n}+\frac{2k\pi}{n}\right)}$\end{tcolorbox}
\begin{tcolorbox}[enhanced,colback=poppink,colframe=poppink,coltext=white,boxrule=0.9pt,arc=3pt,left=4pt,right=4pt,top=3pt,bottom=3pt,boxsep=1pt,halign=left]\bfseries\scriptsize 5. Applications techniques Électricité, Signaux, Mécanique\end{tcolorbox}
\begin{tcolorbox}[enhanced,colback=popgold,colframe=popgold,coltext=white,boxrule=0.9pt,arc=3pt,left=4pt,right=4pt,top=3pt,bottom=3pt,boxsep=1pt,halign=left]\bfseries\scriptsize 6. Rédaction rigoureuse Justifications, Vocabulaire précis\end{tcolorbox}
\end{tcbraster}

\legende{Carte mentale : Leçon 11 -- Nombres complexes (Approche géométrique)}
\end{popfigure}

\coursec{Bilan de la leçon}

\courssub{Définitions et notations clés}

\begin{itemize}
    \item \cle{Forme algébrique} : $z = a + ib$ avec $a = \Re(z)$ (partie réelle) et $b = \Im(z)$ (partie imaginaire).
    \item \cle{Affixe} : le point $M(a\,;\,b)$ (ou le vecteur $\vec{u}(a\,;\,b)$) est l'image de $z$ dans le plan muni d'un repère orthonormé direct $(O\,;\,\vec{\imath},\vec{\jmath})$.
    \item \cle{Module} : $|z| = \sqrt{a^2 + b^2} = \rho \ge 0$. C'est la distance $OM$ (ou la norme de $\vec{u}$).
    \item \cle{Argument} : pour $z \neq 0$, $\arg(z) = \theta$ (en radians, défini à $2k\pi$ près) est une mesure de l'angle $(\vec{\imath}\,;\,\overrightarrow{OM})$. L'argument principal est l'unique argument appartenant à $]-\pi\,;\,\pi]$.
    \item \cle{Forme trigonométrique} : $z = \rho(\cos\theta + i\sin\theta)$ avec $\rho > 0$.
    \item \cle{Forme exponentielle (Euler)} : $z = \rho e^{i\theta}$ avec $\rho > 0$, où $e^{i\theta} = \cos\theta + i\sin\theta$.
    \item \cle{Conjugué} : $\overline{z} = a - ib = \rho e^{-i\theta}$ ; on a $z\overline{z} = |z|^2 = \rho^2$.
    \item \cle{Racines $n$-ièmes de l'unité} : $U_n = \left\{ \varepsilon_k = e^{i\frac{2k\pi}{n}} \;\middle|\; k = 0,\dots,n-1 \right\}$. Pour $n \ge 2$ : $\displaystyle\sum_{k=0}^{n-1} \varepsilon_k = 0$.
\end{itemize}

\courssub{Théorèmes et formules à connaître par cœur}

\begin{center}
\renewcommand{\arraystretch}{1.6}
\begin{tabularx}{\linewidth}{|>{\bfseries}l|X|}\hline
\rowcolor{popblueL}
\textbf{Opération} & \textbf{Forme exponentielle ($z_1=\rho_1 e^{i\theta_1}$, $z_2=\rho_2 e^{i\theta_2}$, $\rho_1,\rho_2>0$)} \\ \hline
Produit & $z_1 z_2 = \rho_1\rho_2 \, e^{i(\theta_1+\theta_2)}$ \\ \hline
Quotient ($z_2\neq 0$) & $\dfrac{z_1}{z_2} = \dfrac{\rho_1}{\rho_2} \, e^{i(\theta_1-\theta_2)}$ \\ \hline
Puissance entière & $z_1^{\,n} = \rho_1^{\,n} e^{in\theta_1}$ \quad ($n \in \mathbb{Z}$) \\ \hline
Racines $n$-ièmes de $z=\rho e^{i\theta}$ & $z_k = \sqrt[n]{\rho} \, e^{i\left(\frac{\theta}{n}+\frac{2k\pi}{n}\right)}$,\quad $k=0,\dots,n-1$ \\ \hline
Conjugué & $\overline{z_1} = \rho_1 e^{-i\theta_1}$ \\ \hline
Inverse ($z_1\neq 0$) & $z_1^{-1} = \dfrac{1}{\rho_1} e^{-i\theta_1}$ \\ \hline
\end{tabularx}
\end{center}

\begin{attention}[Conditions d'application à ne pas oublier]
\begin{itemize}
    \item L'argument n'est défini que pour un complexe \textbf{non nul} : $z = 0$ n'a pas d'argument.
    \item Dans la forme trigonométrique ou exponentielle, le module doit être \textbf{strictement positif} : $-2\left(\cos\frac{\pi}{3}+i\sin\frac{\pi}{3}\right)$ n'est \emph{pas} une forme trigonométrique ; il faut écrire $2\left(\cos\frac{4\pi}{3}+i\sin\frac{4\pi}{3}\right) = 2e^{-i\frac{2\pi}{3}}$.
    \item Les arguments s'additionnent pour un produit, mais le résultat doit être ramené dans $]-\pi\,;\,\pi]$ si l'on demande l'argument principal.
\end{itemize}
\end{attention}

\vspace{0.3cm}
\begin{center}
\renewcommand{\arraystretch}{1.6}
\begin{tabularx}{\linewidth}{|>{\bfseries}l|X|}\hline
\rowcolor{popgreenL}
\textbf{Formules de Moivre / Euler} & \textbf{Résultat} \\ \hline
Formule d'Euler & $e^{i\theta} = \cos\theta + i\sin\theta$ \\ \hline
Formule de Moivre & $(\cos\theta + i\sin\theta)^n = \cos(n\theta) + i\sin(n\theta)$ \quad ($n\in\mathbb{Z}$) \\ \hline
Formules d'Euler (bis) & $\cos\theta = \dfrac{e^{i\theta}+e^{-i\theta}}{2}$ \quad ; \quad $\sin\theta = \dfrac{e^{i\theta}-e^{-i\theta}}{2i}$ \\ \hline
$\cos(n\theta)$ & $\Re\big((\cos\theta + i\sin\theta)^n\big)$ \\ \hline
$\sin(n\theta)$ & $\Im\big((\cos\theta + i\sin\theta)^n\big)$ \\ \hline
\end{tabularx}
\end{center}

\begin{exemplebox}[Linéarisations classiques à savoir refaire]
Avec $\cos\theta = \frac{e^{i\theta}+e^{-i\theta}}{2}$ et le binôme de Newton :
\begin{align*}
\cos^2\theta &= \frac{1+\cos(2\theta)}{2} &
\sin^2\theta &= \frac{1-\cos(2\theta)}{2} \\
\cos^3\theta &= \frac{\cos(3\theta)+3\cos\theta}{4} &
\sin^3\theta &= \frac{3\sin\theta-\sin(3\theta)}{4}
\end{align*}
Détaillons $\cos^3\theta$ : $\cos^3\theta = \frac{1}{8}\left(e^{i\theta}+e^{-i\theta}\right)^3 = \frac{1}{8}\left(e^{3i\theta}+3e^{i\theta}+3e^{-i\theta}+e^{-3i\theta}\right) = \frac{1}{4}\left(\frac{e^{3i\theta}+e^{-3i\theta}}{2}+3\,\frac{e^{i\theta}+e^{-i\theta}}{2}\right) = \frac{\cos(3\theta)+3\cos\theta}{4}$.
\end{exemplebox}

\courssub{Méthodes express}

\begin{methode}[Passer de la forme algébrique à la forme exponentielle]
\begin{enumerate}
    \item Calculer le module : $\rho = \sqrt{a^2+b^2}$.
    \item Déterminer $\theta$ tel que $\cos\theta = \dfrac{a}{\rho}$ et $\sin\theta = \dfrac{b}{\rho}$ (les \textbf{deux} conditions sont nécessaires).
    \item Choisir $\theta \in \left]-\pi\,;\,\pi\right]$ selon le quadrant du point $M(a\,;\,b)$.
    \item Conclure : $z = \rho e^{i\theta}$.
\end{enumerate}
\end{methode}

\begin{methode}[Calculer les racines $n$-ièmes de $z \neq 0$]
\begin{enumerate}
    \item Mettre $z$ sous forme exponentielle : $z = \rho e^{i\theta}$ ($\rho>0$).
    \item Calculer le module commun : $r = \sqrt[n]{\rho}$.
    \item Calculer le premier argument : $\alpha_0 = \dfrac{\theta}{n}$.
    \item Écrire les $n$ racines : $z_k = r \, e^{i\left(\alpha_0 + \frac{2k\pi}{n}\right)}$, pour $k=0,\dots,n-1$.
    \item \textbf{Vérifications utiles} : la somme des racines vaut $0$ (si $n\ge 2$) ; les images des $z_k$ sont les sommets d'un polygone régulier à $n$ côtés inscrit dans le cercle de centre $O$ et de rayon $r$.
\end{enumerate}
\end{methode}

\begin{methode}[Linéariser $\cos^p\theta\,\sin^q\theta$ ou exprimer $\cos(n\theta)$, $\sin(n\theta)$]
\begin{enumerate}
    \item Pour \textbf{linéariser} : remplacer $\cos\theta = \frac{e^{i\theta}+e^{-i\theta}}{2}$ et $\sin\theta = \frac{e^{i\theta}-e^{-i\theta}}{2i}$.
    \item Développer avec le binôme de Newton.
    \item Regrouper les termes en $e^{ik\theta}$ et $e^{-ik\theta}$.
    \item Reconvertir : $\frac{e^{ik\theta}+e^{-ik\theta}}{2} = \cos(k\theta)$ et $\frac{e^{ik\theta}-e^{-ik\theta}}{2i} = \sin(k\theta)$.
    \item Pour \textbf{exprimer} $\cos(n\theta)$ ou $\sin(n\theta)$ : développer $(\cos\theta+i\sin\theta)^n$ avec le binôme, puis identifier parties réelle et imaginaire grâce à la formule de Moivre.
\end{enumerate}
\end{methode}

\begin{exoresolu}[Racines cubiques d'un complexe]
Résoudre dans $\mathbb{C}$ l'équation $z^3 = 4\sqrt{2}\,(-1+i)$ et représenter les solutions.
\tcblower
\textbf{Solution détaillée.}
\begin{enumerate}
    \item \textbf{Forme exponentielle du second membre.} On a $|-1+i| = \sqrt{2}$, donc $|4\sqrt{2}(-1+i)| = 4\sqrt{2}\times\sqrt{2} = 8$. De plus $-1+i = \sqrt{2}\left(\cos\frac{3\pi}{4}+i\sin\frac{3\pi}{4}\right)$, d'où
    \[ 4\sqrt{2}\,(-1+i) = 8\,e^{i\frac{3\pi}{4}}. \]
    \item \textbf{Module des racines} : $r = \sqrt[3]{8} = 2$.
    \item \textbf{Arguments} : $\alpha_k = \dfrac{1}{3}\left(\dfrac{3\pi}{4}+2k\pi\right) = \dfrac{\pi}{4}+\dfrac{2k\pi}{3}$, pour $k=0,1,2$.
    \item \textbf{Les trois solutions} :
    \[ z_0 = 2e^{i\frac{\pi}{4}}, \qquad z_1 = 2e^{i\frac{11\pi}{12}}, \qquad z_2 = 2e^{i\frac{19\pi}{12}} = 2e^{-i\frac{5\pi}{12}}. \]
    \item \textbf{Interprétation géométrique} : les images de $z_0, z_1, z_2$ sont les sommets d'un triangle équilatéral inscrit dans le cercle de centre $O$ et de rayon $2$.
\end{enumerate}
\textbf{Vérification} : $z_0^3 = 8e^{i\frac{3\pi}{4}} = 4\sqrt{2}\,(-1+i)$. \checkmark
\end{exoresolu}

\begin{aretenir}[Checklist avant le Probatoire]
\begin{itemize}
    \item $\square$ Savoir représenter un nombre complexe (point $M$ ou vecteur $\vec{u}$) dans un repère orthonormé direct.
    \item $\square$ Passer sans erreur des formes algébrique, trigonométrique et exponentielle (module $\rho$, argument principal $\theta \in \left]-\pi\,;\,\pi\right]$).
    \item $\square$ Appliquer les règles du produit, quotient, puissance, inverse en forme exponentielle (modules multipliés ou divisés, arguments additionnés ou soustraits).
    \item $\square$ Énoncer et utiliser la formule d'Euler : $e^{i\theta} = \cos\theta + i\sin\theta$.
    \item $\square$ Utiliser la formule de Moivre : $(\cos\theta+i\sin\theta)^n = \cos(n\theta)+i\sin(n\theta)$.
    \item $\square$ Linéariser $\cos^p\theta$, $\sin^q\theta$, $\cos^p\theta\sin^q\theta$ via les formules d'Euler et le binôme de Newton.
    \item $\square$ Exprimer $\cos(n\theta)$ et $\sin(n\theta)$ en fonction de $\cos\theta$ et $\sin\theta$.
    \item $\square$ Déterminer les $n$ racines $n$-ièmes d'un complexe non nul : module $\sqrt[n]{\rho}$, arguments $\frac{\theta}{n}+\frac{2k\pi}{n}$, $k=0,\dots,n-1$.
    \item $\square$ Connaître les propriétés des racines $n$-ièmes de l'unité : somme nulle pour $n\ge 2$, sommets d'un polygone régulier.
    \item $\square$ Résoudre une équation du type $z^n = a$ ($a \in \mathbb{C}^*$) et représenter les solutions géométriquement.
    \item $\square$ Rédiger une démonstration complète : justifier le choix de l'argument principal, préciser $k \in \{0,\dots,n-1\}$, conclure clairement.
\end{itemize}
\end{aretenir}

\findecours

\end{document}
